%% GENERATED by generate_papers.py -- edit the bank_*.py files, not this file.
\documentclass[11pt,a4paper]{article}
\usepackage[a4paper, top=20mm, bottom=20mm, left=18mm, right=18mm, headheight=14pt]{geometry}
\usepackage[T1]{fontenc}
\usepackage{lmodern}
\usepackage{amsmath, amssymb, mathtools}
\usepackage{braket}
\usepackage{array, tabularx, booktabs, multicol}
\usepackage{enumitem}
\usepackage[dvipsnames]{xcolor}
\usepackage[most]{tcolorbox}
\usepackage{needspace}
\usepackage{fancyhdr, lastpage}
\usepackage{pifont}
\usepackage{tikz}
\usepackage[hidelinks]{hyperref}
\newcolumntype{L}{>{\raggedright\arraybackslash}X}
\definecolor{brand}{HTML}{1F4E79}
\definecolor{okgreen}{HTML}{1B7F3B}
\setlength{\parindent}{0pt}
\setlength{\parskip}{3pt}
\newcommand{\optlabel}[1]{\textbf{(#1)}}
\newcommand{\correctlabel}[1]{\textcolor{okgreen}{\textbf{(#1)}\,\ding{51}}}
\newenvironment{qitem}[1]{\begin{list}{}{\setlength{\leftmargin}{3.3em}\setlength{\labelwidth}{3.1em}%
  \setlength{\labelsep}{0.2em}\setlength{\topsep}{3pt}\setlength{\parsep}{3pt}}\item[\textbf{Q.#1}\hfill]}{\end{list}}
\newcommand{\qsep}{\par\vspace{2pt}{\color{black!15}\hrule}\vspace{4pt}}
\newcommand{\qmeta}[2]{{\footnotesize\color{black!60}\textsc{#1} \;$\cdot$\; #2}\par}
\newtcolorbox{passagebox}[1]{enhanced, breakable, colback=brand!4, colframe=brand!70, fonttitle=\bfseries,
  title={#1}, left=6pt, right=6pt, top=4pt, bottom=4pt}
\newtcolorbox{sectionbanner}{enhanced, colback=brand, colframe=brand, coltext=white, fontupper=\large\bfseries,
  left=8pt, right=8pt, top=6pt, bottom=6pt, sharp corners}
\newtcolorbox{infobox}[1]{enhanced, breakable, colback=black!3, colframe=black!45, fonttitle=\bfseries,
  title={#1}, left=6pt, right=6pt, top=4pt, bottom=4pt}
\newtcolorbox{solbox}[1]{enhanced, breakable, colback=okgreen!3, colframe=okgreen!60!black, fonttitle=\bfseries,
  title={Answer: (#1) \quad Detailed solution}, left=6pt, right=6pt, top=3pt, bottom=4pt, before skip=4pt, after skip=6pt}
\newcommand{\QAOAop}{\mathrm{QAOA}}
\newcommand{\bub}[1]{\tikz[baseline=-0.65ex]\node[draw, circle, inner sep=0pt, minimum size=4.6mm, font=\scriptsize]{#1};}
\allowdisplaybreaks

\begin{document}
\pagestyle{fancy}\fancyhf{}
\fancyhead[L]{\small\color{black!70}Quantum Approximate Optimization Algorithm --- Assessment}
\fancyhead[R]{\small\color{black!70}\textbf{Answer Book}}
\fancyfoot[C]{\small Page \thepage\ of \pageref*{LastPage}}
\renewcommand{\headrulewidth}{0.4pt}

\thispagestyle{empty}
\begin{center}
{\color{brand}\rule{\linewidth}{1.2pt}}\\[6pt]
{\Huge\bfseries\color{brand} Quantum Approximate Optimization Algorithm}\\[4pt]
{\Large Assessment of the QAOA algorithm (Max-Cut)}\\[4pt]
{\large\bfseries ANSWER BOOK}\\[2pt]
{\color{brand}\rule{\linewidth}{1.2pt}}
\end{center}
This book accompanies the separate \emph{Question Paper} (80 questions, 3 hours, 240 marks). It contains
\begin{enumerate}[leftmargin=1.4em, itemsep=0pt, topsep=2pt]
\item the \textbf{answer key}, section by section, with the topic tested by every question (page~\pageref{key}),
\item a \textbf{topic-wise map} grouping the questions by theme and a \textbf{quick-marking grid} (page~\pageref{topicmap}), and
\item the \textbf{Answer Bank}: every question restated with the correct option marked \textcolor{okgreen}{\ding{51}} and a
      complete solution with derivations (from page~\pageref{bank}).
\end{enumerate}
\begin{infobox}{Marking scheme (every question)}
Each question has \textbf{four} options (A), (B), (C) and (D). \textbf{ONLY ONE} of these four options is correct.\\[2pt]
\begin{tabular}{@{}ll@{}}
\emph{Full marks}: & $+3$ if ONLY the correct option is chosen;\\
\emph{Zero marks}: & $\phantom{+}0$ if none of the options is chosen (question unanswered);\\
\emph{Negative marks}: & $-1$ in all other cases (wrong option, or more than one option chosen).
\end{tabular}
\end{infobox}

\begin{infobox}{How to evaluate}
For each section count $N_{\mathrm c}$ = number of correct answers and $N_{\mathrm w}$ = number of questions answered wrongly (or with more than
one bubble darkened). Section marks $= 3N_{\mathrm c} - N_{\mathrm w}$; unanswered questions score $0$. The total is out of $240$.
Every numerical answer was re-computed independently (NumPy state vectors in the notebook's convention, brute force over all splits, and CUDA-Q for the QAOA circuits) by
\texttt{verify\_answers.py} in the source folder.
\end{infobox}
\begin{infobox}{Conventions (identical to the Question Paper)}

\begin{itemize}[leftmargin=1.2em, itemsep=1pt, topsep=2pt]
\item Graph $G = (V, E)$ with nodes $0, 1, \dots, n-1$ and edge weights $w_{ij}$ ($w_{ij} = 1$ unless stated); $W = \sum_{(i,j)\in E} w_{ij}$.
      A split is a bitstring $z = z_0z_1\dots z_{n-1}$ ($z_i \in \{0,1\}$ = group of node $i$, \textbf{node 0 written first});
      $\bar z$ is $z$ with every bit flipped. Spins: $s_i = (-1)^{z_i} = 1 - 2z_i$, the eigenvalue of $Z_i$ on $\ket z$.
\item Cut value $C(z) = \sum_{(i,j)\in E} w_{ij}\,[z_i \neq z_j] = \sum_{(i,j)} w_{ij}\frac{1 - s_is_j}{2}$; $C_{\max} = \max_z C(z)$.
\item \textbf{Notebook convention} (\texttt{QAOA.ipynb}): $H_C = \sum_{(i,j)\in E}\frac{w_{ij}}{2}Z_iZ_j$, so $C = \frac W2 - H_C$;
      $H_M = \sum_i X_i$; $\ket{\boldsymbol\gamma,\boldsymbol\beta} = e^{-i\beta_pH_M}e^{-i\gamma_pH_C}\cdots e^{-i\beta_1H_M}e^{-i\gamma_1H_C}\ket{+}^{\otimes n}$;
      $F_p(\boldsymbol\gamma,\boldsymbol\beta) = \braket{\boldsymbol\gamma,\boldsymbol\beta|C|\boldsymbol\gamma,\boldsymbol\beta}$; approximation ratio $r = F_p/C_{\max}$.
      The angle vector is \texttt{params} $= [\gamma_1, \dots, \gamma_p, \beta_1, \dots, \beta_p]$.
\item Gates: $H = \frac{1}{\sqrt2}\begin{psmallmatrix}1&1\\1&-1\end{psmallmatrix}$; $X, Y, Z$ are the Pauli matrices; $R_X(\theta) = e^{-i\theta X/2}$,
      $R_Z(\theta) = e^{-i\theta Z/2}$; CNOT$(c, t)$ flips qubit $t$ when qubit $c$ is 1. Two-qubit matrices use the basis order
      $\ket{00}, \ket{01}, \ket{10}, \ket{11}$ (lower-numbered qubit first); $\ket{\pm} = \frac{1}{\sqrt2}(\ket0 \pm \ket1)$.
\item \textbf{Notebook defaults:} random graph with 15 nodes and edge probability $0.3$ (\texttt{SEED = 42}: 30 edges), $p = 2$, COBYLA with
      80 evaluations, ramp initial angles, 4000 shots, exact objective via \texttt{cudaq.observe}. A \emph{previously recorded run} means the
      outputs stored in the notebook (CUDA-Q 0.16.0, CPU target \texttt{qpp-cpu}).
\item All measurements are in the computational basis; $\ln$ is the natural logarithm. Numerical answers are rounded.
\end{itemize}
\end{infobox}
\newpage
\section*{Answer key}\label{key}
{\small Answer distribution: (A) 20, (B) 20, (C) 20, (D) 20.}

\needspace{14\baselineskip}
\subsection*{Section A --- Mathematics and Graph Theory for QAOA}
{\small\renewcommand{\arraystretch}{1.15}
\begin{tabularx}{\linewidth}{@{}r c L @{\hspace{1em}} r c L@{}}
\toprule
\textbf{Q.} & \textbf{Ans.} & \textbf{Topic tested} & \textbf{Q.} & \textbf{Ans.} & \textbf{Topic tested}\\ \midrule
1 & \textbf{A} & Bitstrings vs.\ partitions & 11 & \textbf{C} & Commuting exponentials \\
2 & \textbf{A} & $C_{\max}$ of standard graphs & 12 & \textbf{A} & Shots for a target standard error \\
3 & \textbf{B} & Random-split baseline $W/2$ & 13 & \textbf{A} & Binomial statistics of $P_{\text{opt}}$ \\
4 & \textbf{B} & Cut as a QUBO polynomial & 14 & \textbf{D} & Complementary optimal strings \\
5 & \textbf{D} & Random graph $G(n, q)$ statistics & 15 & \textbf{B} & QUBO $\to$ Ising conversion \\
6 & \textbf{C} & Handshake lemma & 16 & \textbf{C} & Exact shift rule for derivatives \\
7 & \textbf{C} & Odd cycles vs.\ odd node count & 17 & \textbf{D} & House graph: $C_{\max}$ and count \\
8 & \textbf{B} & Complexity: max cut vs.\ min cut & 18 & \textbf{D} & House graph: random baseline \\
9 & \textbf{A} & Spectrum of $H_C$ on a triangle & 19 & \textbf{B} & Ground energy $= W/2 - C_{\max}$ \\
10 & \textbf{C} & Variance of a random cut & 20 & \textbf{D} & Local search: $|E|/2$, not optimal \\
\bottomrule
\end{tabularx}}

\needspace{14\baselineskip}
\subsection*{Section B --- Quantum Mechanics}
{\small\renewcommand{\arraystretch}{1.15}
\begin{tabularx}{\linewidth}{@{}r c L @{\hspace{1em}} r c L@{}}
\toprule
\textbf{Q.} & \textbf{Ans.} & \textbf{Topic tested} & \textbf{Q.} & \textbf{Ans.} & \textbf{Topic tested}\\ \midrule
21 & \textbf{D} & Uniform state; QAOA enhancement & 31 & \textbf{B} & Conjugation $e^{i\beta X}Ze^{-i\beta X}$ \\
22 & \textbf{B} & Pauli products & 32 & \textbf{C} & State-vector memory on one GPU \\
23 & \textbf{D} & Rotation-gate matrices & 33 & \textbf{C} & Bit-flip probability of $R_X$ \\
24 & \textbf{A} & Relative phase and interference & 34 & \textbf{B} & $\ket+^{\otimes n}$ as top eigenstate of $H_M$ \\
25 & \textbf{D} & Expected cut from amplitudes & 35 & \textbf{B} & Complex-conjugation symmetry \\
26 & \textbf{A} & Commutation of Pauli strings & 36 & \textbf{D} & Linearity of $\langle H_C\rangle$ \\
27 & \textbf{D} & Pauli propagation through CNOT & 37 & \textbf{A} & Mixer at $\beta = \pi/2$; period $\pi/2$ \\
28 & \textbf{A} & Entanglement by the cost gate & 38 & \textbf{A} & Measuring $Z_iZ_j$ from shots \\
29 & \textbf{C} & Global phase vs.\ sign of $\gamma$ & 39 & \textbf{C} & Trace: phase after the cost gate \\
30 & \textbf{B} & $\langle Z_0Z_1\rangle$ for a product state & 40 & \textbf{C} & Trace: cut probability after the mixer \\
\bottomrule
\end{tabularx}}

\needspace{14\baselineskip}
\subsection*{Section C --- QAOA Basics}
{\small\renewcommand{\arraystretch}{1.15}
\begin{tabularx}{\linewidth}{@{}r c L @{\hspace{1em}} r c L@{}}
\toprule
\textbf{Q.} & \textbf{Ans.} & \textbf{Topic tested} & \textbf{Q.} & \textbf{Ans.} & \textbf{Topic tested}\\ \midrule
41 & \textbf{A} & \texttt{cost\_hamiltonian()} output & 51 & \textbf{C} & Parameter-vector layout \\
42 & \textbf{C} & All $\gamma_l = 0$ & 52 & \textbf{D} & The ramp initialisation \\
43 & \textbf{D} & CNOT--$R_Z$--CNOT decomposition & 53 & \textbf{A} & Minimise $\langle H_C\rangle$ = maximise $\langle C\rangle$ \\
44 & \textbf{D} & Why \texttt{rx(2*beta)} & 54 & \textbf{A} & Sign of the default start point \\
45 & \textbf{B} & Gate and CNOT counts & 55 & \textbf{C} & Algorithm $\leftrightarrow$ code \\
46 & \textbf{C} & Default circuit statistics & 56 & \textbf{B} & Triangle solved at $p = 1$ \\
47 & \textbf{D} & $p = 1$ single edge: formula, optimum & 57 & \textbf{B} & $r$ vs.\ best sample \\
48 & \textbf{B} & $p = 1$ on the ring & 58 & \textbf{C} & Best evaluated angles; no global guarantee \\
49 & \textbf{A} & $0.6924$ bound; monotonicity in $p$ & 59 & \textbf{A} & Path graph: $\langle Z_0Z_1\rangle$ \\
50 & \textbf{D} & Limit $p \to \infty$ & 60 & \textbf{B} & Path graph: $\langle C\rangle$, $P_{\text{opt}}$ \\
\bottomrule
\end{tabularx}}

\needspace{14\baselineskip}
\subsection*{Section D --- Advanced QAOA and Implementation}
{\small\renewcommand{\arraystretch}{1.15}
\begin{tabularx}{\linewidth}{@{}r c L @{\hspace{1em}} r c L@{}}
\toprule
\textbf{Q.} & \textbf{Ans.} & \textbf{Topic tested} & \textbf{Q.} & \textbf{Ans.} & \textbf{Topic tested}\\ \midrule
61 & \textbf{A} & Light cone of an edge & 71 & \textbf{C} & Depolarising-noise model \\
62 & \textbf{B} & Ring at $p = 2$: $5/6$ & 72 & \textbf{C} & QAOA variants \\
63 & \textbf{D} & Parameter shift vs.\ finite differences & 73 & \textbf{A} & INTERP initialisation \\
64 & \textbf{D} & COBYLA \texttt{rhobeg} & 74 & \textbf{D} & Heuristic reference: upper bound on $r$ \\
65 & \textbf{C} & Shot noise of the objective & 75 & \textbf{C} & Recorded timing per gate \\
66 & \textbf{B} & Brute-force memory & 76 & \textbf{B} & Run-time scaling to 25 nodes \\
67 & \textbf{D} & Multi-GPU state vector & 77 & \textbf{A} & Quantum advantage: state of knowledge \\
68 & \textbf{D} & Backend / method $\to$ target & 78 & \textbf{B} & Custom weighted edges \\
69 & \textbf{C} & When MPS is efficient & 79 & \textbf{A} & Recorded run: partition probability \\
70 & \textbf{A} & Shot-based objective on hardware & 80 & \textbf{B} & Recorded run: improvement over random \\
\bottomrule
\end{tabularx}}

\needspace{18\baselineskip}
\subsection*{Topic-wise map of the questions}\label{topicmap}
{\small Every question appears under at least one theme; several appear under two.\par}
{\small\renewcommand{\arraystretch}{1.2}
\begin{tabularx}{\linewidth}{@{}>{\raggedright\arraybackslash}p{0.42\linewidth}L@{}}
\toprule
\textbf{Theme} & \textbf{Question numbers}\\ \midrule
Graphs, cuts and Max-Cut & 1, 2, 6, 7, 14, 17, 18, 19, 20, 78 \\
QUBO / Ising formulation and the cost Hamiltonian & 4, 9, 15, 19, 29, 41, 53 \\
Probability, statistics and sampling & 3, 5, 10, 12, 13, 18, 21, 65, 79 \\
Complexity and classical algorithms & 8, 20, 49, 74, 77 \\
Linear algebra: commutators, exponentials, Pauli algebra & 11, 22, 23, 26, 27, 31, 37 \\
Qubits, measurement, phases and entanglement & 24, 25, 28, 30, 33, 34, 35, 36, 38, 39 \\
Circuit construction and gate counts & 43, 44, 45, 46, 55, 63 \\
Initial state, mixer and parameters & 34, 37, 42, 51, 52, 54, 73 \\
Analysis at $p = 1$ and depth dependence & 40, 47, 48, 49, 50, 56, 59, 60, 61, 62 \\
Classical optimizers and gradients & 16, 53, 58, 63, 64, 65, 73 \\
Backends, memory and performance & 32, 66, 67, 68, 69, 75, 76 \\
Noise, hardware, variants and limitations & 70, 71, 72, 77 \\
Interpreting recorded results of \texttt{QAOA.ipynb} & 13, 14, 21, 54, 57, 75, 79, 80 \\
\bottomrule
\end{tabularx}}

\subsection*{Quick-marking grid}
\begin{center}\renewcommand{\arraystretch}{1.3}
\begin{tabular}{|c|c|c|c|c|c|c|c|c|c|}\hline
1\,\textbf{A} & 2\,\textbf{A} & 3\,\textbf{B} & 4\,\textbf{B} & 5\,\textbf{D} & 6\,\textbf{C} & 7\,\textbf{C} & 8\,\textbf{B} & 9\,\textbf{A} & 10\,\textbf{C} \\
\hline
11\,\textbf{C} & 12\,\textbf{A} & 13\,\textbf{A} & 14\,\textbf{D} & 15\,\textbf{B} & 16\,\textbf{C} & 17\,\textbf{D} & 18\,\textbf{D} & 19\,\textbf{B} & 20\,\textbf{D} \\
\hline
21\,\textbf{D} & 22\,\textbf{B} & 23\,\textbf{D} & 24\,\textbf{A} & 25\,\textbf{D} & 26\,\textbf{A} & 27\,\textbf{D} & 28\,\textbf{A} & 29\,\textbf{C} & 30\,\textbf{B} \\
\hline
31\,\textbf{B} & 32\,\textbf{C} & 33\,\textbf{C} & 34\,\textbf{B} & 35\,\textbf{B} & 36\,\textbf{D} & 37\,\textbf{A} & 38\,\textbf{A} & 39\,\textbf{C} & 40\,\textbf{C} \\
\hline
41\,\textbf{A} & 42\,\textbf{C} & 43\,\textbf{D} & 44\,\textbf{D} & 45\,\textbf{B} & 46\,\textbf{C} & 47\,\textbf{D} & 48\,\textbf{B} & 49\,\textbf{A} & 50\,\textbf{D} \\
\hline
51\,\textbf{C} & 52\,\textbf{D} & 53\,\textbf{A} & 54\,\textbf{A} & 55\,\textbf{C} & 56\,\textbf{B} & 57\,\textbf{B} & 58\,\textbf{C} & 59\,\textbf{A} & 60\,\textbf{B} \\
\hline
61\,\textbf{A} & 62\,\textbf{B} & 63\,\textbf{D} & 64\,\textbf{D} & 65\,\textbf{C} & 66\,\textbf{B} & 67\,\textbf{D} & 68\,\textbf{D} & 69\,\textbf{C} & 70\,\textbf{A} \\
\hline
71\,\textbf{C} & 72\,\textbf{C} & 73\,\textbf{A} & 74\,\textbf{D} & 75\,\textbf{C} & 76\,\textbf{B} & 77\,\textbf{A} & 78\,\textbf{B} & 79\,\textbf{A} & 80\,\textbf{B} \\
\hline\end{tabular}\end{center}
\newpage
\section*{Answer Bank --- questions with detailed solutions}\label{bank}

\begin{sectionbanner}SECTION A \enspace$\cdot$\enspace Mathematics and Graph Theory for QAOA \hfill Q.1--20\end{sectionbanner}

\needspace{13\baselineskip}
\begin{qitem}{1}
\qmeta{Single-correct MCQ}{Bitstrings vs.\ partitions}
A graph has $n = 15$ nodes. The number of bitstrings $z \in \{0,1\}^{15}$, and the number of \emph{distinct}
partitions $\{S, V\setminus S\}$ of the node set (unordered; $S = \varnothing$ allowed) that these bitstrings describe, are
\par\smallskip\noindent\begin{tabularx}{\linewidth}{@{}*{4}{L}@{}}
\correctlabel{A}\,$32\,768$ and $16\,384$ & \optlabel{B}\,$16\,384$ and $8\,192$ & \optlabel{C}\,$32\,768$ and $16\,383$ & \optlabel{D}\,$32\,768$ and $32\,768$
\end{tabularx}
\end{qitem}
\begin{solbox}{A}
\textbf{\color{brand}Correct answer.} \textcolor{okgreen}{\textbf{(A)}} $32\,768$ and $16\,384$

\textbf{\color{brand}Concepts required.} A split is a bitstring $z$ ($z_i$ = group of node $i$). Relabelling the two groups (complementing $z$) gives the same partition and the same cut.

\textbf{\color{brand}Relevant formulae.} Number of bitstrings $= 2^n$; number of unordered partitions $= 2^n/2 = 2^{n-1}$ (pairs $\{z, \bar z\}$, $\bar z$ = all bits flipped).

\textbf{\color{brand}Step-by-step explanation.} Each node independently chooses group 0 or 1, so there are $2^{15} = 32\,768$ bitstrings.
A bitstring $z$ and its complement $\bar z$ (all bits flipped) describe the same unordered partition $\{S, V\setminus S\}$:
the two groups only swap their labels. Complementation pairs the $2^{15}$ strings perfectly (no string equals its own
complement), so the number of distinct partitions is $2^{15}/2 = 2^{14} = 16\,384$. The partition $\{\varnothing, V\}$
(strings $00\dots0$ and $11\dots1$, cut value $0$) is one of them because $S = \varnothing$ is allowed.

\textbf{\color{brand}Why the other options are incorrect.}
\begin{itemize}[leftmargin=2.4em, itemsep=1pt, topsep=2pt]
\item[\textbf{(B)}] $16\,384$ bitstrings would be $2^{14}$; each of the 15 nodes has 2 choices, so there are $2^{15}$ bitstrings.
\item[\textbf{(C)}] $16\,383$ excludes the partition $\{\varnothing, V\}$, but the question allows $S = \varnothing$ (cut value 0).
\item[\textbf{(D)}] $32\,768$ partitions counts every bitstring separately, forgetting that $z$ and $\bar z$ describe the same partition.
\end{itemize}
\textbf{\color{brand}Key takeaway.} $2^n$ bitstrings, $2^{n-1}$ partitions: every cut appears twice in a QAOA histogram, as $z$ and $\bar z$.
\end{solbox}
\qsep

\needspace{13\baselineskip}
\begin{qitem}{2}
\qmeta{Matching lists}{$C_{\max}$ of standard graphs}
Match each graph (List-I) with its maximum cut $C_{\max}$ (List-II). All edges have weight 1.
\par\smallskip\noindent{\renewcommand{\arraystretch}{1.35}%
\begin{tabularx}{\linewidth}{@{}c>{\raggedright\arraybackslash}p{0.33\linewidth}@{\hspace{1.2em}}cL@{}}
\toprule
\multicolumn{2}{@{}l}{\textbf{List-I}} & \multicolumn{2}{l@{}}{\textbf{List-II}} \\ \midrule
\textbf{(P)} & cycle $C_9$ (9 nodes, 9 edges) & \textbf{(1)} & $12$ \\
\textbf{(Q)} & complete graph $K_5$ & \textbf{(2)} & $9$ \\
\textbf{(R)} & complete bipartite graph $K_{3,3}$ & \textbf{(3)} & $6$ \\
\textbf{(S)} & Petersen graph (10 nodes, 15 edges, 3-regular) & \textbf{(4)} & $8$ \\
\bottomrule
\end{tabularx}}
\par\smallskip\noindent\begin{tabularx}{\linewidth}{@{}*{2}{L}@{}}
\correctlabel{A}\,P$\to$4, Q$\to$3, R$\to$2, S$\to$1 & \optlabel{B}\,P$\to$2, Q$\to$3, R$\to$1, S$\to$4 \\[3pt]
\optlabel{C}\,P$\to$4, Q$\to$2, R$\to$3, S$\to$1 & \optlabel{D}\,P$\to$4, Q$\to$3, R$\to$1, S$\to$2
\end{tabularx}
\end{qitem}
\begin{solbox}{A}
\textbf{\color{brand}Correct answer.} \textcolor{okgreen}{\textbf{(A)}} P$\to$4, Q$\to$3, R$\to$2, S$\to$1

\textbf{\color{brand}Concepts required.} $C_{\max} = |E|$ iff the graph is bipartite (no odd cycle). Each odd cycle forces at least one uncut edge. For $K_n$ only the group sizes matter.

\textbf{\color{brand}Relevant formulae.} $C_{\max}(C_n) = n$ ($n$ even), $n - 1$ ($n$ odd); $\ C_{\max}(K_n) = \lfloor n/2\rfloor\lceil n/2\rceil$; $\ C_{\max}(K_{a,b}) = ab$.

\textbf{\color{brand}Step-by-step explanation.} \begin{itemize}\itemsep0pt
\item $C_9$: an odd cycle cannot be 2-coloured, so at least one edge stays uncut; alternating colours around the cycle
cuts all but one edge: $C_{\max} = 9 - 1 = 8$ (4).
\item $K_5$: with $k$ nodes on one side every one of the $k(5-k)$ cross pairs is an edge; the maximum over $k$ is
$2 \cdot 3 = 6$ (3). In general $C_{\max}(K_n) = \lfloor n/2\rfloor\lceil n/2\rceil$.
\item $K_{3,3}$ is bipartite: putting the two natural sides in different groups cuts \emph{all} $9$ edges (2).
\item Petersen graph: it contains 5-cycles (odd), so not all 15 edges can be cut. Every one of its 12 five-cycles needs
an uncut edge, and the optimum leaves exactly 3 edges uncut: $C_{\max} = 12$ (1) (checked by brute force over all
$2^{10}$ splits in \texttt{verify\_answers.py}).
\end{itemize}
Hence P$\to$4, Q$\to$3, R$\to$2, S$\to$1.

\textbf{\color{brand}Why the other options are incorrect.}
\begin{itemize}[leftmargin=2.4em, itemsep=1pt, topsep=2pt]
\item[\textbf{(B)}] P$\to$2 gives $C_9$ all 9 edges, as if it were bipartite --- but an odd cycle always leaves one edge uncut; R$\to$1 and S$\to$4 then fail as above.
\item[\textbf{(C)}] Q$\to$2 gives $K_5$ the value 9; the best split of $K_5$ is $2 + 3$ nodes, cutting $2\cdot3 = 6$ edges; and R$\to$3 undercounts the bipartite $K_{3,3}$.
\item[\textbf{(D)}] R$\to$1, S$\to$2 gives $K_{3,3}$ the value 12 (it has only 9 edges) and the Petersen graph 9; the bipartite $K_{3,3}$ has all 9 edges cut, Petersen loses 3 of 15.
\end{itemize}
\textbf{\color{brand}Key takeaway.} Look for odd cycles first: they decide whether every edge can be cut.
\end{solbox}
\qsep

\needspace{13\baselineskip}
\begin{qitem}{3}
\qmeta{Single-correct MCQ}{Random-split baseline $W/2$}
A weighted graph has the edges $(0,1)$ with weight 3, $(1,2)$ with weight 1, $(2,3)$ with weight 2, $(3,0)$ with
weight 4 and $(0,2)$ with weight 2. Every node is put into group 0 or group 1 independently with probability $\tfrac12$.
The expected cut value is
\par\smallskip\noindent\begin{tabularx}{\linewidth}{@{}*{4}{L}@{}}
\optlabel{A}\,$12$ & \correctlabel{B}\,$6$ & \optlabel{C}\,$3$ & \optlabel{D}\,$2.4$
\end{tabularx}
\end{qitem}
\begin{solbox}{B}
\textbf{\color{brand}Correct answer.} \textcolor{okgreen}{\textbf{(B)}} $6$

\textbf{\color{brand}Concepts required.} Linearity of expectation; for independent fair bits each edge is cut with probability $\frac12$.

\textbf{\color{brand}Relevant formulae.} $\mathbb E[C] = \sum_{(i,j)\in E} w_{ij}\,P(z_i\neq z_j) = \tfrac12\sum_{(i,j)} w_{ij} = \tfrac W2$, $\ W$ = total weight.

\textbf{\color{brand}Step-by-step explanation.} Write $C(z) = \sum_{(i,j)\in E} w_{ij}\,[z_i \neq z_j]$. By linearity of expectation
\[ \mathbb{E}[C] = \sum_{(i,j)} w_{ij}\,P(z_i \neq z_j). \]
For independent fair bits $P(z_i \neq z_j) = P(01) + P(10) = \tfrac14 + \tfrac14 = \tfrac12$ for \emph{every} edge, whatever the graph looks like.
So $\mathbb{E}[C] = \tfrac12 W$ with $W = 3 + 1 + 2 + 4 + 2 = 12$, i.e.\ $\mathbb{E}[C] = 6$. This is exactly the
``random-split baseline'' $W/2$ printed by the notebook.

\textbf{\color{brand}Why the other options are incorrect.}
\begin{itemize}[leftmargin=2.4em, itemsep=1pt, topsep=2pt]
\item[\textbf{(A)}] $12$ is $W$ itself: it would require every edge to be cut with certainty.
\item[\textbf{(C)}] $3 = W/4$ uses $P(z_i\neq z_j) = \frac14$, i.e.\ counts only the outcome $01$ and forgets $10$.
\item[\textbf{(D)}] $2.4$ is the average weight $W/|E| = 12/5$, not an expected cut.
\end{itemize}
\textbf{\color{brand}Key takeaway.} The random-split baseline is always $W/2$, independent of the graph's structure.
\end{solbox}
\qsep

\needspace{13\baselineskip}
\begin{qitem}{4}
\qmeta{Single-correct MCQ}{Cut as a QUBO polynomial}
For the path graph $0 - 1 - 2$ (two unit edges), the cut value written as a polynomial in the binary variables
$z_0, z_1, z_2 \in \{0, 1\}$ is
\par\smallskip\noindent\begin{tabularx}{\linewidth}{@{}*{2}{L}@{}}
\optlabel{A}\,$z_0 + 2z_1 + z_2 - z_0z_1 - z_1z_2$ & \correctlabel{B}\,$z_0 + 2z_1 + z_2 - 2z_0z_1 - 2z_1z_2$ \\[3pt]
\optlabel{C}\,$z_0 + z_1 + z_2 - 2z_0z_1 - 2z_1z_2$ & \optlabel{D}\,$2 - z_0 - 2z_1 - z_2 + 2z_0z_1 + 2z_1z_2$
\end{tabularx}
\end{qitem}
\begin{solbox}{B}
\textbf{\color{brand}Correct answer.} \textcolor{okgreen}{\textbf{(B)}} $z_0 + 2z_1 + z_2 - 2z_0z_1 - 2z_1z_2$

\textbf{\color{brand}Concepts required.} XOR of two bits as a polynomial; QUBO = quadratic polynomial in binary variables.

\textbf{\color{brand}Relevant formulae.} $[a\neq b] = a \oplus b = a + b - 2ab$ for $a, b\in\{0,1\}$; $\ C(z) = \sum_{(i,j)} w_{ij}(z_i + z_j - 2z_iz_j)$.

\textbf{\color{brand}Step-by-step explanation.} For bits $a, b$: $[a \neq b] = a \oplus b = a + b - 2ab$ (check: $00 \to 0$, $01 \to 1$, $10 \to 1$, $11 \to 1 + 1 - 2 = 0$).
Summing over the two edges:
\[ C = (z_0 + z_1 - 2z_0z_1) + (z_1 + z_2 - 2z_1z_2) = z_0 + 2z_1 + z_2 - 2z_0z_1 - 2z_1z_2 . \]
The middle node appears in both edges, hence the coefficient $2$ of $z_1$ (in general the linear coefficient of $z_i$ is
the weighted degree of node $i$). Check $z = 010$: $C = 2$ (both edges cut) and the polynomial gives $0 + 2 + 0 - 0 - 0 = 2$.
This is the QUBO form of Max-Cut.

\textbf{\color{brand}Why the other options are incorrect.}
\begin{itemize}[leftmargin=2.4em, itemsep=1pt, topsep=2pt]
\item[\textbf{(A)}] $-z_0z_1 - z_1z_2$ uses $a + b - ab$, which is OR, not XOR: for $z = 111$ it gives $2$ instead of $0$.
\item[\textbf{(C)}] $z_0 + z_1 + z_2 - \dots$ forgets that node 1 belongs to two edges; for $z = 010$ it gives $1$ instead of $2$.
\item[\textbf{(D)}] $2 - z_0 - 2z_1 - \dots$ is $|E| - C$, the number of \emph{uncut} edges.
\end{itemize}
\textbf{\color{brand}Key takeaway.} The linear coefficient of $z_i$ in the cut polynomial is the (weighted) degree of node $i$.
\end{solbox}
\qsep

\needspace{13\baselineskip}
\begin{qitem}{5}
\qmeta{Single-correct MCQ}{Random graph $G(n, q)$ statistics}
The notebook's default graph is a random graph on $n = 15$ nodes in which each of the possible edges is present
independently with probability $0.3$. The expected number of edges and its standard deviation are
\par\smallskip\noindent\begin{tabularx}{\linewidth}{@{}*{4}{L}@{}}
\optlabel{A}\,$45$ and $3.97$ & \optlabel{B}\,$31.5$ and $22.05$ & \optlabel{C}\,$63$ and $6.64$ & \correctlabel{D}\,$31.5$ and $4.70$
\end{tabularx}
\end{qitem}
\begin{solbox}{D}
\textbf{\color{brand}Correct answer.} \textcolor{okgreen}{\textbf{(D)}} $31.5$ and $4.70$

\textbf{\color{brand}Concepts required.} Erd\H{o}s--R\'enyi random graph $G(n, q)$: each of the $\binom n2$ pairs is an edge independently with probability $q$.

\textbf{\color{brand}Relevant formulae.} $|E| \sim \mathrm{Binomial}\big(\binom n2, q\big)$: $\ \mathbb E = \binom n2 q$, $\ \sigma = \sqrt{\binom n2 q(1-q)}$; $\ \binom{15}{2} = 105$.

\textbf{\color{brand}Step-by-step explanation.} There are $\binom{15}{2} = \frac{15\cdot 14}{2} = 105$ unordered node pairs, each an independent Bernoulli($0.3$) trial.
The edge count is Binomial($105, 0.3$):
\[ \mathbb{E}[|E|] = 105 \times 0.3 = 31.5, \qquad \mathrm{Var} = 105 \times 0.3 \times 0.7 = 22.05, \qquad \sigma = \sqrt{22.05} \approx 4.70. \]
The instance actually generated with \texttt{SEED = 42} has $30$ edges, well within one standard deviation (the notebook's
text ``about 31 edges'' refers to this expectation; its printed output says 30).

\textbf{\color{brand}Why the other options are incorrect.}
\begin{itemize}[leftmargin=2.4em, itemsep=1pt, topsep=2pt]
\item[\textbf{(A)}] $45$ and $3.97$ use $\binom{10}{2} = 45$ pairs, i.e.\ the wrong number of nodes.
\item[\textbf{(B)}] $22.05$ is the variance, not the standard deviation (take the square root: $4.70$).
\item[\textbf{(C)}] $63$ and $6.64$ count ordered pairs ($15\cdot14 = 210$); an undirected edge is an unordered pair.
\end{itemize}
\textbf{\color{brand}Key takeaway.} Expected edges $= \binom n2 q$; the actual seed-42 graph (30 edges) is a typical draw.
\end{solbox}
\qsep

\needspace{13\baselineskip}
\begin{qitem}{6}
\qmeta{Single-correct MCQ}{Handshake lemma}
The 15-node graph of the notebook has the degree sequence (nodes $0$ to $14$)
\[ (d_0, \dots, d_{14}) = (2, 3, 5, 4, 5, 4, 6, 1, 4, 3, 4, 2, 4, 6, 7). \]
The number of edges, and the number of nodes of odd degree, are
\par\smallskip\noindent\begin{tabularx}{\linewidth}{@{}*{4}{L}@{}}
\optlabel{A}\,$30$ and $5$ & \optlabel{B}\,$60$ and $6$ & \correctlabel{C}\,$30$ and $6$ & \optlabel{D}\,$30$ and $9$
\end{tabularx}
\end{qitem}
\begin{solbox}{C}
\textbf{\color{brand}Correct answer.} \textcolor{okgreen}{\textbf{(C)}} $30$ and $6$

\textbf{\color{brand}Concepts required.} Handshake lemma; parity of the number of odd-degree nodes.

\textbf{\color{brand}Relevant formulae.} $\sum_{i\in V} d_i = 2|E|$; hence the number of odd-degree nodes is even.

\textbf{\color{brand}Step-by-step explanation.} \textbf{Handshake lemma:} every edge contributes 1 to the degree of each of its two end nodes, so
$\sum_i d_i = 2|E|$. Here $\sum_i d_i = 2+3+5+4+5+4+6+1+4+3+4+2+4+6+7 = 60$, hence $|E| = 30$.
Odd degrees: $3, 5, 5, 1, 3, 7$ (nodes 1, 2, 4, 7, 9, 14) --- six nodes. The number of odd-degree nodes is always even,
because $\sum_i d_i$ is even. (The degrees were recomputed from the edge list generated with \texttt{SEED = 42}.)

\textbf{\color{brand}Why the other options are incorrect.}
\begin{itemize}[leftmargin=2.4em, itemsep=1pt, topsep=2pt]
\item[\textbf{(A)}] Five odd-degree nodes is impossible: their count must be even (the six are nodes 1, 2, 4, 7, 9, 14).
\item[\textbf{(B)}] $60$ is $\sum d_i$, which counts every edge twice.
\item[\textbf{(D)}] $9$ is the number of \emph{even}-degree nodes ($15 - 6$).
\end{itemize}
\textbf{\color{brand}Key takeaway.} Divide the degree sum by two; odd degrees always come in pairs.
\end{solbox}
\qsep

\needspace{13\baselineskip}
\begin{qitem}{7}
\qmeta{Assertion--Reason}{Odd cycles vs.\ odd node count}
\textbf{Assertion (A):} The cycle $C_7$ (7 nodes, 7 unit edges) has $C_{\max} = 6$.\\[1mm]
\textbf{Reason (R):} Every graph with an odd number of nodes has $C_{\max} < |E|$.
\begin{itemize}[leftmargin=2.6em, labelsep=0.4em, itemsep=2pt, topsep=3pt]
\item[\optlabel{A}] Both Assertion (A) and Reason (R) are true, and R is the correct explanation of A
\item[\optlabel{B}] Both A and R are true, but R is \emph{not} the correct explanation of A
\item[\correctlabel{C}] A is true, but R is false
\item[\optlabel{D}] A is false, but R is true
\end{itemize}
\end{qitem}
\begin{solbox}{C}
\textbf{\color{brand}Correct answer.} \textcolor{okgreen}{\textbf{(C)}} A is true, but R is false

\textbf{\color{brand}Concepts required.} $C_{\max} = |E|$ exactly for bipartite graphs; a graph is bipartite iff it has no odd cycle.

\textbf{\color{brand}Relevant formulae.} $C_{\max}(C_n) = n - 1$ for odd $n$; bipartite $\iff$ no odd cycle $\iff$ $C_{\max} = |E|$.

\textbf{\color{brand}Step-by-step explanation.} \textbf{A is true:} $C_{\max} = |E|$ would require every edge to join the two groups, i.e.\ a proper 2-colouring;
going round an odd cycle the colours must alternate an odd number of times and cannot close up. So at least one edge
is uncut, and the alternating colouring $0101010$ cuts exactly $6$ edges.
\textbf{R is false:} $C_{\max} = |E|$ holds exactly for \emph{bipartite} graphs, and bipartiteness has nothing to do with the
parity of the number of nodes. Counter-example: the path on 3 nodes, or the star with 7 nodes (one centre joined to six
leaves) --- both bipartite, both with $C_{\max} = |E|$. What matters is an \emph{odd cycle}, not an odd node count. Answer: (C).

\textbf{\color{brand}Why the other options are incorrect.}
\begin{itemize}[leftmargin=2.4em, itemsep=1pt, topsep=2pt]
\item[\textbf{(A)}] needs A true, R true \emph{and} R explaining A; R is false (an odd number of nodes does not prevent bipartiteness).
\item[\textbf{(B)}] needs A and R both true but R not explaining A; R is false.
\item[\textbf{(D)}] needs A false and R true; A is true ($C_7$ has $C_{\max} = 6$) and R is false.
\end{itemize}
\textbf{\color{brand}Key takeaway.} Odd \emph{cycles}, not odd node counts, block a perfect cut.
\end{solbox}
\qsep

\needspace{13\baselineskip}
\begin{qitem}{8}
\qmeta{Statement I / II}{Complexity: max cut vs.\ min cut}
\textbf{Statement I:} A polynomial-time algorithm that solves Max-Cut \emph{exactly} on every graph is known.\\[1mm]
\textbf{Statement II:} Finding a \emph{minimum} cut of a graph (both groups non-empty, non-negative weights) is NP-hard.
\begin{itemize}[leftmargin=2.6em, labelsep=0.4em, itemsep=2pt, topsep=3pt]
\item[\optlabel{A}] Both Statement I and Statement II are true
\item[\correctlabel{B}] Both Statement I and Statement II are false
\item[\optlabel{C}] Statement I is true, but Statement II is false
\item[\optlabel{D}] Statement I is false, but Statement II is true
\end{itemize}
\end{qitem}
\begin{solbox}{B}
\textbf{\color{brand}Correct answer.} \textcolor{okgreen}{\textbf{(B)}} Both Statement I and Statement II are false

\textbf{\color{brand}Concepts required.} Complexity: Max-Cut is NP-hard; minimum cut is polynomial.

\textbf{\color{brand}Relevant formulae.} Best polynomial-time approximation known: Goemans--Williamson, $\mathbb E[C] \ge 0.878\,C_{\max}$; global min cut: Stoer--Wagner $O(|V||E| + |V|^2\log|V|)$.

\textbf{\color{brand}Step-by-step explanation.} \textbf{I is false:} Max-Cut is NP-hard (one of Karp's 21 NP-complete problems, decision version). No
polynomial-time exact algorithm is known, and none exists unless P = NP. The best known polynomial-time guarantee is the
Goemans--Williamson SDP rounding with expected ratio $\approx 0.878$; exact solution in the notebook uses brute force over $2^n$ splits.
\textbf{II is false:} the global minimum cut is polynomial --- e.g.\ the Stoer--Wagner algorithm, or $n-1$ max-flow/min-cut
computations. Maximising and minimising a cut have completely different complexity. Answer: (B).

\textbf{\color{brand}Why the other options are incorrect.}
\begin{itemize}[leftmargin=2.4em, itemsep=1pt, topsep=2pt]
\item[\textbf{(A)}] needs both statements true; neither statement is true.
\item[\textbf{(C)}] needs I true and II false; I is false: no exact polynomial algorithm for Max-Cut is known.
\item[\textbf{(D)}] needs I false and II true; II is false: min cut is solvable in polynomial time.
\end{itemize}
\textbf{\color{brand}Key takeaway.} Maximising a cut is hard, minimising it is easy --- the direction matters.
\end{solbox}
\qsep

\needspace{13\baselineskip}
\begin{qitem}{9}
\qmeta{Single-correct MCQ}{Spectrum of $H_C$ on a triangle}
For the triangle $K_3$ (unit weights) consider $H_C = \tfrac12\big(Z_0Z_1 + Z_1Z_2 + Z_0Z_2\big)$. Its smallest eigenvalue
and the number of computational basis states with that eigenvalue are
\par\smallskip\noindent\begin{tabularx}{\linewidth}{@{}*{2}{L}@{}}
\correctlabel{A}\,$-\tfrac12$, six basis states & \optlabel{B}\,$-\tfrac12$, three basis states \\[3pt]
\optlabel{C}\,$+\tfrac32$, two basis states & \optlabel{D}\,$-\tfrac32$, two basis states
\end{tabularx}
\end{qitem}
\begin{solbox}{A}
\textbf{\color{brand}Correct answer.} \textcolor{okgreen}{\textbf{(A)}} $-\tfrac12$, six basis states

\textbf{\color{brand}Concepts required.} $H_C$ is diagonal; its value on $\ket z$ is computed from the spins $s_i = (-1)^{z_i}$; frustration.

\textbf{\color{brand}Relevant formulae.} $H_C(z) = \sum_{(i,j)}\tfrac{w_{ij}}{2}s_is_j$; on a triangle $s_0s_1\cdot s_1s_2\cdot s_0s_2 = +1$; $\ C(z) = \tfrac W2 - H_C(z)$.

\textbf{\color{brand}Step-by-step explanation.} $H_C$ is diagonal: on $\ket{z}$ it equals $\tfrac12(s_0s_1 + s_1s_2 + s_0s_2)$ with $s_i = (-1)^{z_i} = \pm1$.
\begin{itemize}\itemsep0pt
\item $z = 000, 111$: all spins equal, every product $+1$: eigenvalue $\tfrac32$ (2 states).
\item one node differs from the other two (6 states such as $001, 010, \dots$): products $-1, -1, +1$, eigenvalue $-\tfrac12$.
\end{itemize}
$-\tfrac32$ would need all three products to be $-1$, but $(s_0s_1)(s_1s_2)(s_0s_2) = s_0^2 s_1^2 s_2^2 = +1$: impossible
(the triangle is \emph{frustrated}). So $E_0 = -\tfrac12$ with 6 basis states, corresponding to $C = \tfrac32 - E_0 = 2 = C_{\max}$.

\textbf{\color{brand}Why the other options are incorrect.}
\begin{itemize}[leftmargin=2.4em, itemsep=1pt, topsep=2pt]
\item[\textbf{(B)}] Three counts partitions ($\{0\}|\{1,2\}$ etc.); each partition is two bitstrings, so six basis states.
\item[\textbf{(C)}] $+\frac32$ is the \emph{largest} eigenvalue (strings $000$, $111$, cut $0$).
\item[\textbf{(D)}] $-\frac32$ needs all three products $s_is_j = -1$, impossible because their product is $+1$ (frustration).
\end{itemize}
\textbf{\color{brand}Key takeaway.} Ground states of $H_C$ = optimal cuts; on a frustrated triangle the ground energy is $-\frac12$, six-fold.
\end{solbox}
\qsep

\needspace{13\baselineskip}
\begin{qitem}{10}
\qmeta{Single-correct MCQ}{Variance of a random cut}
For the notebook's 15-node graph ($30$ unit edges), a split $z$ is chosen uniformly at random from all $2^{15}$
bitstrings. The variance of the cut value $C(z)$ is
\par\smallskip\noindent\begin{tabularx}{\linewidth}{@{}*{4}{L}@{}}
\optlabel{A}\,$15$ & \optlabel{B}\,$3.75$ & \correctlabel{C}\,$7.5$ & \optlabel{D}\,$30$
\end{tabularx}
\end{qitem}
\begin{solbox}{C}
\textbf{\color{brand}Correct answer.} \textcolor{okgreen}{\textbf{(C)}} $7.5$

\textbf{\color{brand}Concepts required.} Variance of a sum; pairwise independence of the edge indicators for a uniformly random split.

\textbf{\color{brand}Relevant formulae.} $X_e = z_i\oplus z_j \sim$ Bernoulli$(\tfrac12)$, $\mathrm{Var}(X_e) = \tfrac14$, $\mathrm{Cov}(X_e, X_f) = 0$ $\Rightarrow$ $\mathrm{Var}(C) = \sum_e w_e^2/4$ ($= |E|/4$ unweighted).

\textbf{\color{brand}Step-by-step explanation.} Write $C = \sum_{e} X_e$ with $X_e = [z_i \neq z_j]$ for edge $e = (i,j)$. Each $X_e$ is Bernoulli($\tfrac12$):
$\mathrm{Var}(X_e) = \tfrac12\cdot\tfrac12 = \tfrac14$.
\textbf{Pairwise independence:} for two edges sharing a node, $e = (i,j)$, $f = (j,k)$: $X_e = z_i \oplus z_j$ and
$X_f = z_j \oplus z_k$ are independent fair bits (given any $z_j$, $z_i$ and $z_k$ are independent and fair), so
$\mathrm{Cov}(X_e, X_f) = 0$; edges with no common node are trivially independent. Hence
\[ \mathrm{Var}(C) = \sum_e \mathrm{Var}(X_e) = \frac{|E|}{4} = \frac{30}{4} = 7.5 . \]
(Exact enumeration of the $2^{15}$ cut values of the actual graph gives $7.5$.) The standard deviation is $\approx 2.74$, so
$C_{\max} = 25$ lies $(25 - 15)/2.74 \approx 3.7$ standard deviations above the random mean $15$.

\textbf{\color{brand}Why the other options are incorrect.}
\begin{itemize}[leftmargin=2.4em, itemsep=1pt, topsep=2pt]
\item[\textbf{(A)}] $15$ is the mean $W/2$, not the variance.
\item[\textbf{(B)}] $3.75$ halves the correct value, e.g.\ by assuming that edges sharing a node are anticorrelated --- they are uncorrelated.
\item[\textbf{(D)}] $30 = |E|$ uses variance 1 per edge (the variance of a $\pm1$ variable), but $X_e \in \{0,1\}$ has variance $\frac14$.
\end{itemize}
\textbf{\color{brand}Key takeaway.} For a random split the cut has mean $|E|/2$ and variance $|E|/4$; good cuts are many standard deviations above the mean.
\end{solbox}
\qsep

\needspace{13\baselineskip}
\begin{qitem}{11}
\qmeta{Statement I / II}{Commuting exponentials}
\textbf{Statement I:} Because all the operators $Z_iZ_j$ commute with one another,\\
\hspace*{1.5em}$e^{-i\gamma H_C} = \prod_{(i,j)\in E} e^{-i\gamma\frac{w_{ij}}{2} Z_iZ_j}$ holds exactly, with the factors in any order.\\[1mm]
\textbf{Statement II:} Because the operators $X_i$ also commute with one another,\\
\hspace*{1.5em}$e^{-i\beta H_M}\,e^{-i\gamma H_C} = e^{-i(\beta H_M + \gamma H_C)}$ holds for all $\beta, \gamma$.
\begin{itemize}[leftmargin=2.6em, labelsep=0.4em, itemsep=2pt, topsep=3pt]
\item[\optlabel{A}] Both Statement I and Statement II are true
\item[\optlabel{B}] Both Statement I and Statement II are false
\item[\correctlabel{C}] Statement I is true, but Statement II is false
\item[\optlabel{D}] Statement I is false, but Statement II is true
\end{itemize}
\end{qitem}
\begin{solbox}{C}
\textbf{\color{brand}Correct answer.} \textcolor{okgreen}{\textbf{(C)}} Statement I is true, but Statement II is false

\textbf{\color{brand}Concepts required.} $e^{A+B} = e^Ae^B$ when $[A,B] = 0$; Baker--Campbell--Hausdorff otherwise.

\textbf{\color{brand}Relevant formulae.} $[Z_iZ_j, Z_kZ_l] = 0$; $\{X_i, Z_iZ_j\} = 0$ (anticommute); $e^{A}e^{B} = e^{A + B + \frac12[A,B] + \dots}$.

\textbf{\color{brand}Step-by-step explanation.} For matrices, $e^{A+B} = e^Ae^B$ holds whenever $[A, B] = 0$, but not in general.
\textbf{I is true:} all $Z_iZ_j$ are diagonal, hence commute, so the exponential of their sum factorises exactly into one
two-qubit gate per edge, in any order. That is why the notebook can loop over the edges in list order.
\textbf{II is false:} commuting \emph{within} $H_M$ only gives $e^{-i\beta H_M} = \prod_i e^{-i\beta X_i}$. The relevant
commutator is $[H_M, H_C] \neq 0$ (e.g.\ $X_i$ anticommutes with $Z_iZ_j$), so $e^{-i\beta H_M}e^{-i\gamma H_C} \neq
e^{-i(\beta H_M + \gamma H_C)}$ in general (Baker--Campbell--Hausdorff). The non-commutativity is essential: if they commuted,
QAOA could not move probability between bitstrings in a useful way. Answer: (C).

\textbf{\color{brand}Why the other options are incorrect.}
\begin{itemize}[leftmargin=2.4em, itemsep=1pt, topsep=2pt]
\item[\textbf{(A)}] needs both statements true; II is false.
\item[\textbf{(B)}] needs both statements false; I is true.
\item[\textbf{(D)}] needs I false and II true; I is true and II is false.
\end{itemize}
\textbf{\color{brand}Key takeaway.} Within a layer everything commutes; between cost and mixer layers nothing does --- that is where QAOA's power lies.
\end{solbox}
\qsep

\needspace{13\baselineskip}
\begin{qitem}{12}
\qmeta{Single-correct MCQ}{Shots for a target standard error}
The expectation value $\langle Z_iZ_j\rangle$ is estimated as the average of the $\pm1$ values $(-1)^{z_i \oplus z_j}$ over
$S$ independent shots. The smallest $S$ that guarantees a standard error of at most $0.01$ \emph{for every state} is
\par\smallskip\noindent\begin{tabularx}{\linewidth}{@{}*{4}{L}@{}}
\correctlabel{A}\,$10\,000$ & \optlabel{B}\,$2\,500$ & \optlabel{C}\,$100$ & \optlabel{D}\,$40\,000$
\end{tabularx}
\end{qitem}
\begin{solbox}{A}
\textbf{\color{brand}Correct answer.} \textcolor{okgreen}{\textbf{(A)}} $10\,000$

\textbf{\color{brand}Concepts required.} Standard error of a sample mean; variance of a $\pm1$ random variable.

\textbf{\color{brand}Relevant formulae.} $Y = \pm1$: $\mathrm{Var}(Y) = 1 - \mu^2 \le 1$; $\ \mathrm{SE} = \sqrt{\mathrm{Var}(Y)/S} \le 1/\sqrt S$.

\textbf{\color{brand}Step-by-step explanation.} Each shot gives a random variable $Y = \pm1$ with mean $\mu = \langle Z_iZ_j\rangle$ and variance
$\mathrm{Var}(Y) = \mathbb{E}[Y^2] - \mu^2 = 1 - \mu^2 \le 1$ (worst case $\mu = 0$). The average of $S$ independent shots has
standard error $\sqrt{(1-\mu^2)/S} \le 1/\sqrt S$. Requiring $1/\sqrt S \le 0.01$ gives $S \ge 10^4$.
(For the $0/1$ indicator $[z_i \neq z_j]$ the variance is at most $\tfrac14$, giving $2\,500$ --- but that estimates
$P(\text{cut})$, whose error is half the error of $\langle Z_iZ_j\rangle = 1 - 2P(\text{cut})$.)

\textbf{\color{brand}Why the other options are incorrect.}
\begin{itemize}[leftmargin=2.4em, itemsep=1pt, topsep=2pt]
\item[\textbf{(B)}] $2\,500$ uses the variance bound $\frac14$ of a $0/1$ variable; $\langle Z_iZ_j\rangle$ is an average of $\pm1$ values, whose variance can be 1.
\item[\textbf{(C)}] $100$ makes $1/S$ (not $1/\sqrt S$) equal to $0.01$.
\item[\textbf{(D)}] $40\,000$ uses $(2/\varepsilon)^2$, confusing the width of the range $[-1, 1]$ with the standard deviation.
\end{itemize}
\textbf{\color{brand}Key takeaway.} Shot noise falls only like $1/\sqrt S$: one more decimal digit costs $100\times$ more shots.
\end{solbox}
\qsep

\needspace{13\baselineskip}
\begin{qitem}{13}
\qmeta{Matching lists}{Binomial statistics of $P_{\text{opt}}$}
In the previously recorded notebook run, an optimal split is sampled with probability $P_{\text{opt}} = 0.06$ per
shot. Treat the shots as independent and match each quantity (List-I) with its value (List-II).
\par\smallskip\noindent{\renewcommand{\arraystretch}{1.35}%
\begin{tabularx}{\linewidth}{@{}c>{\raggedright\arraybackslash}p{0.33\linewidth}@{\hspace{1.2em}}cL@{}}
\toprule
\multicolumn{2}{@{}l}{\textbf{List-I}} & \multicolumn{2}{l@{}}{\textbf{List-II}} \\ \midrule
\textbf{(P)} & standard deviation of the number of optimal shots among $4000$ & \textbf{(1)} & $75$ \\
\textbf{(Q)} & probability that $10$ shots contain \emph{no} optimal split & \textbf{(2)} & $15.0$ \\
\textbf{(R)} & probability that $10$ shots contain \emph{at least one} optimal split & \textbf{(3)} & $0.539$ \\
\textbf{(S)} & smallest number of shots giving $\ge 99\,\%$ probability of at least one optimal split & \textbf{(4)} & $0.461$ \\
\bottomrule
\end{tabularx}}
\par\smallskip\noindent\begin{tabularx}{\linewidth}{@{}*{2}{L}@{}}
\correctlabel{A}\,P$\to$2, Q$\to$3, R$\to$4, S$\to$1 & \optlabel{B}\,P$\to$2, Q$\to$4, R$\to$3, S$\to$1 \\[3pt]
\optlabel{C}\,P$\to$1, Q$\to$4, R$\to$3, S$\to$2 & \optlabel{D}\,P$\to$1, Q$\to$3, R$\to$4, S$\to$2
\end{tabularx}
\end{qitem}
\begin{solbox}{A}
\textbf{\color{brand}Correct answer.} \textcolor{okgreen}{\textbf{(A)}} P$\to$2, Q$\to$3, R$\to$4, S$\to$1

\textbf{\color{brand}Concepts required.} Binomial distribution of the number of successes in independent shots; complement rule.

\textbf{\color{brand}Relevant formulae.} $\sigma = \sqrt{SP(1-P)}$; $\ P(\text{none in } k) = (1-P)^k$; $\ S_{99\%} = \lceil \ln 0.01/\ln(1-P)\rceil$.

\textbf{\color{brand}Step-by-step explanation.} \begin{itemize}\itemsep0pt
\item P: the count is Binomial($4000, 0.06$), $\sigma = \sqrt{4000\cdot0.06\cdot0.94} = \sqrt{225.6} \approx 15.0$ (2).
(Its mean is $240$, exactly the number of optimal shots recorded.)
\item Q: $P(\text{none in } 10) = 0.94^{10} \approx 0.539$ (3).
\item R: complement: $1 - 0.539 = 0.461$ (4).
\item S: need $1 - 0.94^S \ge 0.99 \iff 0.94^S \le 0.01 \iff S \ge \dfrac{\ln 0.01}{\ln 0.94} = \dfrac{-4.605}{-0.0619} \approx 74.4$,
so $S = 75$ (1).
\end{itemize}
Hence P$\to$2, Q$\to$3, R$\to$4, S$\to$1.

\textbf{\color{brand}Why the other options are incorrect.}
\begin{itemize}[leftmargin=2.4em, itemsep=1pt, topsep=2pt]
\item[\textbf{(B)}] Q$\to$4, R$\to$3 swaps an event and its complement: \emph{no} optimal split in 10 shots has probability $0.94^{10} = 0.539$.
\item[\textbf{(C)}] P$\to$1 and S$\to$2 are wrong as in the previous option, and Q/R are also swapped.
\item[\textbf{(D)}] P$\to$1, S$\to$2 swaps the standard deviation and the shot count: with 15 shots the success probability is only $1 - 0.94^{15} = 0.60$.
\end{itemize}
\textbf{\color{brand}Key takeaway.} With $P_{\text{opt}} = 0.06$, about 75 shots already give $99\,\%$ confidence of seeing an optimum.
\end{solbox}
\qsep

\needspace{13\baselineskip}
\begin{qitem}{14}
\qmeta{Single-correct MCQ}{Complementary optimal strings}
The notebook's brute-force search reports four optimal bitstrings (cut value $25$):
\[ \texttt{001101110100010},\quad \texttt{001101110110010},\quad \texttt{110010001001101},\quad \texttt{110010001011101}. \]
The number of distinct optimal partitions, and the bitstring describing the same partition as $\texttt{001101110100010}$, are
\par\smallskip\noindent\begin{tabularx}{\linewidth}{@{}*{2}{L}@{}}
\optlabel{A}\,$2$; $\ \texttt{110010001001101}$ & \optlabel{B}\,$2$; $\ \texttt{010001011101100}$ \\[3pt]
\optlabel{C}\,$4$; $\ \texttt{110010001011101}$ & \correctlabel{D}\,$2$; $\ \texttt{110010001011101}$
\end{tabularx}
\end{qitem}
\begin{solbox}{D}
\textbf{\color{brand}Correct answer.} \textcolor{okgreen}{\textbf{(D)}} $2$; $\ \texttt{110010001011101}$

\textbf{\color{brand}Concepts required.} $C(z) = C(\bar z)$: optimal strings come in complementary pairs.

\textbf{\color{brand}Relevant formulae.} $\bar z_i = 1 - z_i$; $\ [z_i\neq z_j] = [\bar z_i\neq\bar z_j]$; number of optimal partitions $=$ (number of optimal strings)$/2$.

\textbf{\color{brand}Step-by-step explanation.} $C(z) = C(\bar z)$: flipping every bit swaps the group labels and keeps every edge cut or uncut. So optimal strings
come in complementary pairs, and each pair is one partition.
Complement of $\texttt{001101110100010}$: flip each bit $\Rightarrow \texttt{110010001011101}$, which is in the list.
Complement of $\texttt{001101110110010}$: $\texttt{110010001001101}$, also in the list.
So the 4 strings form 2 pairs: \textbf{2 distinct optimal partitions}. (The two partitions differ only in node 10,
which has as many neighbours on each side when the rest is fixed.)

\textbf{\color{brand}Why the other options are incorrect.}
\begin{itemize}[leftmargin=2.4em, itemsep=1pt, topsep=2pt]
\item[\textbf{(A)}] $\texttt{110010001001101}$ is the complement of the \emph{other} string $\texttt{001101110110010}$ (they differ in bit 10).
\item[\textbf{(B)}] $\texttt{010001011101100}$ is the string written backwards (reversal), not the complement.
\item[\textbf{(C)}] $4$ counts bitstrings, not partitions.
\end{itemize}
\textbf{\color{brand}Key takeaway.} Complement = flip every bit; reversal is a different operation and does not preserve cuts.
\end{solbox}
\qsep

\needspace{13\baselineskip}
\begin{qitem}{15}
\qmeta{Single-correct MCQ}{QUBO $\to$ Ising conversion}
The QUBO function $f(x_0, x_1) = 2x_0x_1 - x_0 - x_1$ with $x_i \in \{0,1\}$ is rewritten with spins via
$x_i = (1 - s_i)/2$, $s_i \in \{+1, -1\}$. The result is
\par\smallskip\noindent\begin{tabularx}{\linewidth}{@{}*{2}{L}@{}}
\optlabel{A}\,$\tfrac12\,s_0s_1 - 1 + \tfrac12(s_0 + s_1)$ & \correctlabel{B}\,$\tfrac12\,(s_0s_1 - 1)$ \\[3pt]
\optlabel{C}\,$\tfrac14\,(s_0s_1 - 1)$ & \optlabel{D}\,$\tfrac12\,(1 - s_0s_1)$
\end{tabularx}
\end{qitem}
\begin{solbox}{B}
\textbf{\color{brand}Correct answer.} \textcolor{okgreen}{\textbf{(B)}} $\tfrac12\,(s_0s_1 - 1)$

\textbf{\color{brand}Concepts required.} Change of variables between binary (QUBO) and spin (Ising) form.

\textbf{\color{brand}Relevant formulae.} $x_i = \tfrac{1 - s_i}{2}$, $\ x_ix_j = \tfrac14(1 - s_i - s_j + s_is_j)$, $\ [x_i\neq x_j] = \tfrac12(1 - s_is_j)$.

\textbf{\color{brand}Step-by-step explanation.} Substitute:
$x_0x_1 = \tfrac14(1 - s_0)(1 - s_1) = \tfrac14(1 - s_0 - s_1 + s_0s_1)$, so $2x_0x_1 = \tfrac12(1 - s_0 - s_1 + s_0s_1)$;
$-x_0 - x_1 = -\tfrac12(1 - s_0) - \tfrac12(1 - s_1) = -1 + \tfrac12(s_0 + s_1)$.
Adding, the linear terms cancel:
\[ f = \tfrac12 - \tfrac12(s_0+s_1) + \tfrac12 s_0s_1 - 1 + \tfrac12(s_0+s_1) = \tfrac12(s_0s_1 - 1). \]
Check: $x = 01$: $f = -1$ and $s = (+1, -1)$ gives $\tfrac12(-1-1) = -1$. So $f = -[x_0 \neq x_1]$: minimising $f$ is Max-Cut
on one edge, and its Ising form is (up to a constant) $\tfrac12 Z_0Z_1$ --- one term of $H_C$.

\textbf{\color{brand}Why the other options are incorrect.}
\begin{itemize}[leftmargin=2.4em, itemsep=1pt, topsep=2pt]
\item[\textbf{(A)}] The linear terms do not survive: $-\frac12(s_0+s_1)$ from $2x_0x_1$ cancels $+\frac12(s_0+s_1)$ from $-x_0 - x_1$.
\item[\textbf{(C)}] $\frac14(s_0s_1 - 1)$ forgets the factor 2 of $2x_0x_1$; it gives $f(0,1) = -\frac12$.
\item[\textbf{(D)}] $\frac12(1 - s_0s_1)$ has the wrong overall sign: it equals $+[x_0\neq x_1]$, while $f(0,1) = -1$.
\end{itemize}
\textbf{\color{brand}Key takeaway.} Max-Cut is a QUBO whose Ising form has only $Z_iZ_j$ terms (no fields): exactly the shape of $H_C$.
\end{solbox}
\qsep

\needspace{13\baselineskip}
\begin{qitem}{16}
\qmeta{Single-correct MCQ}{Exact shift rule for derivatives}
Let $f(\theta) = a + b\cos\theta + c\sin\theta$ with constants $a, b, c$. Which expression equals $f'(\theta)$
\emph{exactly} for every $\theta$?
\par\smallskip\noindent\begin{tabularx}{\linewidth}{@{}*{2}{L}@{}}
\optlabel{A}\,$\tfrac12\big[f(\theta + \tfrac\pi4) - f(\theta - \tfrac\pi4)\big]$ & \optlabel{B}\,$f(\theta + \tfrac\pi2) - f(\theta)$ \\[3pt]
\correctlabel{C}\,$\tfrac12\big[f(\theta + \tfrac\pi2) - f(\theta - \tfrac\pi2)\big]$ & \optlabel{D}\,$\tfrac1\pi\big[f(\theta + \tfrac\pi2) - f(\theta - \tfrac\pi2)\big]$
\end{tabularx}
\end{qitem}
\begin{solbox}{C}
\textbf{\color{brand}Correct answer.} \textcolor{okgreen}{\textbf{(C)}} $\tfrac12\big[f(\theta + \tfrac\pi2) - f(\theta - \tfrac\pi2)\big]$

\textbf{\color{brand}Concepts required.} Trigonometric shift identities; expectation values are trigonometric in a single gate angle.

\textbf{\color{brand}Relevant formulae.} $\cos(\theta\pm\tfrac\pi2) = \mp\sin\theta$, $\ \sin(\theta\pm\tfrac\pi2) = \pm\cos\theta$; parameter shift $f' = \tfrac12[f(\theta + \tfrac\pi2) - f(\theta - \tfrac\pi2)]$.

\textbf{\color{brand}Step-by-step explanation.} $f'(\theta) = -b\sin\theta + c\cos\theta$. Using $\cos(\theta \pm \tfrac\pi2) = \mp\sin\theta$ and $\sin(\theta \pm \tfrac\pi2) = \pm\cos\theta$:
\[ f(\theta + \tfrac\pi2) - f(\theta - \tfrac\pi2) = b(-\sin\theta - \sin\theta) + c(\cos\theta + \cos\theta) = 2(-b\sin\theta + c\cos\theta) = 2f'(\theta). \]
So $f'(\theta) = \tfrac12[f(\theta + \tfrac\pi2) - f(\theta - \tfrac\pi2)]$ exactly --- not an approximation. This is the
\emph{parameter-shift rule}: an expectation value depends on the angle of one gate $e^{-i\theta P/2}$ ($P$ a Pauli
operator) exactly in this form, so two circuit runs give the exact derivative.

\textbf{\color{brand}Why the other options are incorrect.}
\begin{itemize}[leftmargin=2.4em, itemsep=1pt, topsep=2pt]
\item[\textbf{(A)}] A shift of $\pi/4$ gives $\frac{1}{\sqrt2}f'$; the factor $\frac12$ is right only for shifts of $\pm\pi/2$.
\item[\textbf{(B)}] $f(\theta + \frac\pi2) - f(\theta) = -b\sin\theta + c\cos\theta - b\cos\theta - c\sin\theta \neq f'$ in general.
\item[\textbf{(D)}] Dividing by $\pi$ treats it as a central finite difference with step $\pi/2$; it gives $\frac2\pi f' \neq f'$.
\end{itemize}
\textbf{\color{brand}Key takeaway.} Large shifts, exact derivative: two circuit runs per gate angle.
\end{solbox}
\qsep

\needspace{20\baselineskip}
\begin{passagebox}{Paragraph for Questions 17 and 18}
The ``house graph'' has the $5$ nodes $0, 1, 2, 3, 4$ and the $6$ unit-weight edges
$(0,1), (1,2), (2,3), (3,0)$ (a square) and $(0,4), (1,4)$ (a roof: node 4 forms a triangle with nodes 0 and 1).
A split is a bitstring $z = z_0z_1z_2z_3z_4$; its cut value $C(z)$ counts the edges whose two end nodes are in different
groups.
\end{passagebox}
\needspace{13\baselineskip}
\begin{qitem}{17}
\qmeta{Paragraph based}{House graph: $C_{\max}$ and count}
The maximum cut $C_{\max}$ of the house graph, and the number of bitstrings $z$ that achieve it, are
\par\smallskip\noindent\begin{tabularx}{\linewidth}{@{}*{4}{L}@{}}
\optlabel{A}\,$5$ and $6$ & \optlabel{B}\,$5$ and $2$ & \optlabel{C}\,$6$ and $2$ & \correctlabel{D}\,$5$ and $4$
\end{tabularx}
\end{qitem}
\begin{solbox}{D}
\textbf{\color{brand}Correct answer.} \textcolor{okgreen}{\textbf{(D)}} $5$ and $4$

\textbf{\color{brand}Concepts required.} Upper bound from an odd cycle; construction of a matching lower bound; counting.

\textbf{\color{brand}Relevant formulae.} Each odd cycle has at least one uncut edge; a bipartite subgraph can be cut completely.

\textbf{\color{brand}Step-by-step explanation.} \textbf{Upper bound:} the triangle $0,1,4$ is an odd cycle, so at least one of its three edges is uncut: $C \le 6 - 1 = 5$.
\textbf{Achieving 5:} the square $0{-}1{-}2{-}3$ is bipartite; put $\{0, 2\}$ and $\{1, 3\}$ in different groups (all 4 square
edges cut). Node 4 is joined to 0 and 1, which are now in different groups, so exactly one roof edge is cut \emph{whichever}
group node 4 joins: total $4 + 1 = 5$.
\textbf{Counting:} a cut of 5 must cut all four square edges (if a square edge were uncut, at most $3 + 2 = 5$ with both roof
edges cut, which needs $z_4 \neq z_0$ and $z_4 \neq z_1$, so $z_0 = z_1$ --- but then edge $(0,1)$ is uncut too, giving at most 4).
So $z_0z_1z_2z_3 \in \{0101, 1010\}$ and $z_4$ is free: $2 \times 2 = 4$ strings: $01010, 01011, 10100, 10101$.

\textbf{\color{brand}Why the other options are incorrect.}
\begin{itemize}[leftmargin=2.4em, itemsep=1pt, topsep=2pt]
\item[\textbf{(A)}] Six strings would require more optimal patterns of the square; only $0101$ and $1010$ cut all four square edges.
\item[\textbf{(B)}] $2$ counts partitions ($\{0,2,4\}|\{1,3\}$ and $\{0,2\}|\{1,3,4\}$) instead of bitstrings.
\item[\textbf{(C)}] $6$ would cut every edge, impossible because of the triangle $0, 1, 4$.
\end{itemize}
\textbf{\color{brand}Key takeaway.} Bound with odd cycles, then construct; node 4 is free because its two neighbours are always on opposite sides.
\end{solbox}
\qsep

\needspace{13\baselineskip}
\begin{qitem}{18}
\qmeta{Paragraph based}{House graph: random baseline}
A split of the house graph is drawn uniformly at random. The probability that it is optimal, and the ratio
$\mathbb{E}[C]/C_{\max}$ of the random split, are
\par\smallskip\noindent\begin{tabularx}{\linewidth}{@{}*{4}{L}@{}}
\optlabel{A}\,$\tfrac1{16}$ and $0.6$ & \optlabel{B}\,$\tfrac14$ and $0.6$ & \optlabel{C}\,$\tfrac18$ and $0.5$ & \correctlabel{D}\,$\tfrac18$ and $0.6$
\end{tabularx}
\end{qitem}
\begin{solbox}{D}
\textbf{\color{brand}Correct answer.} \textcolor{okgreen}{\textbf{(D)}} $\tfrac18$ and $0.6$

\textbf{\color{brand}Concepts required.} Uniform distribution over $2^n$ bitstrings; random baseline $W/2$.

\textbf{\color{brand}Relevant formulae.} $P(\text{optimal}) = \#\text{optimal strings}/2^n$; $\ r_{\text{random}} = (W/2)/C_{\max}$.

\textbf{\color{brand}Step-by-step explanation.} There are $2^5 = 32$ equally likely bitstrings and $4$ are optimal (previous question): $P = \tfrac{4}{32} = \tfrac18$.
The expected cut of a uniform split is $W/2 = 6/2 = 3$ (every edge is cut with probability $\tfrac12$), so the random
approximation ratio is $3/5 = 0.6$. Any useful algorithm --- QAOA included --- must beat $r = 0.6$ on this graph.

\textbf{\color{brand}Why the other options are incorrect.}
\begin{itemize}[leftmargin=2.4em, itemsep=1pt, topsep=2pt]
\item[\textbf{(A)}] $\frac1{16}$ counts 2 optimal strings (partitions) instead of 4.
\item[\textbf{(B)}] $\frac14$ divides by $16 = 2^4$, forgetting one node.
\item[\textbf{(C)}] $0.5$ divides $W/2 = 3$ by $|E| = 6$ instead of by $C_{\max} = 5$.
\end{itemize}
\textbf{\color{brand}Key takeaway.} Random guessing already achieves $r = (W/2)/C_{\max}$, here $0.6$.
\end{solbox}
\qsep

\needspace{13\baselineskip}
\begin{qitem}{19}
\qmeta{Single-correct MCQ}{Ground energy $= W/2 - C_{\max}$}
For the Petersen graph (10 nodes, 15 unit edges, $C_{\max} = 12$), the smallest eigenvalue of
$H_C = \sum_{(i,j)\in E}\tfrac12 Z_iZ_j$ is
\par\smallskip\noindent\begin{tabularx}{\linewidth}{@{}*{4}{L}@{}}
\optlabel{A}\,$-7.5$ & \correctlabel{B}\,$-4.5$ & \optlabel{C}\,$-12$ & \optlabel{D}\,$+4.5$
\end{tabularx}
\end{qitem}
\begin{solbox}{B}
\textbf{\color{brand}Correct answer.} \textcolor{okgreen}{\textbf{(B)}} $-4.5$

\textbf{\color{brand}Concepts required.} Relation between the cut and the Ising energy on basis states.

\textbf{\color{brand}Relevant formulae.} $C(z) = \tfrac W2 - H_C(z)$ $\Rightarrow$ $E_0(H_C) = \tfrac W2 - C_{\max}$.

\textbf{\color{brand}Step-by-step explanation.} On a basis state $\ket z$, $C(z) = \tfrac W2 - H_C(z)$ with $W = 15$, because
$[z_i \neq z_j] = \tfrac12(1 - s_is_j)$ and summing gives $C = \tfrac{|E|}{2} - \tfrac12\sum s_is_j$.
Minimising $H_C$ is maximising $C$: $E_0 = \tfrac W2 - C_{\max} = 7.5 - 12 = -4.5$.
$-7.5$ would require all 15 products $s_is_j = -1$ (every edge cut), impossible because the graph has odd cycles.

\textbf{\color{brand}Why the other options are incorrect.}
\begin{itemize}[leftmargin=2.4em, itemsep=1pt, topsep=2pt]
\item[\textbf{(A)}] $-7.5$ assumes all 15 edges can be cut; the Petersen graph has 5-cycles.
\item[\textbf{(C)}] $-12$ confuses the energy with the cut value; the shift by $W/2$ is missing.
\item[\textbf{(D)}] $+4.5$ has the wrong sign: large cuts correspond to \emph{low} energy.
\end{itemize}
\textbf{\color{brand}Key takeaway.} Ground energy $= W/2 - C_{\max}$: minimising $H_C$ and maximising $C$ are the same problem.
\end{solbox}
\qsep

\needspace{13\baselineskip}
\begin{qitem}{20}
\qmeta{Assertion--Reason}{Local search: $|E|/2$, not optimal}
The notebook's \texttt{local\_search\_maxcut} repeatedly flips single nodes while this strictly increases the cut
(unit weights), starting from random splits.\\[1mm]
\textbf{Assertion (A):} Every split returned by one run of this greedy procedure is an optimal split.\\[1mm]
\textbf{Reason (R):} When no single flip increases the cut, every node has at least half of its incident edges cut.
\begin{itemize}[leftmargin=2.6em, labelsep=0.4em, itemsep=2pt, topsep=3pt]
\item[\optlabel{A}] Both Assertion (A) and Reason (R) are true, and R is the correct explanation of A
\item[\optlabel{B}] Both A and R are true, but R is \emph{not} the correct explanation of A
\item[\optlabel{C}] A is true, but R is false
\item[\correctlabel{D}] A is false, but R is true
\end{itemize}
\end{qitem}
\begin{solbox}{D}
\textbf{\color{brand}Correct answer.} \textcolor{okgreen}{\textbf{(D)}} A is false, but R is true

\textbf{\color{brand}Concepts required.} Local versus global optimum; the $|E|/2$ guarantee of single-flip local search.

\textbf{\color{brand}Relevant formulae.} Flip gain at node $v$: $\Delta C = u_v - c_v$; local optimum $\Rightarrow c_v \ge d_v/2$ $\Rightarrow C \ge |E|/2$.

\textbf{\color{brand}Step-by-step explanation.} \textbf{R is true:} flipping node $v$ turns its cut edges into uncut ones and vice versa, so the cut changes by
$u_v - c_v$ ($u_v$, $c_v$ = numbers of uncut and cut edges at $v$). At a local optimum no flip strictly improves, so
$u_v - c_v \le 0$, i.e.\ $c_v \ge d_v/2$. Summing over $v$ (each cut edge counted twice): $2C \ge \sum_v d_v/2 = |E|$, so a
local optimum always has $C \ge |E|/2$.
\textbf{A is false:} a local optimum need not be global --- greedy search can get stuck. That is exactly why the notebook
restarts it 200 times and labels its result ``heuristic local search (best known, NOT guaranteed optimal)''. Answer: (D).

\textbf{\color{brand}Why the other options are incorrect.}
\begin{itemize}[leftmargin=2.4em, itemsep=1pt, topsep=2pt]
\item[\textbf{(A)}] needs A true, R true \emph{and} R explaining A; A is false: a local optimum need not be global.
\item[\textbf{(B)}] needs A and R both true but R not explaining A; A is false.
\item[\textbf{(C)}] needs A true and R false; A is false and R is true.
\end{itemize}
\textbf{\color{brand}Key takeaway.} Greedy local search guarantees only $|E|/2$; restarts help but prove nothing.
\end{solbox}
\qsep

\newpage

\begin{sectionbanner}SECTION B \enspace$\cdot$\enspace Quantum Mechanics \hfill Q.21--40\end{sectionbanner}

\needspace{13\baselineskip}
\begin{qitem}{21}
\qmeta{Single-correct MCQ}{Uniform state; QAOA enhancement}
The 15-qubit state $\ket{+}^{\otimes 15}$ is measured in the computational basis. The notebook's graph has 4 optimal
bitstrings. The probability of obtaining an optimal bitstring, and the factor by which the previously recorded QAOA value
$P_{\text{opt}} = 0.06$ exceeds it, are
\par\smallskip\noindent\begin{tabularx}{\linewidth}{@{}*{2}{L}@{}}
\optlabel{A}\,$6.10\times10^{-5}$; about $980$ & \optlabel{B}\,$3.05\times10^{-5}$; about $1970$ \\[3pt]
\optlabel{C}\,$2.44\times10^{-4}$; about $245$ & \correctlabel{D}\,$1.22\times10^{-4}$; about $490$
\end{tabularx}
\end{qitem}
\begin{solbox}{D}
\textbf{\color{brand}Correct answer.} \textcolor{okgreen}{\textbf{(D)}} $1.22\times10^{-4}$; about $490$

\textbf{\color{brand}Concepts required.} Born rule; uniform superposition; probability of a set of outcomes.

\textbf{\color{brand}Relevant formulae.} $\ket+^{\otimes n} = 2^{-n/2}\sum_z\ket z$, $\ P(z) = 2^{-n}$; $\ P(\text{set}) = |\text{set}|/2^n$.

\textbf{\color{brand}Step-by-step explanation.} $\ket{+}^{\otimes n} = 2^{-n/2}\sum_z \ket z$: every one of the $2^{15} = 32\,768$ bitstrings has amplitude $2^{-15/2}$ and
probability $|2^{-15/2}|^2 = 1/32\,768$ (Born rule). Four optimal strings: $P = 4/32\,768 = 1.22\times10^{-4}$.
Enhancement: $0.06 / (1.22\times10^{-4}) = 0.06 \times 8192 \approx 491.5$. QAOA concentrates probability on good cuts;
the initial uniform superposition is simply random guessing.

\textbf{\color{brand}Why the other options are incorrect.}
\begin{itemize}[leftmargin=2.4em, itemsep=1pt, topsep=2pt]
\item[\textbf{(A)}] $6.10\times10^{-5}$ counts 2 partitions as 2 bitstrings; there are 4 optimal bitstrings.
\item[\textbf{(B)}] $3.05\times10^{-5} = 1/32\,768$ counts only one optimal string, not four.
\item[\textbf{(C)}] $2.44\times10^{-4}$ divides by $2^{14}$ (partitions); measurements return bitstrings, of which there are $2^{15}$.
\end{itemize}
\textbf{\color{brand}Key takeaway.} $\ket+^{\otimes n}$ is random guessing; QAOA's job is to raise $P_{\text{opt}}$ far above $|\text{opt}|/2^n$.
\end{solbox}
\qsep

\needspace{13\baselineskip}
\begin{qitem}{22}
\qmeta{Single-correct MCQ}{Pauli products}
The two-qubit operator $(Z\otimes Z)(X\otimes X)$ equals
\par\smallskip\noindent\begin{tabularx}{\linewidth}{@{}*{4}{L}@{}}
\optlabel{A}\,$i\,Y\otimes Y$ & \correctlabel{B}\,$-\,Y\otimes Y$ & \optlabel{C}\,$Y\otimes Y$ & \optlabel{D}\,$-i\,Y\otimes Y$
\end{tabularx}
\end{qitem}
\begin{solbox}{B}
\textbf{\color{brand}Correct answer.} \textcolor{okgreen}{\textbf{(B)}} $-\,Y\otimes Y$

\textbf{\color{brand}Concepts required.} Pauli multiplication table; tensor products multiply factor-wise.

\textbf{\color{brand}Relevant formulae.} $ZX = iY$, $XZ = -iY$, $XY = iZ$, $YZ = iX$; $\ (A\otimes B)(C\otimes D) = AC\otimes BD$.

\textbf{\color{brand}Step-by-step explanation.} Tensor products multiply factor by factor: $(Z\otimes Z)(X\otimes X) = (ZX)\otimes(ZX)$.
Single-qubit Pauli algebra: $ZX = iY$ (check: $ZX = \begin{psmallmatrix}1&0\\0&-1\end{psmallmatrix}\begin{psmallmatrix}0&1\\1&0\end{psmallmatrix}
= \begin{psmallmatrix}0&1\\-1&0\end{psmallmatrix} = i\begin{psmallmatrix}0&-i\\i&0\end{psmallmatrix} = iY$).
So $(iY)\otimes(iY) = i^2\,Y\otimes Y = -Y\otimes Y$. Also $(X\otimes X)(Z\otimes Z) = (XZ)\otimes(XZ) = (-iY)\otimes(-iY) = -Y\otimes Y$:
the same --- $Z\otimes Z$ and $X\otimes X$ commute, although $Z$ and $X$ anticommute.

\textbf{\color{brand}Why the other options are incorrect.}
\begin{itemize}[leftmargin=2.4em, itemsep=1pt, topsep=2pt]
\item[\textbf{(A)}] $iY\otimes Y$ keeps only one factor $i$ (two qubits give $i\cdot i$).
\item[\textbf{(C)}] $Y\otimes Y$ forgets the factor $i^2 = -1$.
\item[\textbf{(D)}] $-iY\otimes Y$ mixes $ZX = iY$ on one qubit with $XZ = -iY$ on the other; both factors are $ZX$.
\end{itemize}
\textbf{\color{brand}Key takeaway.} Count the phases qubit by qubit; two anticommuting positions make the products commute.
\end{solbox}
\qsep

\needspace{13\baselineskip}
\begin{qitem}{23}
\qmeta{Matching lists}{Rotation-gate matrices}
Match each operator (List-I) with its matrix (List-II). Two-qubit matrices use the basis order $\ket{00}, \ket{01},
\ket{10}, \ket{11}$.
\par\smallskip\noindent{\renewcommand{\arraystretch}{1.35}%
\begin{tabularx}{\linewidth}{@{}c>{\raggedright\arraybackslash}p{0.33\linewidth}@{\hspace{1.2em}}cL@{}}
\toprule
\multicolumn{2}{@{}l}{\textbf{List-I}} & \multicolumn{2}{l@{}}{\textbf{List-II}} \\ \midrule
\textbf{(P)} & $R_X(2\beta)$ & \textbf{(1)} & $\operatorname{diag}\big(e^{-i\theta/2}, e^{i\theta/2}, e^{i\theta/2}, e^{-i\theta/2}\big)$ \\
\textbf{(Q)} & $R_Z(\theta)$ & \textbf{(2)} & $\begin{pmatrix}\cos\beta & -i\sin\beta\\ -i\sin\beta & \cos\beta\end{pmatrix}$ \\
\textbf{(R)} & $e^{-i\theta\, Z\otimes Z/2}$ & \textbf{(3)} & $\tfrac{1}{\sqrt2}\begin{pmatrix}1 & 1\\ 1 & -1\end{pmatrix}$ \\
\textbf{(S)} & $H$ & \textbf{(4)} & $\operatorname{diag}\big(e^{-i\theta/2}, e^{i\theta/2}\big)$ \\
\bottomrule
\end{tabularx}}
\par\smallskip\noindent\begin{tabularx}{\linewidth}{@{}*{2}{L}@{}}
\optlabel{A}\,P$\to$3, Q$\to$4, R$\to$1, S$\to$2 & \optlabel{B}\,P$\to$2, Q$\to$1, R$\to$4, S$\to$3 \\[3pt]
\optlabel{C}\,P$\to$4, Q$\to$2, R$\to$1, S$\to$3 & \correctlabel{D}\,P$\to$2, Q$\to$4, R$\to$1, S$\to$3
\end{tabularx}
\end{qitem}
\begin{solbox}{D}
\textbf{\color{brand}Correct answer.} \textcolor{okgreen}{\textbf{(D)}} P$\to$2, Q$\to$4, R$\to$1, S$\to$3

\textbf{\color{brand}Concepts required.} Exponentials of Pauli operators; half-angle convention of rotation gates.

\textbf{\color{brand}Relevant formulae.} $e^{-i\phi P} = \cos\phi\,I - i\sin\phi\,P$ ($P^2 = I$); $R_X(\theta) = e^{-i\theta X/2}$, $R_Z(\theta) = e^{-i\theta Z/2}$; $Z\otimes Z = \operatorname{diag}(1,-1,-1,1)$.

\textbf{\color{brand}Step-by-step explanation.} For a Pauli $P$ (so $P^2 = I$): $e^{-i\phi P} = \cos\phi\, I - i\sin\phi\, P$.
\begin{itemize}\itemsep0pt
\item P: $R_X(2\beta) = e^{-i\beta X} = \cos\beta\, I - i\sin\beta\, X$ (2).
\item Q: $R_Z(\theta) = e^{-i\theta Z/2} = \operatorname{diag}(e^{-i\theta/2}, e^{+i\theta/2})$ (4).
\item R: $Z\otimes Z = \operatorname{diag}(+1, -1, -1, +1)$, so $e^{-i\theta Z\otimes Z/2} = \operatorname{diag}(e^{-i\theta/2}, e^{i\theta/2}, e^{i\theta/2}, e^{-i\theta/2})$ (1):
equal bits get $e^{-i\theta/2}$, different bits get $e^{+i\theta/2}$.
\item S: Hadamard (3).
\end{itemize}

\textbf{\color{brand}Why the other options are incorrect.}
\begin{itemize}[leftmargin=2.4em, itemsep=1pt, topsep=2pt]
\item[\textbf{(A)}] P$\to$3 assigns $H$ to $R_X(2\beta)$; no value of $\beta$ gives $H$ (e.g.\ $\beta = \pi/4$ gives $\frac{1}{\sqrt2}\begin{psmallmatrix}1&-i\\-i&1\end{psmallmatrix}$).
\item[\textbf{(B)}] Q$\to$1, R$\to$4 swaps the one-qubit $R_Z$ ($2\times2$) and the two-qubit $ZZ$ rotation ($4\times4$).
\item[\textbf{(C)}] P$\to$4, Q$\to$2 swaps $R_X$ and $R_Z$: an $X$ rotation has off-diagonal entries, a $Z$ rotation is diagonal.
\end{itemize}
\textbf{\color{brand}Key takeaway.} $Z$-type rotations are diagonal (phases), $X$-type rotations mix $\ket0$ and $\ket1$.
\end{solbox}
\qsep

\needspace{13\baselineskip}
\begin{qitem}{24}
\qmeta{Single-correct MCQ}{Relative phase and interference}
The state $R_Z(\theta)\ket{+}$ is passed through a Hadamard gate and then measured in the computational basis. The
probability of the outcome 0 is
\par\smallskip\noindent\begin{tabularx}{\linewidth}{@{}*{4}{L}@{}}
\correctlabel{A}\,$\cos^2(\theta/2)$ & \optlabel{B}\,$\sin^2(\theta/2)$ & \optlabel{C}\,$\cos^2\theta$ & \optlabel{D}\,$\tfrac12$
\end{tabularx}
\end{qitem}
\begin{solbox}{A}
\textbf{\color{brand}Correct answer.} \textcolor{okgreen}{\textbf{(A)}} $\cos^2(\theta/2)$

\textbf{\color{brand}Concepts required.} Relative phase is invisible in the $Z$ basis but becomes a probability after a Hadamard (interference).

\textbf{\color{brand}Relevant formulae.} $H\tfrac{1}{\sqrt2}(\ket0 + e^{i\theta}\ket1) = \tfrac12[(1 + e^{i\theta})\ket0 + (1 - e^{i\theta})\ket1]$; $|1 + e^{i\theta}|^2 = 4\cos^2\tfrac\theta2$.

\textbf{\color{brand}Step-by-step explanation.} $R_Z(\theta)\ket+ = \tfrac{1}{\sqrt2}(e^{-i\theta/2}\ket0 + e^{i\theta/2}\ket1) = e^{-i\theta/2}\,\tfrac{1}{\sqrt2}(\ket0 + e^{i\theta}\ket1)$.
Apply $H$: $\tfrac{1}{\sqrt2}(\ket0 + e^{i\theta}\ket1) \mapsto \tfrac12\big[(1 + e^{i\theta})\ket0 + (1 - e^{i\theta})\ket1\big]$.
$P(0) = \tfrac14|1 + e^{i\theta}|^2 = \tfrac14(2 + 2\cos\theta) = \cos^2(\theta/2)$.
Without the Hadamard the result would be $\tfrac12$ for every $\theta$: a $Z$-rotation only changes a \emph{relative phase},
which becomes visible only after interference. In QAOA the mixer plays this interfering role for the cost phases.

\textbf{\color{brand}Why the other options are incorrect.}
\begin{itemize}[leftmargin=2.4em, itemsep=1pt, topsep=2pt]
\item[\textbf{(B)}] $\sin^2(\theta/2)$ is $P(1)$, not $P(0)$.
\item[\textbf{(C)}] $\cos^2\theta$ doubles the angle: the relative phase is $\theta$, and the probability is $\cos^2(\theta/2)$.
\item[\textbf{(D)}] $\frac12$ is the result \emph{without} the Hadamard.
\end{itemize}
\textbf{\color{brand}Key takeaway.} Phases need interference to be seen --- in QAOA the mixer does this.
\end{solbox}
\qsep

\needspace{13\baselineskip}
\begin{qitem}{25}
\qmeta{Single-correct MCQ}{Expected cut from amplitudes}
For one unit edge $(0,1)$ the qubits are in the state $\ket\psi = \tfrac12\ket{00} + \tfrac{\sqrt3}{2}\ket{01}$. The expected
cut value $\langle C\rangle = \big\langle \tfrac12(1 - Z_0Z_1)\big\rangle$ is
\par\smallskip\noindent\begin{tabularx}{\linewidth}{@{}*{4}{L}@{}}
\optlabel{A}\,$\tfrac14$ & \optlabel{B}\,$-\tfrac12$ & \optlabel{C}\,$\tfrac{\sqrt3}{2}$ & \correctlabel{D}\,$\tfrac34$
\end{tabularx}
\end{qitem}
\begin{solbox}{D}
\textbf{\color{brand}Correct answer.} \textcolor{okgreen}{\textbf{(D)}} $\tfrac34$

\textbf{\color{brand}Concepts required.} Expectation of a diagonal observable from the Born probabilities.

\textbf{\color{brand}Relevant formulae.} $\langle Z_0Z_1\rangle = \sum_z P(z)\,s_0s_1$; $\ \langle C_{01}\rangle = \tfrac12(1 - \langle Z_0Z_1\rangle) = P(z_0\neq z_1)$.

\textbf{\color{brand}Step-by-step explanation.} $Z_0Z_1$ is diagonal with eigenvalue $+1$ on $\ket{00}$ and $-1$ on $\ket{01}$. Born rule: $P(00) = \tfrac14$, $P(01) = \tfrac34$.
$\langle Z_0Z_1\rangle = \tfrac14(+1) + \tfrac34(-1) = -\tfrac12$, so $\langle C\rangle = \tfrac12\big(1 - (-\tfrac12)\big) = \tfrac34$.
Directly: the edge is cut exactly in outcome $01$, which has probability $\tfrac34$.

\textbf{\color{brand}Why the other options are incorrect.}
\begin{itemize}[leftmargin=2.4em, itemsep=1pt, topsep=2pt]
\item[\textbf{(A)}] $\frac14$ is $P(00)$, the probability that the edge is \emph{not} cut.
\item[\textbf{(B)}] $-\frac12$ is $\langle Z_0Z_1\rangle$, not the cut.
\item[\textbf{(C)}] $\frac{\sqrt3}{2}$ is an amplitude; probabilities are squared moduli.
\end{itemize}
\textbf{\color{brand}Key takeaway.} $\langle C\rangle$ for an edge is just the probability that the edge is cut.
\end{solbox}
\qsep

\needspace{13\baselineskip}
\begin{qitem}{26}
\qmeta{Single-correct MCQ}{Commutation of Pauli strings}
Which of the following operators commutes with $Z_0Z_1$ (on two or more qubits)?
\par\smallskip\noindent\begin{tabularx}{\linewidth}{@{}*{4}{L}@{}}
\correctlabel{A}\,$X_0X_1$ & \optlabel{B}\,$Y_0Z_1$ & \optlabel{C}\,$X_0$ & \optlabel{D}\,$X_1$
\end{tabularx}
\end{qitem}
\begin{solbox}{A}
\textbf{\color{brand}Correct answer.} \textcolor{okgreen}{\textbf{(A)}} $X_0X_1$

\textbf{\color{brand}Concepts required.} Commutation of Pauli strings: count anticommuting positions.

\textbf{\color{brand}Relevant formulae.} $PQ = (-1)^{k}QP$, $k$ = number of positions where the single-qubit factors anticommute.

\textbf{\color{brand}Step-by-step explanation.} For Pauli products, count the positions where the two factors anticommute: an \emph{even} count means the products
commute, odd means they anticommute.
$X_0X_1$ vs $Z_0Z_1$: positions 0 ($X$ vs $Z$, anticommute) and 1 ($X$ vs $Z$, anticommute) --- two, so they commute.
$X_0$: one anticommuting position. $Y_0Z_1$: position 0 anticommutes ($Y$ vs $Z$), position 1 commutes --- one.
$X_1$: one. So only $X_0X_1$ commutes with $Z_0Z_1$ (consistent with Q.22).

\textbf{\color{brand}Why the other options are incorrect.}
\begin{itemize}[leftmargin=2.4em, itemsep=1pt, topsep=2pt]
\item[\textbf{(B)}] $Y_0Z_1$: $Y_0$ vs $Z_0$ anticommute, $Z_1$ vs $Z_1$ commute --- one position, so they anticommute.
\item[\textbf{(C)}] $X_0$ anticommutes with $Z_0$ at one position.
\item[\textbf{(D)}] $X_1$ anticommutes with $Z_1$ at one position.
\end{itemize}
\textbf{\color{brand}Key takeaway.} Even number of clashes $\Rightarrow$ commute; this is why the mixer does not commute with the cost terms.
\end{solbox}
\qsep

\needspace{13\baselineskip}
\begin{qitem}{27}
\qmeta{Matching lists}{Pauli propagation through CNOT}
Conjugation by a CNOT with control qubit 0 and target qubit 1 maps a Pauli operator $P$ to
$\mathrm{CNOT}\,P\,\mathrm{CNOT}$. Match each $P$ (List-I) with its image (List-II).
\par\smallskip\noindent{\renewcommand{\arraystretch}{1.35}%
\begin{tabularx}{\linewidth}{@{}c>{\raggedright\arraybackslash}p{0.33\linewidth}@{\hspace{1.2em}}cL@{}}
\toprule
\multicolumn{2}{@{}l}{\textbf{List-I}} & \multicolumn{2}{l@{}}{\textbf{List-II}} \\ \midrule
\textbf{(P)} & $X\otimes I$ & \textbf{(1)} & $Z\otimes Z$ \\
\textbf{(Q)} & $I\otimes X$ & \textbf{(2)} & $X\otimes X$ \\
\textbf{(R)} & $Z\otimes I$ & \textbf{(3)} & $I\otimes X$ \\
\textbf{(S)} & $I\otimes Z$ & \textbf{(4)} & $Z\otimes I$ \\
\bottomrule
\end{tabularx}}
\par\smallskip\noindent\begin{tabularx}{\linewidth}{@{}*{2}{L}@{}}
\optlabel{A}\,P$\to$3, Q$\to$2, R$\to$1, S$\to$4 & \optlabel{B}\,P$\to$2, Q$\to$3, R$\to$1, S$\to$4 \\[3pt]
\optlabel{C}\,P$\to$3, Q$\to$2, R$\to$4, S$\to$1 & \correctlabel{D}\,P$\to$2, Q$\to$3, R$\to$4, S$\to$1
\end{tabularx}
\end{qitem}
\begin{solbox}{D}
\textbf{\color{brand}Correct answer.} \textcolor{okgreen}{\textbf{(D)}} P$\to$2, Q$\to$3, R$\to$4, S$\to$1

\textbf{\color{brand}Concepts required.} Propagation of Pauli operators through a CNOT (Clifford conjugation).

\textbf{\color{brand}Relevant formulae.} CNOT$_{0\to1}$: $X_0 \to X_0X_1$, $X_1 \to X_1$, $Z_0 \to Z_0$, $Z_1 \to Z_0Z_1$.

\textbf{\color{brand}Step-by-step explanation.} Rules (CNOT$^2 = I$, so conjugation is $P \mapsto \mathrm{CNOT}\,P\,\mathrm{CNOT}$):
\begin{itemize}\itemsep0pt
\item $X$ on the control \emph{spreads} to the target: $X\otimes I \to X\otimes X$ (2). (Flipping the control first changes whether the target is flipped.)
\item $X$ on the target is unchanged: $I\otimes X \to I\otimes X$ (3) (the target flip commutes with ``flip if control is 1'').
\item $Z$ on the control is unchanged: $Z\otimes I \to Z\otimes I$ (4) (the control value is never changed).
\item $Z$ on the target spreads \emph{back} to the control: $I\otimes Z \to Z\otimes Z$ (1): after the CNOT the target holds $z_0\oplus z_1$.
\end{itemize}
The last rule is the heart of the cost gate: CNOT$\cdot(I\otimes R_Z(\theta))\cdot$CNOT $= e^{-i\theta\,Z\otimes Z/2}$.

\textbf{\color{brand}Why the other options are incorrect.}
\begin{itemize}[leftmargin=2.4em, itemsep=1pt, topsep=2pt]
\item[\textbf{(A)}] P$\to$3, Q$\to$2, R$\to$1, S$\to$4 reverses both propagation directions: it lets $X$ move from target to control and $Z$ from control to target (e.g.\ it changes $Z\otimes I$, although the CNOT never alters the control bit).
\item[\textbf{(B)}] R$\to$1, S$\to$4 lets $Z$ spread from the control to the target; $Z$ spreads the other way (target $\to$ control).
\item[\textbf{(C)}] P$\to$3, Q$\to$2 lets $X$ spread from the target to the control; $X$ spreads from control to target.
\end{itemize}
\textbf{\color{brand}Key takeaway.} $X$ flows control $\to$ target, $Z$ flows target $\to$ control; CNOT-$R_Z$-CNOT is a $ZZ$ rotation.
\end{solbox}
\qsep

\needspace{13\baselineskip}
\begin{qitem}{28}
\qmeta{Single-correct MCQ}{Entanglement by the cost gate}
The state $e^{-i\gamma\, Z\otimes Z/2}\ket{+}\ket{+}$ is maximally entangled for which $\gamma \in (0, \pi)$?
\par\smallskip\noindent\begin{tabularx}{\linewidth}{@{}*{2}{L}@{}}
\correctlabel{A}\,$\gamma = \pi/2$ & \optlabel{B}\,$\gamma = 3\pi/4$ \\[3pt]
\optlabel{C}\,$\gamma = \pi/4$ & \optlabel{D}\,for no $\gamma$: it is always a product state
\end{tabularx}
\end{qitem}
\begin{solbox}{A}
\textbf{\color{brand}Correct answer.} \textcolor{okgreen}{\textbf{(A)}} $\gamma = \pi/2$

\textbf{\color{brand}Concepts required.} Product-state test and concurrence of a two-qubit pure state.

\textbf{\color{brand}Relevant formulae.} $\ket\psi = \sum c_{ab}\ket{ab}$ is a product state $\iff c_{00}c_{11} = c_{01}c_{10}$; concurrence $\mathcal C = 2|c_{00}c_{11} - c_{01}c_{10}|$.

\textbf{\color{brand}Step-by-step explanation.} Amplitudes: $c_{00} = c_{11} = \tfrac12 e^{-i\gamma/2}$, $c_{01} = c_{10} = \tfrac12 e^{i\gamma/2}$.
A two-qubit pure state is a product state iff $c_{00}c_{11} - c_{01}c_{10} = 0$; the concurrence is
$\mathcal C = 2|c_{00}c_{11} - c_{01}c_{10}| = 2\cdot\tfrac14|e^{-i\gamma} - e^{i\gamma}| = |\sin\gamma|$.
$\mathcal C = 1$ (maximal) requires $\sin\gamma = \pm1$, i.e.\ $\gamma = \pi/2$ in $(0, \pi)$; at $\gamma = \pi/4$ or $3\pi/4$, $\mathcal C = 0.71$.
The cost layer is the only entangling part of QAOA; the mixer acts on single qubits.

\textbf{\color{brand}Why the other options are incorrect.}
\begin{itemize}[leftmargin=2.4em, itemsep=1pt, topsep=2pt]
\item[\textbf{(B)}] At $3\pi/4$, also $\mathcal C \approx 0.71$.
\item[\textbf{(C)}] At $\gamma = \pi/4$, $\mathcal C = \sin(\pi/4) \approx 0.71$: entangled, but not maximally.
\item[\textbf{(D)}] For $\gamma \notin \pi\mathbb Z$, $c_{00}c_{11} = \frac14 e^{-i\gamma} \neq \frac14 e^{i\gamma} = c_{01}c_{10}$: the state is entangled.
\end{itemize}
\textbf{\color{brand}Key takeaway.} Cost layers entangle (most strongly at $\gamma w = \pi/2$); mixers do not.
\end{solbox}
\qsep

\needspace{13\baselineskip}
\begin{qitem}{29}
\qmeta{Statement I / II}{Global phase vs.\ sign of $\gamma$}
\textbf{Statement I:} Replacing $H_C$ by $H_C + c\,I$ (any real constant $c$) changes the QAOA state
$\ket{\boldsymbol\gamma, \boldsymbol\beta}$ only by a global phase.\\[1mm]
\textbf{Statement II:} Consequently, a QAOA circuit that uses $e^{-i\gamma C}$ with $C = \tfrac W2 - H_C$ has the same optimal
angles $(\gamma^*, \beta^*)$ as the notebook's circuit with $e^{-i\gamma H_C}$.
\begin{itemize}[leftmargin=2.6em, labelsep=0.4em, itemsep=2pt, topsep=3pt]
\item[\optlabel{A}] Both Statement I and Statement II are true
\item[\optlabel{B}] Both Statement I and Statement II are false
\item[\correctlabel{C}] Statement I is true, but Statement II is false
\item[\optlabel{D}] Statement I is false, but Statement II is true
\end{itemize}
\end{qitem}
\begin{solbox}{C}
\textbf{\color{brand}Correct answer.} \textcolor{okgreen}{\textbf{(C)}} Statement I is true, but Statement II is false

\textbf{\color{brand}Concepts required.} Global phases are unobservable; the sign of the generator matters.

\textbf{\color{brand}Relevant formulae.} $e^{-i\gamma(H + cI)} = e^{-i\gamma c}e^{-i\gamma H}$; $\ C = \tfrac W2 I - H_C$ $\Rightarrow$ $e^{-i\gamma C} \simeq e^{+i\gamma H_C}$.

\textbf{\color{brand}Step-by-step explanation.} \textbf{I is true:} $cI$ commutes with everything, so $e^{-i\gamma(H_C + cI)} = e^{-i\gamma c}\,e^{-i\gamma H_C}$; the factors
$e^{-i\gamma_l c}$ multiply the final state by an unobservable global phase.
\textbf{II is false:} $C = \tfrac W2 I - H_C$ contains $-H_C$, so $e^{-i\gamma C} = e^{-i\gamma W/2}\,e^{+i\gamma H_C}$. Apart from a global phase
this is the notebook's cost unitary with $\gamma \to -\gamma$. The optimal angles are related by $\gamma^* \to -\gamma^*$ (the
constant is harmless, the \emph{sign} is not). This is why textbook formulas (written with $e^{-i\gamma C}$) must be translated
before they are used with the notebook's convention. Answer: (C).

\textbf{\color{brand}Why the other options are incorrect.}
\begin{itemize}[leftmargin=2.4em, itemsep=1pt, topsep=2pt]
\item[\textbf{(A)}] needs both statements true; II is false: the cut operator contains $-H_C$, which reverses the sign of $\gamma$.
\item[\textbf{(B)}] needs both statements false; I is true.
\item[\textbf{(D)}] needs I false and II true; I is true and II is false.
\end{itemize}
\textbf{\color{brand}Key takeaway.} Constants are harmless, signs are not: $\gamma_{\text{textbook}} = -\gamma_{\text{notebook}}$.
\end{solbox}
\qsep

\needspace{13\baselineskip}
\begin{qitem}{30}
\qmeta{Single-correct MCQ}{$\langle Z_0Z_1\rangle$ for a product state}
Two qubits are in a product state in which $P(z_0 = 1) = 0.2$ and $P(z_1 = 1) = 0.7$. The value of $\langle Z_0Z_1\rangle$ is
\par\smallskip\noindent\begin{tabularx}{\linewidth}{@{}*{4}{L}@{}}
\optlabel{A}\,$+0.24$ & \correctlabel{B}\,$-0.24$ & \optlabel{C}\,$0.14$ & \optlabel{D}\,$0.62$
\end{tabularx}
\end{qitem}
\begin{solbox}{B}
\textbf{\color{brand}Correct answer.} \textcolor{okgreen}{\textbf{(B)}} $-0.24$

\textbf{\color{brand}Concepts required.} Expectation values factorise for product states.

\textbf{\color{brand}Relevant formulae.} $\langle Z\rangle = P(0) - P(1)$; product state: $\langle Z_0Z_1\rangle = \langle Z_0\rangle\langle Z_1\rangle$; $\ \langle Z_0Z_1\rangle = 1 - 2P(z_0\neq z_1)$.

\textbf{\color{brand}Step-by-step explanation.} For a product state the expectation factorises: $\langle Z_0Z_1\rangle = \langle Z_0\rangle\langle Z_1\rangle$, with
$\langle Z\rangle = P(0) - P(1)$. So $\langle Z_0\rangle = 0.8 - 0.2 = 0.6$ and $\langle Z_1\rangle = 0.3 - 0.7 = -0.4$, giving $-0.24$.
Check via the cut probability: $P(z_0 \neq z_1) = 0.2\cdot0.3 + 0.8\cdot0.7 = 0.62$ and $\langle Z_0Z_1\rangle = 1 - 2\cdot0.62 = -0.24$.

\textbf{\color{brand}Why the other options are incorrect.}
\begin{itemize}[leftmargin=2.4em, itemsep=1pt, topsep=2pt]
\item[\textbf{(A)}] $+0.24$ has the wrong sign: $\langle Z_1\rangle = 0.3 - 0.7 < 0$.
\item[\textbf{(C)}] $0.14 = 0.2\times0.7$ is $P(z_0 = 1, z_1 = 1)$.
\item[\textbf{(D)}] $0.62$ is $P(z_0\neq z_1)$, the probability that the edge is cut, not $\langle Z_0Z_1\rangle$.
\end{itemize}
\textbf{\color{brand}Key takeaway.} Translate between $\langle Z_iZ_j\rangle$ and $P(\text{cut})$ with $\langle Z_iZ_j\rangle = 1 - 2P(\text{cut})$.
\end{solbox}
\qsep

\needspace{13\baselineskip}
\begin{qitem}{31}
\qmeta{Single-correct MCQ}{Conjugation $e^{i\beta X}Ze^{-i\beta X}$}
Conjugating $Z$ by the mixer gate gives $e^{i\beta X}\,Z\,e^{-i\beta X} =$
\par\smallskip\noindent\begin{tabularx}{\linewidth}{@{}*{4}{L}@{}}
\optlabel{A}\,$\cos 2\beta\; Z - \sin 2\beta\; Y$ & \correctlabel{B}\,$\cos 2\beta\; Z + \sin 2\beta\; Y$ & \optlabel{C}\,$\cos\beta\; Z + \sin\beta\; Y$ & \optlabel{D}\,$\cos 2\beta\; Z + \sin 2\beta\; X$
\end{tabularx}
\end{qitem}
\begin{solbox}{B}
\textbf{\color{brand}Correct answer.} \textcolor{okgreen}{\textbf{(B)}} $\cos 2\beta\; Z + \sin 2\beta\; Y$

\textbf{\color{brand}Concepts required.} Heisenberg picture: conjugation of a Pauli by a rotation about another axis.

\textbf{\color{brand}Relevant formulae.} $e^{i\beta X}Z = Ze^{-i\beta X}$; $\ e^{-2i\beta X} = \cos2\beta - i\sin2\beta\,X$; $\ ZX = iY$.

\textbf{\color{brand}Step-by-step explanation.} Because $XZ = -ZX$, moving $Z$ through $e^{i\beta X}$ flips the sign of the exponent: $e^{i\beta X}Z = Z\,e^{-i\beta X}$. Hence
\[ e^{i\beta X}Ze^{-i\beta X} = Z\,e^{-2i\beta X} = Z(\cos 2\beta - i\sin 2\beta\, X) = \cos2\beta\, Z - i\sin2\beta\, ZX . \]
With $ZX = iY$: $-i\sin2\beta\,(iY) = +\sin 2\beta\, Y$. Result $\cos2\beta\,Z + \sin2\beta\,Y$: a rotation of the Bloch vector by $2\beta$
about the $x$ axis (the angle doubles, as for every $e^{-i\phi P}$). This identity is the first step of the $p = 1$ derivation.

\textbf{\color{brand}Why the other options are incorrect.}
\begin{itemize}[leftmargin=2.4em, itemsep=1pt, topsep=2pt]
\item[\textbf{(A)}] $-\sin2\beta\,Y$ comes from using $ZX = -iY$; actually $ZX = +iY$.
\item[\textbf{(C)}] $\cos\beta\,Z + \sin\beta\,Y$ misses the doubling: conjugation by $e^{-i\beta X}$ rotates by $2\beta$.
\item[\textbf{(D)}] $X$ commutes with $e^{i\beta X}$; rotating $Z$ about $x$ produces a $Y$ component, never an $X$ component.
\end{itemize}
\textbf{\color{brand}Key takeaway.} Rotations by $e^{-i\phi P}$ turn the Bloch vector by $2\phi$.
\end{solbox}
\qsep

\needspace{13\baselineskip}
\begin{qitem}{32}
\qmeta{Single-correct MCQ}{State-vector memory on one GPU}
A state vector needs $16 \times 2^n$ bytes in fp64 and $8\times 2^n$ bytes in fp32. Ignoring simulator workspace, the
largest $n$ whose state fits on a single 80 GB ($80\times10^9$ bytes) GPU is
\par\smallskip\noindent\begin{tabularx}{\linewidth}{@{}*{2}{L}@{}}
\optlabel{A}\,$33$ (fp64) and $34$ (fp32) & \optlabel{B}\,$36$ (fp64) and $37$ (fp32) \\[3pt]
\correctlabel{C}\,$32$ (fp64) and $33$ (fp32) & \optlabel{D}\,$32$ (fp64) and $64$ (fp32)
\end{tabularx}
\end{qitem}
\begin{solbox}{C}
\textbf{\color{brand}Correct answer.} \textcolor{okgreen}{\textbf{(C)}} $32$ (fp64) and $33$ (fp32)

\textbf{\color{brand}Concepts required.} State-vector memory grows as $2^n$; precision sets bytes per amplitude.

\textbf{\color{brand}Relevant formulae.} fp64: $16\times2^n$ bytes (complex128); fp32: $8\times2^n$ bytes (complex64).

\textbf{\color{brand}Step-by-step explanation.} fp64: $16\times2^{32} = 6.87\times10^{10}$ bytes $= 68.7$ GB $\le 80$ GB, but $16\times2^{33} = 137.4$ GB $> 80$ GB, so $n = 32$.
fp32 halves the memory: $8\times2^{33} = 68.7$ GB fits, $8\times2^{34} = 137.4$ GB does not, so $n = 33$.
Halving the precision buys exactly one qubit. The notebook's 15-qubit state needs $16\times2^{15} = 524\,288$ bytes (``524.29 KB'').

\textbf{\color{brand}Why the other options are incorrect.}
\begin{itemize}[leftmargin=2.4em, itemsep=1pt, topsep=2pt]
\item[\textbf{(A)}] $n = 33$ in fp64 needs $137$ GB, more than 80 GB.
\item[\textbf{(B)}] $36$ qubits need $1.1$ TB in fp64.
\item[\textbf{(D)}] Halving the memory per amplitude adds one qubit, not 32.
\end{itemize}
\textbf{\color{brand}Key takeaway.} Every added qubit doubles the memory; fp32 buys exactly one qubit.
\end{solbox}
\qsep

\needspace{13\baselineskip}
\begin{qitem}{33}
\qmeta{Single-correct MCQ}{Bit-flip probability of $R_X$}
The gate $R_X(2\beta)$ with $\beta = \pi/8$ is applied to $\ket0$. The probability of measuring 1 is
\par\smallskip\noindent\begin{tabularx}{\linewidth}{@{}*{2}{L}@{}}
\optlabel{A}\,$\sin(\pi/8) \approx 0.383$ & \optlabel{B}\,$\tfrac12$ \\[3pt]
\correctlabel{C}\,$\tfrac{2 - \sqrt2}{4} \approx 0.146$ & \optlabel{D}\,$\tfrac{2 + \sqrt2}{4} \approx 0.854$
\end{tabularx}
\end{qitem}
\begin{solbox}{C}
\textbf{\color{brand}Correct answer.} \textcolor{okgreen}{\textbf{(C)}} $\tfrac{2 - \sqrt2}{4} \approx 0.146$

\textbf{\color{brand}Concepts required.} Action of $R_X$ on $\ket0$; Born rule; half-angle identity.

\textbf{\color{brand}Relevant formulae.} $R_X(2\beta)\ket0 = \cos\beta\ket0 - i\sin\beta\ket1$; $\ \sin^2\phi = \tfrac{1 - \cos2\phi}{2}$.

\textbf{\color{brand}Step-by-step explanation.} $R_X(2\beta)\ket0 = \cos\beta\ket0 - i\sin\beta\ket1$ (first column of the matrix in Q.23), so $P(1) = \sin^2\beta$.
For $\beta = \pi/8$: $\sin^2(\pi/8) = \tfrac{1 - \cos(\pi/4)}{2} = \tfrac{1 - \sqrt2/2}{2} = \tfrac{2 - \sqrt2}{4} \approx 0.146$.
$\sin(\pi/8)$ is the amplitude's modulus, not the probability.

\textbf{\color{brand}Why the other options are incorrect.}
\begin{itemize}[leftmargin=2.4em, itemsep=1pt, topsep=2pt]
\item[\textbf{(A)}] $0.383 = \sin(\pi/8)$ is the modulus of the amplitude, not its square.
\item[\textbf{(B)}] $\frac12$ would need $\beta = \pi/4$.
\item[\textbf{(D)}] $0.854 = \cos^2(\pi/8)$ is $P(0)$.
\end{itemize}
\textbf{\color{brand}Key takeaway.} $R_X(2\beta)$ flips a bit with probability $\sin^2\beta$.
\end{solbox}
\qsep

\needspace{13\baselineskip}
\begin{qitem}{34}
\qmeta{Single-correct MCQ}{$\ket+^{\otimes n}$ as top eigenstate of $H_M$}
Which statement about $\ket{+}^{\otimes n}$ and $H_M = \sum_{i} X_i$ is correct?
\begin{itemize}[leftmargin=2.6em, labelsep=0.4em, itemsep=2pt, topsep=3pt]
\item[\optlabel{A}] It is an eigenstate of $H_M$ with eigenvalue $0$
\item[\correctlabel{B}] It is the eigenstate of $H_M$ with the \emph{largest} eigenvalue, $+n$
\item[\optlabel{C}] It is not an eigenstate of $H_M$
\item[\optlabel{D}] It is the ground state of $H_M$, with eigenvalue $-n$
\end{itemize}
\end{qitem}
\begin{solbox}{B}
\textbf{\color{brand}Correct answer.} \textcolor{okgreen}{\textbf{(B)}} It is the eigenstate of $H_M$ with the \emph{largest} eigenvalue, $+n$

\textbf{\color{brand}Concepts required.} Eigenstates of $X$ and of sums of commuting single-qubit operators.

\textbf{\color{brand}Relevant formulae.} $X\ket\pm = \pm\ket\pm$; $\ H_M\ket+^{\otimes n} = n\ket+^{\otimes n}$; spectrum of $H_M$: $\{-n, -n+2, \dots, n\}$.

\textbf{\color{brand}Step-by-step explanation.} $X\ket+ = +\ket+$, so each $X_i$ gives $+1$ on $\ket+^{\otimes n}$ and $H_M\ket+^{\otimes n} = n\,\ket+^{\otimes n}$.
The eigenvalues of $H_M$ range from $-n$ (state $\ket-^{\otimes n}$) to $+n$, so $\ket+^{\otimes n}$ is the \emph{highest} eigenstate of
$H_M$ (equivalently the ground state of $-H_M$). This matters for the annealing interpretation: adiabatic evolution keeps a
state in the \emph{same} position of the spectrum, so with $+H_M$ the matching target is the highest state of the final
Hamiltonian (see Q.54).

\textbf{\color{brand}Why the other options are incorrect.}
\begin{itemize}[leftmargin=2.4em, itemsep=1pt, topsep=2pt]
\item[\textbf{(A)}] Eigenvalue $0$ (for even $n$) belongs to states with equally many $\ket+$ and $\ket-$ factors.
\item[\textbf{(C)}] $\ket+^{\otimes n}$ is an eigenstate: each $X_i$ leaves it unchanged.
\item[\textbf{(D)}] The ground state of $\sum_iX_i$ is $\ket-^{\otimes n}$ with eigenvalue $-n$; $\ket+^{\otimes n}$ is the ground state of $-H_M$.
\end{itemize}
\textbf{\color{brand}Key takeaway.} The start state is the \emph{top} state of $+H_M$ --- remember this when reading QAOA as annealing.
\end{solbox}
\qsep

\needspace{13\baselineskip}
\begin{qitem}{35}
\qmeta{Assertion--Reason}{Complex-conjugation symmetry}
\textbf{Assertion (A):} A state $\ket\psi$ and its complex conjugate $\ket{\psi^*}$ (every amplitude conjugated) give
identical computational-basis measurement statistics.\\[1mm]
\textbf{Reason (R):} The initial QAOA state $\ket{+}^{\otimes n}$ has only real amplitudes.
\begin{itemize}[leftmargin=2.6em, labelsep=0.4em, itemsep=2pt, topsep=3pt]
\item[\optlabel{A}] Both Assertion (A) and Reason (R) are true, and R is the correct explanation of A
\item[\correctlabel{B}] Both A and R are true, but R is \emph{not} the correct explanation of A
\item[\optlabel{C}] A is true, but R is false
\item[\optlabel{D}] A is false, but R is true
\end{itemize}
\end{qitem}
\begin{solbox}{B}
\textbf{\color{brand}Correct answer.} \textcolor{okgreen}{\textbf{(B)}} Both A and R are true, but R is \emph{not} the correct explanation of A

\textbf{\color{brand}Concepts required.} Measurement probabilities depend only on moduli of amplitudes; complex conjugation symmetry of QAOA.

\textbf{\color{brand}Relevant formulae.} $|a^*|^2 = |a|^2$; real $H_C$, $H_M$, real start state $\Rightarrow$ $\ket{\boldsymbol\gamma,\boldsymbol\beta}^* = \ket{-\boldsymbol\gamma,-\boldsymbol\beta}$.

\textbf{\color{brand}Step-by-step explanation.} \textbf{A is true:} the outcome probabilities are $|a_z|^2$ and $|a_z^*|^2 = |a_z|^2$ for every $z$, for \emph{any} state.
\textbf{R is true:} all amplitudes of $\ket+^{\otimes n}$ equal $2^{-n/2}$.
But R does not explain A: A holds for every state, real or not; the explanation is $|a^*| = |a|$.
(R \emph{is} the reason for a different fact: since $H_C$, $H_M$ and $\ket+^{\otimes n}$ are real,
$\ket{\boldsymbol\gamma,\boldsymbol\beta}^* = \ket{-\boldsymbol\gamma, -\boldsymbol\beta}$, so $F_p(-\boldsymbol\gamma, -\boldsymbol\beta) = F_p(\boldsymbol\gamma, \boldsymbol\beta)$.) Answer: (B).

\textbf{\color{brand}Why the other options are incorrect.}
\begin{itemize}[leftmargin=2.4em, itemsep=1pt, topsep=2pt]
\item[\textbf{(A)}] needs A true, R true \emph{and} R explaining A; R does not explain A, which holds for every state.
\item[\textbf{(C)}] needs A true and R false; R is true.
\item[\textbf{(D)}] needs A false and R true; A is true.
\end{itemize}
\textbf{\color{brand}Key takeaway.} $F_p(-\boldsymbol\gamma, -\boldsymbol\beta) = F_p(\boldsymbol\gamma, \boldsymbol\beta)$: the landscape is point-symmetric.
\end{solbox}
\qsep

\needspace{13\baselineskip}
\begin{qitem}{36}
\qmeta{Single-correct MCQ}{Linearity of $\langle H_C\rangle$}
For the triangle $K_3$ (unit weights) a state has $\langle Z_0Z_1\rangle = -0.6$, $\langle Z_1Z_2\rangle = -0.2$ and
$\langle Z_0Z_2\rangle = 0.4$. Its expected cut $\langle C\rangle$ is
\par\smallskip\noindent\begin{tabularx}{\linewidth}{@{}*{4}{L}@{}}
\optlabel{A}\,$-0.2$ & \optlabel{B}\,$1.9$ & \optlabel{C}\,$1.3$ & \correctlabel{D}\,$1.7$
\end{tabularx}
\end{qitem}
\begin{solbox}{D}
\textbf{\color{brand}Correct answer.} \textcolor{okgreen}{\textbf{(D)}} $1.7$

\textbf{\color{brand}Concepts required.} Linearity of expectation values; $\langle C\rangle$ from Pauli expectations.

\textbf{\color{brand}Relevant formulae.} $\langle H_C\rangle = \sum\tfrac{w_{ij}}{2}\langle Z_iZ_j\rangle$; $\ \langle C\rangle = \tfrac W2 - \langle H_C\rangle$.

\textbf{\color{brand}Step-by-step explanation.} By linearity $\langle H_C\rangle = \tfrac12(-0.6 - 0.2 + 0.4) = -0.2$, and $\langle C\rangle = \tfrac W2 - \langle H_C\rangle = 1.5 + 0.2 = 1.7$.
Edge by edge: $\langle C_{ij}\rangle = \tfrac12(1 - \langle Z_iZ_j\rangle)$ gives $0.8 + 0.6 + 0.3 = 1.7$.
This is how \texttt{cudaq.observe} works: it evaluates $\langle H\rangle = \sum_k c_k\langle P_k\rangle$ term by term. The value is below
$C_{\max} = 2$, as it must be.

\textbf{\color{brand}Why the other options are incorrect.}
\begin{itemize}[leftmargin=2.4em, itemsep=1pt, topsep=2pt]
\item[\textbf{(A)}] $-0.2$ is $\langle H_C\rangle$, not $\langle C\rangle$.
\item[\textbf{(B)}] $1.9 = 1.5 + 0.4$ forgets the factor $\frac12$ in front of the $Z_iZ_j$ terms.
\item[\textbf{(C)}] $1.3 = 1.5 - 0.2$ uses $+\langle H_C\rangle$ instead of $-\langle H_C\rangle$.
\end{itemize}
\textbf{\color{brand}Key takeaway.} Per edge: $\langle C_{ij}\rangle = \frac12(1 - \langle Z_iZ_j\rangle)$, then add.
\end{solbox}
\qsep

\needspace{13\baselineskip}
\begin{qitem}{37}
\qmeta{Single-correct MCQ}{Mixer at $\beta = \pi/2$; period $\pi/2$}
On $n$ qubits, $e^{-i(\pi/2) H_M}$ with $H_M = \sum_i X_i$ equals ..., and therefore (since $C(z) = C(\bar z)$) $F_p$ is
periodic in every single $\beta_l$ with period ...
\par\smallskip\noindent\begin{tabularx}{\linewidth}{@{}*{2}{L}@{}}
\correctlabel{A}\,$(-i)^n X^{\otimes n}$; $\ \pi/2$ & \optlabel{B}\,$(-i)^n X^{\otimes n}$; $\ \pi$ \\[3pt]
\optlabel{C}\,$I$; $\ \pi/2$ & \optlabel{D}\,$X^{\otimes n}$; $\ 2\pi$
\end{tabularx}
\end{qitem}
\begin{solbox}{A}
\textbf{\color{brand}Correct answer.} \textcolor{okgreen}{\textbf{(A)}} $(-i)^n X^{\otimes n}$; $\ \pi/2$

\textbf{\color{brand}Concepts required.} Mixer at special angles; global bit-flip symmetry of the cut.

\textbf{\color{brand}Relevant formulae.} $e^{-i\frac\pi2X} = -iX$; $\ X^{\otimes n}Z_iZ_jX^{\otimes n} = Z_iZ_j$; $\ C(\bar z) = C(z)$.

\textbf{\color{brand}Step-by-step explanation.} $e^{-i(\pi/2)X} = \cos\tfrac\pi2 I - i\sin\tfrac\pi2 X = -iX$, and the mixer factorises over qubits, so
$e^{-i(\pi/2)H_M} = (-iX)^{\otimes n} = (-i)^n X^{\otimes n}$: a global phase times ``flip every bit''.
Shifting $\beta_l \to \beta_l + \pi/2$ inserts $X^{\otimes n}$ after layer $l$. $X^{\otimes n}$ commutes with every $Z_iZ_j$ (two sign
flips) and with $H_M$, so it can be moved to the end, where it maps $\ket z \to \ket{\bar z}$; since $C(\bar z) = C(z)$ the
expectation is unchanged. Hence period $\pi/2$ in each $\beta_l$ --- one quarter of the naive $2\pi$.

\textbf{\color{brand}Why the other options are incorrect.}
\begin{itemize}[leftmargin=2.4em, itemsep=1pt, topsep=2pt]
\item[\textbf{(B)}] The operator is right, but the period is $\pi/2$: the flip it produces is already a symmetry.
\item[\textbf{(C)}] $e^{-i(\pi/2)H_M}$ is not the identity; it flips every bit.
\item[\textbf{(D)}] The global phase $(-i)^n$ is missing (harmless) and the period is wrong: $2\pi$ is only the naive period.
\end{itemize}
\textbf{\color{brand}Key takeaway.} Search $\beta_l$ in an interval of length $\pi/2$ only.
\end{solbox}
\qsep

\needspace{13\baselineskip}
\begin{qitem}{38}
\qmeta{Statement I / II}{Measuring $Z_iZ_j$ from shots}
\textbf{Statement I:} $Z_iZ_j$ has eigenvalue $-1$ exactly on the computational basis states with $z_i \neq z_j$.\\[1mm]
\textbf{Statement II:} $\langle Z_iZ_j\rangle$ can be estimated from the same computational-basis shots that are used to sample
candidate splits.
\begin{itemize}[leftmargin=2.6em, labelsep=0.4em, itemsep=2pt, topsep=3pt]
\item[\correctlabel{A}] Both Statement I and Statement II are true
\item[\optlabel{B}] Both Statement I and Statement II are false
\item[\optlabel{C}] Statement I is true, but Statement II is false
\item[\optlabel{D}] Statement I is false, but Statement II is true
\end{itemize}
\end{qitem}
\begin{solbox}{A}
\textbf{\color{brand}Correct answer.} \textcolor{okgreen}{\textbf{(A)}} Both Statement I and Statement II are true

\textbf{\color{brand}Concepts required.} Diagonal observables are measured in the computational basis.

\textbf{\color{brand}Relevant formulae.} $Z_iZ_j\ket z = (-1)^{z_i\oplus z_j}\ket z$; $\ \widehat{\langle Z_iZ_j\rangle} = \tfrac1S\sum_{\text{shots}}(-1)^{z_i\oplus z_j}$.

\textbf{\color{brand}Step-by-step explanation.} \textbf{I is true:} $Z_iZ_j\ket z = s_is_j\ket z$ with $s = (-1)^z$, and $s_is_j = -1 \iff z_i \neq z_j$.
\textbf{II is true:} $Z_iZ_j$ is diagonal in the computational basis, so each shot $z$ directly yields the value
$(-1)^{z_i\oplus z_j}$; averaging over shots estimates $\langle Z_iZ_j\rangle$ (and $\langle H_C\rangle$) with no extra basis change.
This is what happens on hardware targets, where the notebook estimates the objective from shots. Answer: (A).

\textbf{\color{brand}Why the other options are incorrect.}
\begin{itemize}[leftmargin=2.4em, itemsep=1pt, topsep=2pt]
\item[\textbf{(B)}] needs both statements false; both are true.
\item[\textbf{(C)}] needs I true and II false; II is true.
\item[\textbf{(D)}] needs I false and II true; I is true.
\end{itemize}
\textbf{\color{brand}Key takeaway.} One set of shots gives both the candidate splits and an estimate of the objective.
\end{solbox}
\qsep

\needspace{20\baselineskip}
\begin{passagebox}{Paragraph for Questions 39 and 40}
A single edge $(0,1)$ with weight 1 is treated with one QAOA layer in the notebook's convention. Starting from
$\ket{00}$ the gates are: $H$ on both qubits; CNOT (control 0, target 1); $R_Z(\theta)$ on qubit 1 with $\theta = \gamma w = \gamma$;
CNOT (control 0, target 1); $R_X(2\beta)$ on both qubits. The angles are $\gamma = \pi/2$ and $\beta = \pi/8$.
Amplitudes are listed in the order $\ket{00}, \ket{01}, \ket{10}, \ket{11}$ (qubit 0 written first).
\end{passagebox}
\needspace{13\baselineskip}
\begin{qitem}{39}
\qmeta{Paragraph based}{Trace: phase after the cost gate}
Immediately after the second CNOT (before the mixer), the phase of the amplitude of $\ket{01}$ relative to the amplitude
of $\ket{00}$ is
\par\smallskip\noindent\begin{tabularx}{\linewidth}{@{}*{4}{L}@{}}
\optlabel{A}\,$45^\circ$ & \optlabel{B}\,$180^\circ$ & \correctlabel{C}\,$90^\circ$ & \optlabel{D}\,$0^\circ$
\end{tabularx}
\end{qitem}
\begin{solbox}{C}
\textbf{\color{brand}Correct answer.} \textcolor{okgreen}{\textbf{(C)}} $90^\circ$

\textbf{\color{brand}Concepts required.} Cost gate as a diagonal phase: equal bits and different bits get opposite phases.

\textbf{\color{brand}Relevant formulae.} $e^{-i\gamma Z_0Z_1/2}\ket{z_0z_1} = e^{-i\frac\gamma2(-1)^{z_0\oplus z_1}}\ket{z_0z_1}$.

\textbf{\color{brand}Step-by-step explanation.} After the Hadamards every amplitude is $\tfrac12$. CNOT$\cdot(I\otimes R_Z(\gamma))\cdot$CNOT $= e^{-i\gamma Z_0Z_1/2}$: the first CNOT
writes $z_0\oplus z_1$ on qubit 1, $R_Z(\gamma)$ gives phase $e^{-i\gamma/2}$ if that bit is 0 and $e^{+i\gamma/2}$ if it is 1, and the
second CNOT restores the bits. With $\gamma = \pi/2$:
$\ket{00}, \ket{11} \to \tfrac12 e^{-i\pi/4}$ and $\ket{01}, \ket{10} \to \tfrac12 e^{+i\pi/4}$.
Relative phase of $\ket{01}$ to $\ket{00}$: $e^{i\pi/4}/e^{-i\pi/4} = e^{i\pi/2}$, i.e.\ $90^\circ$ ($= \gamma$). The measurement
probabilities are still all $\tfrac14$.

\textbf{\color{brand}Why the other options are incorrect.}
\begin{itemize}[leftmargin=2.4em, itemsep=1pt, topsep=2pt]
\item[\textbf{(A)}] $45^\circ$ is the phase of each amplitude, $\pm\gamma/2$; the \emph{relative} phase is $\gamma$.
\item[\textbf{(B)}] $180^\circ$ would need $\gamma = \pi$.
\item[\textbf{(D)}] $0^\circ$ would mean the cost layer does nothing; $\ket{00}$ and $\ket{01}$ get opposite phases.
\end{itemize}
\textbf{\color{brand}Key takeaway.} The cost layer marks cut and uncut strings with a relative phase $\gamma w$.
\end{solbox}
\qsep

\needspace{13\baselineskip}
\begin{qitem}{40}
\qmeta{Paragraph based}{Trace: cut probability after the mixer}
After the mixer, the probability of measuring a \emph{cut} bitstring ($01$ or $10$) is
\par\smallskip\noindent\begin{tabularx}{\linewidth}{@{}*{4}{L}@{}}
\optlabel{A}\,$1$ & \optlabel{B}\,$\tfrac34$ & \correctlabel{C}\,$0$ & \optlabel{D}\,$\tfrac12$
\end{tabularx}
\end{qitem}
\begin{solbox}{C}
\textbf{\color{brand}Correct answer.} \textcolor{okgreen}{\textbf{(C)}} $0$

\textbf{\color{brand}Concepts required.} $p = 1$ single-edge formula; the sign relation between $\beta$ and $\gamma$.

\textbf{\color{brand}Relevant formulae.} $\langle Z_0Z_1\rangle = \sin4\beta\sin\gamma$, $\ P(\text{cut}) = \tfrac12(1 - \sin4\beta\sin\gamma)$.

\textbf{\color{brand}Step-by-step explanation.} For one edge with one layer in this convention $\langle Z_0Z_1\rangle = \sin4\beta\,\sin\gamma$ (derived in the Answer Book of Q.47;
the derivation uses Q.31 twice and the conjugation of $Y_0Z_1$ by the cost gate). At $\gamma = \pi/2$, $\beta = \pi/8$:
$\sin(\pi/2)\sin(\pi/2) = 1$, so $\langle Z_0Z_1\rangle = +1$ and $P(\text{cut}) = \tfrac12(1 - \langle Z_0Z_1\rangle) = 0$.
Explicitly, the final state has support only on $\ket{00}$ and $\ket{11}$ (state-vector simulation gives the probabilities
$0.5, 0, 0, 0.5$ for $\ket{00}, \ket{01}, \ket{10}, \ket{11}$). With $\beta = -\pi/8$ instead, the probability of a cut bitstring is $1$. The sign of $\beta$
relative to $\gamma$ decides whether the layer helps or hurts.

\textbf{\color{brand}Why the other options are incorrect.}
\begin{itemize}[leftmargin=2.4em, itemsep=1pt, topsep=2pt]
\item[\textbf{(A)}] $1$ is the result for $\beta = -\pi/8$, not $+\pi/8$.
\item[\textbf{(B)}] $\frac34$ is the ring's best value per edge at $p = 1$, unrelated here.
\item[\textbf{(D)}] $\frac12$ would mean the layer did nothing (as at $\gamma = 0$ or $\beta = 0$).
\end{itemize}
\textbf{\color{brand}Key takeaway.} Same $|\beta|$, opposite sign: the layer can solve the edge perfectly or make it perfectly wrong.
\end{solbox}
\qsep

\newpage

\begin{sectionbanner}SECTION C \enspace$\cdot$\enspace QAOA Basics \hfill Q.41--60\end{sectionbanner}

\needspace{13\baselineskip}
\begin{qitem}{41}
\qmeta{Single-correct MCQ}{\texttt{cost\_hamiltonian()} output}
For \texttt{edges = [(0, 1, 2.0), (1, 2, 1.0)]} the notebook's function \texttt{cost\_hamiltonian()} returns the operator
\par\smallskip\noindent\begin{tabularx}{\linewidth}{@{}*{2}{L}@{}}
\correctlabel{A}\,$1.0\,Z_0Z_1 + 0.5\,Z_1Z_2$ & \optlabel{B}\,$-1.0\,Z_0Z_1 - 0.5\,Z_1Z_2$ \\[3pt]
\optlabel{C}\,$2.0\,Z_0Z_1 + 1.0\,Z_1Z_2$ & \optlabel{D}\,$1.5\,I - 1.0\,Z_0Z_1 - 0.5\,Z_1Z_2$
\end{tabularx}
\end{qitem}
\begin{solbox}{A}
\textbf{\color{brand}Correct answer.} \textcolor{okgreen}{\textbf{(A)}} $1.0\,Z_0Z_1 + 0.5\,Z_1Z_2$

\textbf{\color{brand}Concepts required.} The notebook's cost Hamiltonian has coefficient $w/2$ per edge and no constant; the cut operator is $C = \frac W2 - H_C$.

\textbf{\color{brand}Relevant formulae.} $H_C = \sum_{(i,j)\in E}\tfrac{w_{ij}}{2}Z_iZ_j$ (code: \texttt{0.5 * w * spin.z(i) * spin.z(j)}); $\ W = \sum w_{ij}$.

\textbf{\color{brand}Step-by-step explanation.} The code builds \texttt{term = 0.5 * w * cudaq.spin.z(i) * cudaq.spin.z(j)} for every edge and adds the terms:
$H_C = \sum \tfrac{w_{ij}}{2}Z_iZ_j = \tfrac{2.0}{2}Z_0Z_1 + \tfrac{1.0}{2}Z_1Z_2 = 1.0\,Z_0Z_1 + 0.5\,Z_1Z_2$.
No constant term is included; the cut is recovered as $C = \tfrac W2 - H_C$ with $W = 3$, i.e.\ $C = 1.5\,I - H_C$ --- the third
option is the cut operator $C$, not what the function returns.

\textbf{\color{brand}Why the other options are incorrect.}
\begin{itemize}[leftmargin=2.4em, itemsep=1pt, topsep=2pt]
\item[\textbf{(B)}] The negative sign belongs to $C$, not to $H_C$; with it, minimising would find the \emph{minimum} cut.
\item[\textbf{(C)}] $2.0\,Z_0Z_1 + 1.0\,Z_1Z_2$ omits the factor $\frac12$ that comes from $[z_i\neq z_j] = \frac12(1 - s_is_j)$.
\item[\textbf{(D)}] $1.5\,I - H_C$ is the cut operator $C$; the function returns $H_C$ only.
\end{itemize}
\textbf{\color{brand}Key takeaway.} Read the code: \texttt{cost\_hamiltonian()} $= \sum \frac w2 Z_iZ_j$, and $C = W/2 - H_C$.
\end{solbox}
\qsep

\needspace{13\baselineskip}
\begin{qitem}{42}
\qmeta{Single-correct MCQ}{All $\gamma_l = 0$}
In the notebook's convention, if all cost angles are zero ($\gamma_1 = \dots = \gamma_p = 0$), the expected cut
$F_p(\boldsymbol\gamma, \boldsymbol\beta)$ equals
\par\smallskip\noindent\begin{tabularx}{\linewidth}{@{}*{2}{L}@{}}
\optlabel{A}\,$W/2 + \tfrac12\sum_l \sin 4\beta_l$ & \optlabel{B}\,$C_{\max}$ \\[3pt]
\correctlabel{C}\,$W/2$ for every choice of $\boldsymbol\beta$ & \optlabel{D}\,$0$
\end{tabularx}
\end{qitem}
\begin{solbox}{C}
\textbf{\color{brand}Correct answer.} \textcolor{okgreen}{\textbf{(C)}} $W/2$ for every choice of $\boldsymbol\beta$

\textbf{\color{brand}Concepts required.} Without cost phases, the mixer acts on its own eigenstate.

\textbf{\color{brand}Relevant formulae.} $e^{-i\beta H_M}\ket+^{\otimes n} = e^{-i\beta n}\ket+^{\otimes n}$; $\ \langle C\rangle_{\ket+^{\otimes n}} = W/2$.

\textbf{\color{brand}Step-by-step explanation.} With $\gamma_l = 0$ every cost layer is the identity, so the state is $\prod_l e^{-i\beta_l H_M}\ket+^{\otimes n}$.
$\ket+^{\otimes n}$ is an eigenstate of $H_M$ with eigenvalue $n$ (Q.34), so each mixer only multiplies it by the phase $e^{-i\beta_l n}$.
The state remains $\ket+^{\otimes n}$ (up to a global phase), all $2^n$ splits are equally likely, and $\langle C\rangle = W/2$
(the random baseline; $15$ for the notebook's graph). The cost layer creates the relative phases that the mixer turns into
probability changes; without it nothing happens.

\textbf{\color{brand}Why the other options are incorrect.}
\begin{itemize}[leftmargin=2.4em, itemsep=1pt, topsep=2pt]
\item[\textbf{(A)}] There is no $\beta$ dependence at all when every $\gamma_l = 0$; the extra term is invented.
\item[\textbf{(B)}] $C_{\max}$ would require the circuit to find the optimum without any information about the graph.
\item[\textbf{(D)}] $0$ is the value of the strings $00\dots0$ and $11\dots1$, not of the uniform superposition.
\end{itemize}
\textbf{\color{brand}Key takeaway.} No cost phases, no progress: the mixer alone leaves the random baseline $W/2$.
\end{solbox}
\qsep

\needspace{13\baselineskip}
\begin{qitem}{43}
\qmeta{Single-correct MCQ}{CNOT--$R_Z$--CNOT decomposition}
In the notebook's kernel each edge $(i, j)$ with weight $w$ is implemented as
\texttt{x.ctrl(q[i], q[j]); rz(gamma * w, q[j]); x.ctrl(q[i], q[j])}. With $R_Z(\theta) = e^{-i\theta Z/2}$, this sequence equals
\par\smallskip\noindent\begin{tabularx}{\linewidth}{@{}*{4}{L}@{}}
\optlabel{A}\,$e^{+i\gamma\,(w/2)\,Z_iZ_j}$ & \optlabel{B}\,$e^{-i\gamma\, w\, Z_iZ_j}$ & \optlabel{C}\,$e^{-i\gamma\,(w/2)\,Z_j}$ & \correctlabel{D}\,$e^{-i\gamma\,(w/2)\,Z_iZ_j}$
\end{tabularx}
\end{qitem}
\begin{solbox}{D}
\textbf{\color{brand}Correct answer.} \textcolor{okgreen}{\textbf{(D)}} $e^{-i\gamma\,(w/2)\,Z_iZ_j}$

\textbf{\color{brand}Concepts required.} CNOT--$R_Z$--CNOT computes the parity $z_i\oplus z_j$ into a phase.

\textbf{\color{brand}Relevant formulae.} $R_Z(\theta)\ket b = e^{-i\theta(-1)^b/2}\ket b$; $\ \mathrm{CNOT}(I\otimes R_Z(\theta))\mathrm{CNOT} = e^{-i\theta Z\otimes Z/2}$.

\textbf{\color{brand}Step-by-step explanation.} Act on a basis state $\ket{z_i z_j}$. The first CNOT gives $\ket{z_i,\, z_i\oplus z_j}$. $R_Z(\theta)$ on qubit $j$ multiplies by
$e^{-i\theta/2}$ if $z_i \oplus z_j = 0$ and by $e^{+i\theta/2}$ if $z_i \oplus z_j = 1$, i.e.\ by $e^{-i(\theta/2)\, s_is_j}$ with $s = (-1)^z$.
The second CNOT restores $\ket{z_iz_j}$. So the sequence is the diagonal operator $e^{-i(\theta/2)Z_iZ_j}$, and with
$\theta = \gamma w$: $e^{-i\gamma(w/2)Z_iZ_j}$ --- exactly the factor of $e^{-i\gamma H_C}$ for that edge. (Equivalently: CNOT
conjugation maps $I\otimes Z$ to $Z\otimes Z$, see Q.27.)

\textbf{\color{brand}Why the other options are incorrect.}
\begin{itemize}[leftmargin=2.4em, itemsep=1pt, topsep=2pt]
\item[\textbf{(A)}] $R_Z(\theta)$ with $\theta = \gamma w > 0$ gives $e^{-i\dots}$; the sign flips only for negative angles.
\item[\textbf{(B)}] $e^{-i\gamma wZ_iZ_j}$ forgets the half-angle convention of $R_Z$; it would correspond to \texttt{rz(2 * gamma * w)}.
\item[\textbf{(C)}] $e^{-i\gamma(w/2)Z_j}$ ignores the CNOTs, which turn the single-qubit $Z_j$ into the parity $Z_iZ_j$.
\end{itemize}
\textbf{\color{brand}Key takeaway.} Two CNOTs and one $R_Z(\gamma w)$ per edge implement exactly $e^{-i\gamma\frac w2Z_iZ_j}$.
\end{solbox}
\qsep

\needspace{13\baselineskip}
\begin{qitem}{44}
\qmeta{Single-correct MCQ}{Why \texttt{rx(2*beta)}}
Why does the notebook apply \texttt{rx(2.0 * beta, q)} (and not \texttt{rx(beta, q)}) in the mixer layer?
\begin{itemize}[leftmargin=2.6em, labelsep=0.4em, itemsep=2pt, topsep=3pt]
\item[\optlabel{A}] Because the mixer must be applied twice per layer
\item[\optlabel{B}] Because the mixer Hamiltonian is $H_M = 2\sum_i X_i$
\item[\optlabel{C}] Because $\beta$ is measured in units of $\pi/2$
\item[\correctlabel{D}] Because $R_X(\theta) = e^{-i\theta X/2}$, so $R_X(2\beta) = e^{-i\beta X}$, the factor of $e^{-i\beta H_M}$ for one qubit
\end{itemize}
\end{qitem}
\begin{solbox}{D}
\textbf{\color{brand}Correct answer.} \textcolor{okgreen}{\textbf{(D)}} Because $R_X(\theta) = e^{-i\theta X/2}$, so $R_X(2\beta) = e^{-i\beta X}$, the factor of $e^{-i\beta H_M}$ for one qubit

\textbf{\color{brand}Concepts required.} Half-angle convention of rotation gates.

\textbf{\color{brand}Relevant formulae.} $R_X(\theta) = e^{-i\theta X/2}$ $\Rightarrow$ $R_X(2\beta) = e^{-i\beta X}$; $\ e^{-i\beta H_M} = \prod_i e^{-i\beta X_i}$.

\textbf{\color{brand}Step-by-step explanation.} The target is $U_M(\beta) = e^{-i\beta H_M}$ with $H_M = \sum_i X_i$. Since the $X_i$ commute,
$e^{-i\beta H_M} = \prod_i e^{-i\beta X_i}$. CUDA-Q's \texttt{rx}, like the standard definition, is $R_X(\theta) = e^{-i\theta X/2}$ (half-angle
convention). Matching $e^{-i\theta X/2} = e^{-i\beta X}$ gives $\theta = 2\beta$. The same half-angle convention is why the cost gate
uses \texttt{rz(gamma * w)} for the coefficient $w/2$ (previous question).

\textbf{\color{brand}Why the other options are incorrect.}
\begin{itemize}[leftmargin=2.4em, itemsep=1pt, topsep=2pt]
\item[\textbf{(A)}] Each qubit receives exactly one \texttt{rx} per layer.
\item[\textbf{(B)}] $H_M = \sum_iX_i$ (code docstring and Section 1); the 2 comes from the gate convention, not from $H_M$.
\item[\textbf{(C)}] $\beta$ is in radians; the period $\pi/2$ (Q.37) is a property of $F_p$, not a unit.
\end{itemize}
\textbf{\color{brand}Key takeaway.} Factor 2 in \texttt{rx(2*beta)}, factor $w$ (not $w/2$) in \texttt{rz(gamma*w)}: both from $e^{-i\theta P/2}$.
\end{solbox}
\qsep

\needspace{13\baselineskip}
\begin{qitem}{45}
\qmeta{Single-correct MCQ}{Gate and CNOT counts}
Using the notebook's gate structure ($H$ on every qubit, then per layer 3 gates per edge and one $R_X$ per qubit), a
graph with $n = 20$ nodes and $|E| = 57$ edges at depth $p = 3$ needs how many gates (excluding measurements), and how many of
them are CNOTs?
\par\smallskip\noindent\begin{tabularx}{\linewidth}{@{}*{4}{L}@{}}
\optlabel{A}\,$573$ and $342$ & \correctlabel{B}\,$593$ and $342$ & \optlabel{C}\,$593$ and $171$ & \optlabel{D}\,$402$ and $228$
\end{tabularx}
\end{qitem}
\begin{solbox}{B}
\textbf{\color{brand}Correct answer.} \textcolor{okgreen}{\textbf{(B)}} $593$ and $342$

\textbf{\color{brand}Concepts required.} Gate count of the notebook's QAOA circuit.

\textbf{\color{brand}Relevant formulae.} gates $= n + p\,(3|E| + n)$; CNOTs $= 2p|E|$; $R_Z = p|E|$; $R_X = pn$.

\textbf{\color{brand}Step-by-step explanation.} The notebook's estimate is \texttt{gates = N + p * (3 * len(edges) + N)}.
Total: $n + p(3|E| + n) = 20 + 3\,(171 + 20) = 20 + 573 = 593$.
CNOTs: 2 per edge per layer, $2p|E| = 2\cdot3\cdot57 = 342$. (The remaining $593 - 342 = 251$ are $20$ $H$, $171$ $R_Z$ and $60$ $R_X$.)
For the default ($n = 15$, $|E| = 30$, $p = 2$) the same formula gives $225$ gates and $120$ CNOTs, as in the recorded statistics.

\textbf{\color{brand}Why the other options are incorrect.}
\begin{itemize}[leftmargin=2.4em, itemsep=1pt, topsep=2pt]
\item[\textbf{(A)}] $573$ omits the initial layer of $n = 20$ Hadamards.
\item[\textbf{(C)}] $171$ counts one CNOT per edge per layer; each edge needs two.
\item[\textbf{(D)}] $402$ and $228$ are the counts for $p = 2$, not $p = 3$.
\end{itemize}
\textbf{\color{brand}Key takeaway.} Two-qubit gates grow as $2p|E|$ --- the dominant cost on hardware.
\end{solbox}
\qsep

\needspace{13\baselineskip}
\begin{qitem}{46}
\qmeta{Matching lists}{Default circuit statistics}
For the notebook's default run ($n = 15$ nodes, $|E| = 30$ edges, $p = 2$), match each quantity (List-I) with its value
(List-II).
\par\smallskip\noindent{\renewcommand{\arraystretch}{1.35}%
\begin{tabularx}{\linewidth}{@{}c>{\raggedright\arraybackslash}p{0.33\linewidth}@{\hspace{1.2em}}cL@{}}
\toprule
\multicolumn{2}{@{}l}{\textbf{List-I}} & \multicolumn{2}{l@{}}{\textbf{List-II}} \\ \midrule
\textbf{(P)} & number of variational parameters & \textbf{(1)} & $105$ \\
\textbf{(Q)} & number of $R_Z$ gates & \textbf{(2)} & $120$ \\
\textbf{(R)} & number of single-qubit gates ($H$, $R_Z$ and $R_X$ together) & \textbf{(3)} & $4$ \\
\textbf{(S)} & number of CNOT (\texttt{cx}) gates & \textbf{(4)} & $60$ \\
\bottomrule
\end{tabularx}}
\par\smallskip\noindent\begin{tabularx}{\linewidth}{@{}*{2}{L}@{}}
\optlabel{A}\,P$\to$4, Q$\to$3, R$\to$1, S$\to$2 & \optlabel{B}\,P$\to$3, Q$\to$4, R$\to$2, S$\to$1 \\[3pt]
\correctlabel{C}\,P$\to$3, Q$\to$4, R$\to$1, S$\to$2 & \optlabel{D}\,P$\to$3, Q$\to$2, R$\to$1, S$\to$4
\end{tabularx}
\end{qitem}
\begin{solbox}{C}
\textbf{\color{brand}Correct answer.} \textcolor{okgreen}{\textbf{(C)}} P$\to$3, Q$\to$4, R$\to$1, S$\to$2

\textbf{\color{brand}Concepts required.} Counting parameters and gates of the default circuit.

\textbf{\color{brand}Relevant formulae.} parameters $2p$; $R_Z$: $p|E|$; single-qubit: $n + p|E| + pn$; CNOT: $2p|E|$.

\textbf{\color{brand}Step-by-step explanation.} \begin{itemize}\itemsep0pt
\item P: $2p = 4$ angles $(\gamma_1, \gamma_2, \beta_1, \beta_2)$ (3).
\item Q: one $R_Z$ per edge per layer: $p|E| = 60$ (4).
\item R: $n$ Hadamards $+ p|E|$ $R_Z$ $+ pn$ $R_X$ $= 15 + 60 + 30 = 105$ (1).
\item S: two CNOTs per edge per layer: $2p|E| = 120$ (2).
\end{itemize}
These agree with the recorded \texttt{gate counts by type}: \texttt{h 15, cx 120, rz 60, rx 30} ($225$ gates in total).

\textbf{\color{brand}Why the other options are incorrect.}
\begin{itemize}[leftmargin=2.4em, itemsep=1pt, topsep=2pt]
\item[\textbf{(A)}] P$\to$4 would mean 60 parameters; the circuit has only $2p = 4$ angles regardless of the graph size.
\item[\textbf{(B)}] R$\to$2, S$\to$1 calls the 120 CNOTs single-qubit gates.
\item[\textbf{(D)}] Q$\to$2, S$\to$4 swaps $R_Z$ and CNOT counts: there are two CNOTs but only one $R_Z$ per edge per layer.
\end{itemize}
\textbf{\color{brand}Key takeaway.} Parameters: $2p$ (tiny); gates: proportional to $p(|E| + n)$.
\end{solbox}
\qsep

\needspace{13\baselineskip}
\begin{qitem}{47}
\qmeta{Single-correct MCQ}{$p = 1$ single edge: formula, optimum}
For a single unit edge and $p = 1$ in the notebook's convention, the expected cut is
$F_1(\gamma, \beta) = \tfrac12 - \tfrac12\sin4\beta\,\sin\gamma$. Its maximum value, and one point where it is reached, are
\par\smallskip\noindent\begin{tabularx}{\linewidth}{@{}*{2}{L}@{}}
\optlabel{A}\,$\tfrac34$, at $(\gamma, \beta) = (\pi/2,\, -\pi/8)$ & \optlabel{B}\,$\tfrac12$, reached for all angles \\[3pt]
\optlabel{C}\,$1$, at $(\gamma, \beta) = (\pi/2,\, \pi/8)$ & \correctlabel{D}\,$1$, at $(\gamma, \beta) = (\pi/2,\, -\pi/8)$
\end{tabularx}
\end{qitem}
\begin{solbox}{D}
\textbf{\color{brand}Correct answer.} \textcolor{okgreen}{\textbf{(D)}} $1$, at $(\gamma, \beta) = (\pi/2,\, -\pi/8)$

\textbf{\color{brand}Concepts required.} Heisenberg-picture derivation of $\langle Z_0Z_1\rangle$ for one edge and $p = 1$; sign convention.

\textbf{\color{brand}Relevant formulae.} $U_M^\dagger Z U_M = \cos2\beta\,Z + \sin2\beta\,Y$; $\ U_C^\dagger Y_0Z_1U_C = \cos\gamma\,Y_0Z_1 + \sin\gamma\,X_0$; $\ F_1 = \tfrac12(1 - \sin4\beta\sin\gamma)$.

\textbf{\color{brand}Step-by-step explanation.} \textbf{Derivation of the formula.} $\langle Z_0Z_1\rangle = \bra{+ +}U_C^\dagger U_M^\dagger Z_0Z_1 U_MU_C\ket{+ +}$.
Mixer (Q.31 on each qubit, $c = \cos2\beta$, $s = \sin2\beta$): $U_M^\dagger Z_0Z_1U_M = (cZ_0 + sY_0)(cZ_1 + sY_1)$.
Cost $U_C = e^{-i\gamma Z_0Z_1/2}$: $Z_0Z_1$ and $Y_0Y_1$ commute with it; $Y_0Z_1$ anticommutes, so
$U_C^\dagger Y_0Z_1U_C = Y_0Z_1e^{-i\gamma Z_0Z_1} = \cos\gamma\,Y_0Z_1 - i\sin\gamma\,Y_0Z_0 = \cos\gamma\,Y_0Z_1 + \sin\gamma\,X_0$
(using $YZ = iX$); likewise $Z_0Y_1 \to \cos\gamma\,Z_0Y_1 + \sin\gamma\,X_1$. In $\ket{++}$ only $X_0$, $X_1$ have non-zero
expectation ($=1$), so $\langle Z_0Z_1\rangle = 2cs\sin\gamma = \sin4\beta\sin\gamma$ and $F_1 = \tfrac12(1 - \sin4\beta\sin\gamma)$.
\textbf{Maximum.} $F_1 = 1$ needs $\sin4\beta\sin\gamma = -1$: e.g.\ $\gamma = \pi/2$, $\beta = -\pi/8$ (or $\gamma = -\pi/2$, $\beta = \pi/8$).
Then the edge is cut with certainty: for one edge, $p = 1$ already solves Max-Cut.

\textbf{\color{brand}Why the other options are incorrect.}
\begin{itemize}[leftmargin=2.4em, itemsep=1pt, topsep=2pt]
\item[\textbf{(A)}] At $(\pi/2, -\pi/8)$ the value is $1$, not $\frac34$; $\frac34$ is the ring value (Q.48).
\item[\textbf{(B)}] $F_1$ does depend on the angles; it ranges over $[0, 1]$.
\item[\textbf{(C)}] $(\pi/2, \pi/8)$ gives $\sin4\beta\sin\gamma = +1$ and $F_1 = 0$: the edge is \emph{never} cut (Q.40). Textbooks using $e^{-i\gamma C}$ quote this point, but the notebook's convention flips the sign of $\gamma$.
\end{itemize}
\textbf{\color{brand}Key takeaway.} In the notebook's convention $\gamma$ and $\beta$ need opposite signs to increase the cut.
\end{solbox}
\qsep

\needspace{13\baselineskip}
\begin{qitem}{48}
\qmeta{Single-correct MCQ}{$p = 1$ on the ring}
With optimally chosen angles, the best approximation ratio of $p = 1$ QAOA on the ring (cycle) $C_8$ is
\par\smallskip\noindent\begin{tabularx}{\linewidth}{@{}*{4}{L}@{}}
\optlabel{A}\,$1$ & \correctlabel{B}\,$3/4$ & \optlabel{C}\,$5/6$ & \optlabel{D}\,$0.6924$
\end{tabularx}
\end{qitem}
\begin{solbox}{B}
\textbf{\color{brand}Correct answer.} \textcolor{okgreen}{\textbf{(B)}} $3/4$

\textbf{\color{brand}Concepts required.} $p = 1$ formula for triangle-free regular graphs; symmetry between the edges of a ring.

\textbf{\color{brand}Relevant formulae.} $\langle C_{uv}\rangle = \tfrac12 - \tfrac14\sin4\beta\,\sin\gamma\,(\cos^{d_u-1}\gamma + \cos^{d_v-1}\gamma)$; ring: $d = 2$.

\textbf{\color{brand}Step-by-step explanation.} The ring is 2-regular and triangle-free. Each edge has $d_u = d_v = 2$, so (formula printed in the passage for Q.59 and Q.60,
which holds for every triangle-free graph) $\langle Z_uZ_v\rangle = \sin4\beta\sin\gamma\cos\gamma = \tfrac12\sin4\beta\,\sin2\gamma$, and per edge
\[ \langle C_{uv}\rangle = \tfrac12 - \tfrac14\sin4\beta\,\sin2\gamma \le \tfrac12 + \tfrac14 = \tfrac34, \]
reached at $\sin4\beta\sin2\gamma = -1$, e.g.\ $\gamma = \pi/4$, $\beta = -\pi/8$. All 8 edges are equivalent:
$F_1 = 8\cdot\tfrac34 = 6$ and $C_{\max} = 8$ (even ring, bipartite): $r = 6/8 = 3/4$. A numerical optimisation over the angles
gives the same value.

\textbf{\color{brand}Why the other options are incorrect.}
\begin{itemize}[leftmargin=2.4em, itemsep=1pt, topsep=2pt]
\item[\textbf{(A)}] $1$ is impossible at $p = 1$: $\frac14|\sin4\beta\sin2\gamma| \le \frac14$ caps each edge at $\frac34$.
\item[\textbf{(C)}] $5/6$ is the ring value at $p = 2$ (Q.62).
\item[\textbf{(D)}] $0.6924$ is the 3-regular worst-case bound, not the ring value.
\end{itemize}
\textbf{\color{brand}Key takeaway.} Rings: $r_p = \frac{2p+1}{2p+2}$ for $p = 1, 2, 3$: $\frac34, \frac56, \frac78$.
\end{solbox}
\qsep

\needspace{13\baselineskip}
\begin{qitem}{49}
\qmeta{Statement I / II}{$0.6924$ bound; monotonicity in $p$}
\textbf{Statement I:} For $p = 1$ with optimal angles, QAOA achieves $\langle C\rangle \ge 0.6924\,C_{\max}$ on every 3-regular graph.\\[1mm]
\textbf{Statement II:} For every graph, the optimal value $F_p^* = \max_{\boldsymbol\gamma,\boldsymbol\beta}F_p$ satisfies $F_{p+1}^* \ge F_p^*$.
\begin{itemize}[leftmargin=2.6em, labelsep=0.4em, itemsep=2pt, topsep=3pt]
\item[\correctlabel{A}] Both Statement I and Statement II are true
\item[\optlabel{B}] Both Statement I and Statement II are false
\item[\optlabel{C}] Statement I is true, but Statement II is false
\item[\optlabel{D}] Statement I is false, but Statement II is true
\end{itemize}
\end{qitem}
\begin{solbox}{A}
\textbf{\color{brand}Correct answer.} \textcolor{okgreen}{\textbf{(A)}} Both Statement I and Statement II are true

\textbf{\color{brand}Concepts required.} Known performance guarantee at $p = 1$; nesting of QAOA circuits of different depth.

\textbf{\color{brand}Relevant formulae.} $F_1^* \ge 0.6924\,C_{\max}$ (3-regular); $\ F_{p+1}(\boldsymbol\gamma, 0, \boldsymbol\beta, 0) = F_p(\boldsymbol\gamma, \boldsymbol\beta)$ $\Rightarrow$ $F_{p+1}^* \ge F_p^*$.

\textbf{\color{brand}Step-by-step explanation.} \textbf{I is true:} Farhi, Goldstone and Gutmann (2014) showed that $p = 1$ QAOA gives at least $0.6924\,C_{\max}$ on
3-regular graphs (worst case over local neighbourhoods).
\textbf{II is true:} the depth-$(p+1)$ circuit contains the depth-$p$ circuit as a special case: setting $\gamma_{p+1} = \beta_{p+1} = 0$
makes the extra layer the identity. So the maximum over a larger set cannot be smaller. (With a \emph{finite} optimizer
budget the \emph{found} value can still be worse --- more parameters are harder to optimise.) Answer: (A).

\textbf{\color{brand}Why the other options are incorrect.}
\begin{itemize}[leftmargin=2.4em, itemsep=1pt, topsep=2pt]
\item[\textbf{(B)}] needs both statements false; both are true.
\item[\textbf{(C)}] needs I true and II false; II is true.
\item[\textbf{(D)}] needs I false and II true; I is true.
\end{itemize}
\textbf{\color{brand}Key takeaway.} Optimal values never decrease with depth; found values can, if the optimizer struggles.
\end{solbox}
\qsep

\needspace{13\baselineskip}
\begin{qitem}{50}
\qmeta{Single-correct MCQ}{Limit $p \to \infty$}
Which statement about the limit of large depth $p$ (angles optimised for each $p$) is correct?
\begin{itemize}[leftmargin=2.6em, labelsep=0.4em, itemsep=2pt, topsep=3pt]
\item[\optlabel{A}] $F_p^*/C_{\max} \to 0.878$, the Goemans--Williamson ratio
\item[\optlabel{B}] $F_p^* = C_{\max}$ already at $p = 1$ for every graph
\item[\optlabel{C}] $F_p^* \to W/2$, because deep circuits randomise the state
\item[\correctlabel{D}] $F_p^* \to C_{\max}$ as $p \to \infty$, because a Trotterised adiabatic path is one admissible choice of angles
\end{itemize}
\end{qitem}
\begin{solbox}{D}
\textbf{\color{brand}Correct answer.} \textcolor{okgreen}{\textbf{(D)}} $F_p^* \to C_{\max}$ as $p \to \infty$, because a Trotterised adiabatic path is one admissible choice of angles

\textbf{\color{brand}Concepts required.} Adiabatic theorem and Trotterisation as a special QAOA schedule.

\textbf{\color{brand}Relevant formulae.} $e^{-iT\int H(s)ds} \approx \prod_{l=1}^{p} e^{-i\beta_lH_M}e^{-i\gamma_lH_C}$ with small steps; $\ \lim_{p\to\infty}F_p^* = C_{\max}$.

\textbf{\color{brand}Step-by-step explanation.} Farhi et al.\ argue: a slow interpolation from the mixer to the cost Hamiltonian (adiabatic theorem), discretised into
small time steps (Trotterisation), is a QAOA circuit with special angles. As $p \to \infty$ these special angles drive the
initial state to an optimal split, so the \emph{optimal} angles do at least as well: $F_p^* \to C_{\max}$.
$0.878$ is a classical polynomial-time guarantee, not a QAOA limit; $p = 1$ is exact only for special graphs (one edge, the
triangle).

\textbf{\color{brand}Why the other options are incorrect.}
\begin{itemize}[leftmargin=2.4em, itemsep=1pt, topsep=2pt]
\item[\textbf{(A)}] $0.878$ is the classical Goemans--Williamson guarantee; it is not a limit of QAOA.
\item[\textbf{(B)}] $p = 1$ is exact only for special graphs (one edge, triangle); e.g.\ rings reach only $\frac34$ (Q.48).
\item[\textbf{(C)}] Deep circuits with \emph{optimised} angles do not randomise the state; random angles would.
\end{itemize}
\textbf{\color{brand}Key takeaway.} QAOA is exact in the infinite-depth limit, approximate at any fixed depth.
\end{solbox}
\qsep

\needspace{13\baselineskip}
\begin{qitem}{51}
\qmeta{Single-correct MCQ}{Parameter-vector layout}
In the notebook the angle vector is $\texttt{params} = [\gamma_1, \dots, \gamma_p, \beta_1, \dots, \beta_p]$; the kernel reads
\texttt{gamma = params[layer]} and \texttt{beta = params[layers + layer]}. For $p = 3$ and
\texttt{params = [a, b, c, d, e, f]}, the mixer angle $\beta_2$ of the second layer is
\par\smallskip\noindent\begin{tabularx}{\linewidth}{@{}*{4}{L}@{}}
\optlabel{A}\,$d$ & \optlabel{B}\,$b$ & \correctlabel{C}\,$e$ & \optlabel{D}\,$f$
\end{tabularx}
\end{qitem}
\begin{solbox}{C}
\textbf{\color{brand}Correct answer.} \textcolor{okgreen}{\textbf{(C)}} $e$

\textbf{\color{brand}Concepts required.} Blocked parameter layout and zero-based layer index in the kernel.

\textbf{\color{brand}Relevant formulae.} $\gamma_{l+1} = $ \texttt{params[l]}, $\ \beta_{l+1} = $ \texttt{params[p + l]}, $l = 0, \dots, p-1$.

\textbf{\color{brand}Step-by-step explanation.} Layers are numbered \texttt{layer = 0, 1, 2} in the code. The second layer is \texttt{layer = 1}:
\texttt{gamma = params[1] = b} and \texttt{beta = params[3 + 1] = params[4] = e}.
So $\beta_2 = e$ (and $\gamma_2 = b$). The vector is \emph{blocked} (all $\gamma$ first, then all $\beta$), not interleaved
$(\gamma_1, \beta_1, \gamma_2, \beta_2, \dots)$.

\textbf{\color{brand}Why the other options are incorrect.}
\begin{itemize}[leftmargin=2.4em, itemsep=1pt, topsep=2pt]
\item[\textbf{(A)}] $d$ would be right for an interleaved layout $(\gamma_1, \beta_1, \gamma_2, \beta_2, \dots)$, which the notebook does not use.
\item[\textbf{(B)}] $b$ is $\gamma_2$, the cost angle of the second layer.
\item[\textbf{(D)}] $f$ is $\beta_3$, the mixer angle of the third layer.
\end{itemize}
\textbf{\color{brand}Key takeaway.} \texttt{params} = all $\gamma$'s, then all $\beta$'s.
\end{solbox}
\qsep

\needspace{13\baselineskip}
\begin{qitem}{52}
\qmeta{Single-correct MCQ}{The ramp initialisation}
With \texttt{INITIAL\_PARAMETERS = "ramp"} the notebook starts from \texttt{np.linspace(0.2, 0.6, p)} for the $\gamma$'s and
\texttt{np.linspace(0.6, 0.2, p)} for the $\beta$'s. For $p = 5$, the initial values of $\gamma_3$ and $\beta_2$ are
\par\smallskip\noindent\begin{tabularx}{\linewidth}{@{}*{4}{L}@{}}
\optlabel{A}\,$0.3$ and $0.5$ & \optlabel{B}\,$0.36$ and $0.52$ & \optlabel{C}\,$0.4$ and $0.3$ & \correctlabel{D}\,$0.4$ and $0.5$
\end{tabularx}
\end{qitem}
\begin{solbox}{D}
\textbf{\color{brand}Correct answer.} \textcolor{okgreen}{\textbf{(D)}} $0.4$ and $0.5$

\textbf{\color{brand}Concepts required.} \texttt{np.linspace} includes both end points.

\textbf{\color{brand}Relevant formulae.} \texttt{linspace(a, b, p)}$_k = a + k\frac{b - a}{p - 1}$, $k = 0, \dots, p-1$.

\textbf{\color{brand}Step-by-step explanation.} \texttt{np.linspace(a, b, p)} returns $p$ equally spaced values \emph{including both end points}; the step is $(b - a)/(p - 1)$.
$\gamma$: step $(0.6 - 0.2)/4 = 0.1$: $(0.2, 0.3, 0.4, 0.5, 0.6)$, so $\gamma_3 = 0.4$.
$\beta$: $(0.6, 0.5, 0.4, 0.3, 0.2)$, so $\beta_2 = 0.5$.
(A step of $0.4/5 = 0.08$, which excludes an end point, would give $0.36$ and $0.52$.)

\textbf{\color{brand}Why the other options are incorrect.}
\begin{itemize}[leftmargin=2.4em, itemsep=1pt, topsep=2pt]
\item[\textbf{(A)}] $0.3$ is $\gamma_2$, not $\gamma_3$.
\item[\textbf{(B)}] $0.36$ and $0.52$ use a step of $(b - a)/p$, which would not end exactly at $b$.
\item[\textbf{(C)}] $0.3$ is $\beta_4$; $\beta$ decreases from $0.6$, so $\beta_2 = 0.5$.
\end{itemize}
\textbf{\color{brand}Key takeaway.} The ramp has rising $\gamma$ and falling $\beta$ with step $0.4/(p-1)$.
\end{solbox}
\qsep

\needspace{13\baselineskip}
\begin{qitem}{53}
\qmeta{Assertion--Reason}{Minimise $\langle H_C\rangle$ = maximise $\langle C\rangle$}
\textbf{Assertion (A):} The notebook's objective function returns $\langle H_C\rangle$ and SciPy \emph{minimises} it; this is
equivalent to \emph{maximising} the expected cut.\\[1mm]
\textbf{Reason (R):} $\langle C\rangle = W/2 - \langle H_C\rangle$, and $W/2$ does not depend on the angles.
\begin{itemize}[leftmargin=2.6em, labelsep=0.4em, itemsep=2pt, topsep=3pt]
\item[\correctlabel{A}] Both Assertion (A) and Reason (R) are true, and R is the correct explanation of A
\item[\optlabel{B}] Both A and R are true, but R is \emph{not} the correct explanation of A
\item[\optlabel{C}] A is true, but R is false
\item[\optlabel{D}] A is false, but R is true
\end{itemize}
\end{qitem}
\begin{solbox}{A}
\textbf{\color{brand}Correct answer.} \textcolor{okgreen}{\textbf{(A)}} Both Assertion (A) and Reason (R) are true, and R is the correct explanation of A

\textbf{\color{brand}Concepts required.} Minimising $\langle H_C\rangle$ $\Leftrightarrow$ maximising $\langle C\rangle$.

\textbf{\color{brand}Relevant formulae.} $\langle C\rangle = \tfrac W2 - \langle H_C\rangle$ (code: \texttt{history\_cut.append(total\_weight / 2 - energy)}).

\textbf{\color{brand}Step-by-step explanation.} \textbf{R is true:} $[z_i \neq z_j] = \tfrac12(1 - s_is_j)$, so $C = \tfrac W2 I - H_C$ as operators, and taking expectations gives R.
\textbf{A is true because of R:} $\langle C\rangle$ is a constant minus $\langle H_C\rangle$, so the angles that minimise $\langle H_C\rangle$ are exactly
those that maximise $\langle C\rangle$. The code records \texttt{history\_cut.append(total\_weight / 2 - energy)} after each
evaluation and returns \texttt{energy} to \texttt{minimize}. R is the correct explanation. Answer: (A).

\textbf{\color{brand}Why the other options are incorrect.}
\begin{itemize}[leftmargin=2.4em, itemsep=1pt, topsep=2pt]
\item[\textbf{(B)}] needs A and R both true but R not explaining A; R is precisely the reason.
\item[\textbf{(C)}] needs A true and R false; R is true.
\item[\textbf{(D)}] needs A false and R true; A is true.
\end{itemize}
\textbf{\color{brand}Key takeaway.} A constant shift and a sign turn energy minimisation into cut maximisation.
\end{solbox}
\qsep

\needspace{13\baselineskip}
\begin{qitem}{54}
\qmeta{Single-correct MCQ}{Sign of the default start point}
At the notebook's default starting point $\boldsymbol\gamma = (0.2, 0.6)$, $\boldsymbol\beta = (0.6, 0.2)$, an exact
state-vector calculation for its 15-node graph gives $\langle C\rangle \approx 9.18$, \emph{below} the random baseline $W/2 = 15$.
The best explanation is
\begin{itemize}[leftmargin=2.6em, labelsep=0.4em, itemsep=2pt, topsep=3pt]
\item[\correctlabel{A}] In the convention $e^{-i\beta H_M}e^{-i\gamma H_C}$, the start state $\ket+^{\otimes n}$ is the \emph{highest} eigenstate of $H_M$, so an annealing-like ramp with $\gamma, \beta > 0$ heads for the highest eigenstate of $H_C$, the \emph{smallest} cut; heading for the largest cut needs $\gamma$ and $\beta$ of opposite sign
\item[\optlabel{B}] The value is meaningless: \texttt{cudaq.observe} returns random numbers before the optimizer has converged
\item[\optlabel{C}] Every $p = 2$ starting point gives $\langle C\rangle < W/2$; only $p \ge 3$ can beat random guessing
\item[\optlabel{D}] It indicates a bug, because a QAOA state can never have $\langle C\rangle < W/2$
\end{itemize}
\end{qitem}
\begin{solbox}{A}
\textbf{\color{brand}Correct answer.} \textcolor{okgreen}{\textbf{(A)}} In the convention $e^{-i\beta H_M}e^{-i\gamma H_C}$, the start state $\ket+^{\otimes n}$ is the \emph{highest} eigenstate of $H_M$, so an annealing-like ramp with $\gamma, \beta > 0$ heads for the highest eigenstate of $H_C$, the \emph{smallest} cut; heading for the largest cut needs $\gamma$ and $\beta$ of opposite sign

\textbf{\color{brand}Concepts required.} Sign conventions; adiabatic interpretation; the start state is the top eigenstate of $H_M$.

\textbf{\color{brand}Relevant formulae.} Small angles, one edge: $\Delta\langle C\rangle \approx -2\beta\gamma$ (from $\tfrac12 - \tfrac12\sin4\beta\sin\gamma$); adiabatic evolution preserves the position in the spectrum.

\textbf{\color{brand}Step-by-step explanation.} \textbf{Small-angle check.} For one edge $F_1 = \tfrac12 - \tfrac12\sin4\beta\sin\gamma$ (Q.47): with $\gamma > 0$ and small $\beta > 0$ the
cut \emph{decreases}; it increases only if $\sin4\beta\,\sin\gamma < 0$.
\textbf{Annealing picture.} $e^{-i\beta H_M}e^{-i\gamma H_C}$ with positive angles is one Trotter step of evolution under
$\beta H_M + \gamma H_C$. $\ket+^{\otimes n}$ is the \emph{top} eigenstate of $H_M$ (Q.34); an adiabatic ramp keeps the top
eigenstate, which at the end is the top eigenstate of $H_C$ --- the \emph{minimum} cut. For the maximum cut one needs
$H_M \to -H_C$, i.e.\ $\gamma$ and $\beta$ of opposite sign.
\textbf{Not a bug.} $\langle C\rangle$ can be anywhere between $C_{\min}$ and $C_{\max}$; the optimizer (COBYLA) moves away from the poor
start, and the recorded run still ended at $\langle C\rangle = 20.53$ ($r = 0.82$). The value $9.18$ is computed from the notebook's
own conventions (\texttt{verify\_answers.py} reproduces it with NumPy and CUDA-Q).

\textbf{\color{brand}Why the other options are incorrect.}
\begin{itemize}[leftmargin=2.4em, itemsep=1pt, topsep=2pt]
\item[\textbf{(B)}] \texttt{cudaq.observe} is deterministic on a state-vector simulator: it returns the exact $\langle H_C\rangle$ at the given angles.
\item[\textbf{(C)}] Many $p = 2$ angles beat $W/2$ --- the optimizer reached $20.53$ with $p = 2$.
\item[\textbf{(D)}] QAOA states can have any $\langle C\rangle$ between $C_{\min}$ and $C_{\max}$; values below $W/2$ are allowed.
\end{itemize}
\textbf{\color{brand}Key takeaway.} Check the sign of the start: with $e^{-i\gamma H_C}$ choose $\gamma\beta < 0$ to head for large cuts.
\end{solbox}
\qsep

\needspace{13\baselineskip}
\begin{qitem}{55}
\qmeta{Matching lists}{Algorithm $\leftrightarrow$ code}
Match each part of the QAOA algorithm (List-I) with the code that implements it in the notebook (List-II).
\par\smallskip\noindent{\renewcommand{\arraystretch}{1.35}%
\begin{tabularx}{\linewidth}{@{}c>{\raggedright\arraybackslash}p{0.33\linewidth}@{\hspace{1.2em}}cL@{}}
\toprule
\multicolumn{2}{@{}l}{\textbf{List-I}} & \multicolumn{2}{l@{}}{\textbf{List-II}} \\ \midrule
\textbf{(P)} & initial state $\ket+^{\otimes n}$ & \textbf{(1)} & \texttt{rx(2.0 * beta, qubits[q])} \\
\textbf{(Q)} & cost unitary for one edge & \textbf{(2)} & \texttt{cudaq.observe(kernels["state"], H\_C, ...).expectation()} \\
\textbf{(R)} & mixer unitary for one qubit & \textbf{(3)} & \texttt{h(qubits)} \\
\textbf{(S)} & objective $\langle H_C\rangle$ & \textbf{(4)} & \texttt{x.ctrl(...); rz(gamma * weights[e], ...); x.ctrl(...)} \\
\bottomrule
\end{tabularx}}
\par\smallskip\noindent\begin{tabularx}{\linewidth}{@{}*{2}{L}@{}}
\optlabel{A}\,P$\to$3, Q$\to$1, R$\to$4, S$\to$2 & \optlabel{B}\,P$\to$3, Q$\to$4, R$\to$2, S$\to$1 \\[3pt]
\correctlabel{C}\,P$\to$3, Q$\to$4, R$\to$1, S$\to$2 & \optlabel{D}\,P$\to$4, Q$\to$3, R$\to$1, S$\to$2
\end{tabularx}
\end{qitem}
\begin{solbox}{C}
\textbf{\color{brand}Correct answer.} \textcolor{okgreen}{\textbf{(C)}} P$\to$3, Q$\to$4, R$\to$1, S$\to$2

\textbf{\color{brand}Concepts required.} Mapping the algorithm's building blocks to CUDA-Q calls.

\textbf{\color{brand}Relevant formulae.} $\ket+^{\otimes n} = H^{\otimes n}\ket0^{\otimes n}$; $\ e^{-i\gamma\frac w2Z_iZ_j} = $ CNOT $R_Z(\gamma w)$ CNOT; $\ e^{-i\beta X} = R_X(2\beta)$.

\textbf{\color{brand}Step-by-step explanation.} \begin{itemize}\itemsep0pt
\item P: \texttt{h(qubits)} applies $H$ to every qubit of the freshly allocated $\ket0^{\otimes n}$, giving $\ket+^{\otimes n}$ (3).
\item Q: CNOT, $R_Z(\gamma w)$, CNOT $= e^{-i\gamma(w/2)Z_iZ_j}$ (4).
\item R: $R_X(2\beta) = e^{-i\beta X}$ (1).
\item S: \texttt{cudaq.observe} computes $\langle\psi|H_C|\psi\rangle$ for the state-preparation kernel (exactly on simulators, from shots
when \texttt{shots\_count} is given) (2).
\end{itemize}

\textbf{\color{brand}Why the other options are incorrect.}
\begin{itemize}[leftmargin=2.4em, itemsep=1pt, topsep=2pt]
\item[\textbf{(A)}] Q$\to$1, R$\to$4 swaps cost and mixer: the mixer contains no two-qubit gates.
\item[\textbf{(B)}] R$\to$2, S$\to$1: \texttt{cudaq.observe} computes an expectation value; it does not apply gates to the circuit.
\item[\textbf{(D)}] P$\to$4, Q$\to$3: the Hadamard layer is the initial state; it is not the cost layer.
\end{itemize}
\textbf{\color{brand}Key takeaway.} \texttt{h} $\to$ start, \texttt{x.ctrl/rz/x.ctrl} $\to$ cost, \texttt{rx} $\to$ mixer, \texttt{observe} $\to$ objective.
\end{solbox}
\qsep

\needspace{13\baselineskip}
\begin{qitem}{56}
\qmeta{Single-correct MCQ}{Triangle solved at $p = 1$}
For the triangle $K_3$ (unit weights, $C_{\max} = 2$), the best value of $F_1$ over all angles is
\par\smallskip\noindent\begin{tabularx}{\linewidth}{@{}*{2}{L}@{}}
\optlabel{A}\,$\sqrt3$, i.e.\ $r \approx 0.866$ & \correctlabel{B}\,$2$, i.e.\ $r = 1$ \\[3pt]
\optlabel{C}\,$1.75$, i.e.\ $r = 0.875$ & \optlabel{D}\,$1.5$, i.e.\ $r = 0.75$
\end{tabularx}
\end{qitem}
\begin{solbox}{B}
\textbf{\color{brand}Correct answer.} \textcolor{okgreen}{\textbf{(B)}} $2$, i.e.\ $r = 1$

\textbf{\color{brand}Concepts required.} When $H_C$ has only two eigenvalues, one layer can rotate the state fully into the optimal subspace.

\textbf{\color{brand}Relevant formulae.} $H_C = \tfrac32\Pi_{\{000,111\}} - \tfrac12\Pi_{\text{opt}}$; optimum $(\gamma^*, \beta^*) = (\arctan\tfrac{1}{\sqrt2}, -\tfrac12\arctan\tfrac1{\sqrt2})$, $F_1 = 2$.

\textbf{\color{brand}Step-by-step explanation.} $H_C$ has only two eigenvalues (Q.9): $+\tfrac32$ on $\{000, 111\}$ and $-\tfrac12$ on the six optimal strings. Hence
$e^{-i\gamma H_C}$ only puts a relative phase $e^{-2i\gamma}$ between the two non-optimal strings and the six optimal ones, and the
mixer can then interfere the two non-optimal amplitudes away completely. A numerical maximisation gives $F_1 = 2$ exactly, for
example at $\gamma^* = \arctan(1/\sqrt2) \approx 0.6155$, $\beta^* = -\gamma^*/2 \approx -0.3077$ (notebook convention; checked
in \texttt{verify\_answers.py}). So $p = 1$ already finds the optimum with probability 1. The value $1.5 = W/2$ is only the
random baseline.

\textbf{\color{brand}Why the other options are incorrect.}
\begin{itemize}[leftmargin=2.4em, itemsep=1pt, topsep=2pt]
\item[\textbf{(A)}] $\sqrt3$ is not attained as a maximum; numerical optimisation reaches the full $2$.
\item[\textbf{(C)}] $1.75$ is not the maximum either; the triangle is solved exactly at $p = 1$.
\item[\textbf{(D)}] $1.5 = W/2$ is the random baseline (all angles zero).
\end{itemize}
\textbf{\color{brand}Key takeaway.} Small, highly symmetric instances can be solved exactly at $p = 1$; rings cannot.
\end{solbox}
\qsep

\needspace{13\baselineskip}
\begin{qitem}{57}
\qmeta{Single-correct MCQ}{$r$ vs.\ best sample}
The previously recorded run reports: approximation ratio $r = \langle C\rangle/C_{\max} = 0.8213$, best sampled cut $25$
($= C_{\max}$), and $P(\text{optimal split}) = 0.06$ with $4000$ shots. Which interpretation is correct?
\begin{itemize}[leftmargin=2.6em, labelsep=0.4em, itemsep=2pt, topsep=3pt]
\item[\optlabel{A}] $82\,\%$ of the shots returned an optimal split
\item[\correctlabel{B}] $r$ measures the \emph{average} cut of the distribution; the best of $4000$ samples can still be optimal, and here $240$ shots were optimal
\item[\optlabel{C}] The best sampled cut can never exceed $\langle C\rangle = 20.53$, so the value $25$ indicates an error
\item[\optlabel{D}] The best sampled split has cut $0.8213 \times 25 \approx 20.5$
\end{itemize}
\end{qitem}
\begin{solbox}{B}
\textbf{\color{brand}Correct answer.} \textcolor{okgreen}{\textbf{(B)}} $r$ measures the \emph{average} cut of the distribution; the best of $4000$ samples can still be optimal, and here $240$ shots were optimal

\textbf{\color{brand}Concepts required.} Expectation versus best sample; probability of the optimum; QAOA as estimator and as sampler.

\textbf{\color{brand}Relevant formulae.} $r = \langle C\rangle/C_{\max}$; $\ \mathbb E[\#\text{optimal shots}] = P_{\text{opt}}\,S = 0.06\times4000 = 240$.

\textbf{\color{brand}Step-by-step explanation.} $\langle C\rangle = \sum_z |\braket{z|\boldsymbol\gamma,\boldsymbol\beta}|^2\,C(z) = 20.53$ is an average over the whole output distribution;
$r = 20.53/25 = 0.82$ summarises that average. Individual samples scatter around it: some are better than average.
Among $4000$ shots, $P_{\text{opt}} = 0.06$ means about $0.06\times4000 = 240$ optimal shots (exactly $240$ were recorded), so the best
sample reaches $C_{\max}$. QAOA used as a \emph{sampler} is judged by the best sample (and $P_{\text{opt}}$), used as an \emph{estimator}
by $r$.

\textbf{\color{brand}Why the other options are incorrect.}
\begin{itemize}[leftmargin=2.4em, itemsep=1pt, topsep=2pt]
\item[\textbf{(A)}] $82\,\%$ of shots optimal would mean $P_{\text{opt}} = 0.82$; the recorded $P_{\text{opt}}$ is $0.06$.
\item[\textbf{(C)}] Individual samples scatter around the mean; the maximum of 4000 samples is typically far above the mean.
\item[\textbf{(D)}] $0.8213\times25$ is the \emph{expected} cut $20.53$, not the best sample.
\end{itemize}
\textbf{\color{brand}Key takeaway.} Report both: $r$ (average quality) and the best sample with $P_{\text{opt}}$ (what a user actually gets).
\end{solbox}
\qsep

\needspace{13\baselineskip}
\begin{qitem}{58}
\qmeta{Statement I / II}{Best evaluated angles; no global guarantee}
\textbf{Statement I:} The notebook's final sampling run uses the angles of the best evaluation stored in
\texttt{history\_cut}, which need not be the optimizer's final point \texttt{opt.x}.\\[1mm]
\textbf{Statement II:} With COBYLA and a budget of 80 evaluations, the notebook is guaranteed to find the global maximum of $F_2$.
\begin{itemize}[leftmargin=2.6em, labelsep=0.4em, itemsep=2pt, topsep=3pt]
\item[\optlabel{A}] Both Statement I and Statement II are true
\item[\optlabel{B}] Both Statement I and Statement II are false
\item[\correctlabel{C}] Statement I is true, but Statement II is false
\item[\optlabel{D}] Statement I is false, but Statement II is true
\end{itemize}
\end{qitem}
\begin{solbox}{C}
\textbf{\color{brand}Correct answer.} \textcolor{okgreen}{\textbf{(C)}} Statement I is true, but Statement II is false

\textbf{\color{brand}Concepts required.} The notebook's choice of final angles; local optimizers give no global guarantee.

\textbf{\color{brand}Relevant formulae.} \texttt{best\_index = argmax(history\_cut)}; \texttt{best\_params = history\_params[best\_index]}.

\textbf{\color{brand}Step-by-step explanation.} \textbf{I is true:} the code computes \texttt{best\_index = int(np.argmax(history\_cut))} and samples with
\texttt{best\_params = history\_params[best\_index]}; this protects against an optimizer that ends on a worse point.
\textbf{II is false:} COBYLA is a local, derivative-free method; QAOA landscapes have many local optima, and the recorded run
stopped with ``\emph{the objective function has been evaluated MAXFUN times}'' --- i.e.\ because of the budget, not because a global
optimality condition was met. Answer: (C).

\textbf{\color{brand}Why the other options are incorrect.}
\begin{itemize}[leftmargin=2.4em, itemsep=1pt, topsep=2pt]
\item[\textbf{(A)}] needs both statements true; II is false.
\item[\textbf{(B)}] needs both statements false; I is true.
\item[\textbf{(D)}] needs I false and II true; I is true and II is false.
\end{itemize}
\textbf{\color{brand}Key takeaway.} COBYLA stopped at its evaluation budget; ``converged'' and ``globally optimal'' are not the same.
\end{solbox}
\qsep

\needspace{20\baselineskip}
\begin{passagebox}{Paragraph for Questions 59 and 60}
QAOA with $p = 1$ is applied to the path graph $0 - 1 - 2$ (edges $(0,1)$ and $(1,2)$, unit weights,
$C_{\max} = 2$) in the notebook's convention, with $\gamma = \pi/2$ and $\beta = -\pi/8$. For a triangle-free graph and $p = 1$
the edge correlator in this convention is
\[ \langle Z_uZ_v\rangle = \tfrac12\,\sin4\beta\,\sin\gamma\,\big(\cos^{d_u - 1}\gamma + \cos^{d_v - 1}\gamma\big), \]
where $d_u$, $d_v$ are the degrees of the two end nodes.
\end{passagebox}
\needspace{13\baselineskip}
\begin{qitem}{59}
\qmeta{Paragraph based}{Path graph: $\langle Z_0Z_1\rangle$}
The value of $\langle Z_0Z_1\rangle$ is
\par\smallskip\noindent\begin{tabularx}{\linewidth}{@{}*{4}{L}@{}}
\correctlabel{A}\,$-\tfrac12$ & \optlabel{B}\,$-1$ & \optlabel{C}\,$+\tfrac12$ & \optlabel{D}\,$0$
\end{tabularx}
\end{qitem}
\begin{solbox}{A}
\textbf{\color{brand}Correct answer.} \textcolor{okgreen}{\textbf{(A)}} $-\tfrac12$

\textbf{\color{brand}Concepts required.} Evaluating the $p = 1$ triangle-free formula for nodes of different degree.

\textbf{\color{brand}Relevant formulae.} $\langle Z_uZ_v\rangle = \tfrac12\sin4\beta\sin\gamma(\cos^{d_u-1}\gamma + \cos^{d_v-1}\gamma)$; $\ \cos^0\gamma = 1$.

\textbf{\color{brand}Step-by-step explanation.} Edge $(0,1)$: $d_0 = 1$ (node 0 has one neighbour), $d_1 = 2$. With $\gamma = \pi/2$: $\sin\gamma = 1$, $\cos\gamma = 0$, so
$\cos^{0}\gamma = 1$ and $\cos^{1}\gamma = 0$. With $\beta = -\pi/8$: $\sin4\beta = \sin(-\pi/2) = -1$.
\[ \langle Z_0Z_1\rangle = \tfrac12\cdot(-1)\cdot1\cdot(1 + 0) = -\tfrac12 . \]
(Note $0^0 = 1$: the term $\cos^{d_u-1}\gamma$ for a leaf node is $1$ for every $\gamma$.) By symmetry $\langle Z_1Z_2\rangle = -\tfrac12$ too.

\textbf{\color{brand}Why the other options are incorrect.}
\begin{itemize}[leftmargin=2.4em, itemsep=1pt, topsep=2pt]
\item[\textbf{(B)}] $-1$ uses $\cos^{d-1}\gamma = 1$ for both ends; for node 1 ($d = 2$) it is $\cos\gamma = 0$.
\item[\textbf{(C)}] $+\frac12$ has the wrong sign: $\sin4\beta = \sin(-\pi/2) = -1$.
\item[\textbf{(D)}] $0$ treats $\cos^0\gamma$ as $0$ when $\cos\gamma = 0$; but $x^0 = 1$.
\end{itemize}
\textbf{\color{brand}Key takeaway.} Leaves contribute the factor $1$; neighbours with $\cos\gamma = 0$ switch the other term off.
\end{solbox}
\qsep

\needspace{13\baselineskip}
\begin{qitem}{60}
\qmeta{Paragraph based}{Path graph: $\langle C\rangle$, $P_{\text{opt}}$}
The expected cut $\langle C\rangle$ and the probability of sampling an optimal split are
\par\smallskip\noindent\begin{tabularx}{\linewidth}{@{}*{4}{L}@{}}
\optlabel{A}\,$1.5$ and $0.3125$ & \correctlabel{B}\,$1.5$ and $0.625$ & \optlabel{C}\,$1.25$ and $0.625$ & \optlabel{D}\,$2$ and $1$
\end{tabularx}
\end{qitem}
\begin{solbox}{B}
\textbf{\color{brand}Correct answer.} \textcolor{okgreen}{\textbf{(B)}} $1.5$ and $0.625$

\textbf{\color{brand}Concepts required.} Summing edge contributions; probability of the optimal set; complementary optimal strings.

\textbf{\color{brand}Relevant formulae.} $\langle C\rangle = \sum_e\tfrac12(1 - \langle Z_uZ_v\rangle)$; $\ P_{\text{opt}} = P(010) + P(101)$.

\textbf{\color{brand}Step-by-step explanation.} $\langle C\rangle = \sum_{\text{edges}}\tfrac12(1 - \langle Z_uZ_v\rangle) = 2\cdot\tfrac12\big(1 + \tfrac12\big) = 1.5$, so $r = 1.5/2 = 0.75$.
The optimal splits of the path are $010$ and $101$ (both edges cut). A state-vector calculation gives
$P(010) = P(101) = 0.3125$ and $0.0625$ for each of the other six strings, so $P_{\text{opt}} = 0.625$.
Consistency check: $\langle C\rangle = 2\cdot0.625 + 1\cdot P(C = 1) = 1.25 + 1\cdot 4\cdot0.0625 = 1.5$ (the strings $001, 011, 100, 110$ cut one
edge; $000, 111$ cut none).

\textbf{\color{brand}Why the other options are incorrect.}
\begin{itemize}[leftmargin=2.4em, itemsep=1pt, topsep=2pt]
\item[\textbf{(A)}] $0.3125$ counts only one of the two optimal strings ($010$ and its complement $101$).
\item[\textbf{(C)}] $1.25$ adds only the optimal strings' contribution and forgets the four strings that cut one edge.
\item[\textbf{(D)}] $2$ and $1$ would need $\langle Z_uZ_v\rangle = -1$ on both edges; one layer is not enough for the path.
\end{itemize}
\textbf{\color{brand}Key takeaway.} $r = 0.75$ but $P_{\text{opt}} = 0.625$: a sampler can do better than the average suggests.
\end{solbox}
\qsep

\newpage

\begin{sectionbanner}SECTION D \enspace$\cdot$\enspace Advanced QAOA and Implementation \hfill Q.61--80\end{sectionbanner}

\needspace{13\baselineskip}
\begin{qitem}{61}
\qmeta{Single-correct MCQ}{Light cone of an edge}
On a ring graph, the expectation $\langle Z_uZ_v\rangle$ for one edge $(u, v)$ of a depth-$p$ QAOA state depends on at most how
many qubits when $p = 3$?
\par\smallskip\noindent\begin{tabularx}{\linewidth}{@{}*{4}{L}@{}}
\correctlabel{A}\,$8$ & \optlabel{B}\,all $n$ qubits & \optlabel{C}\,$7$ & \optlabel{D}\,$6$
\end{tabularx}
\end{qitem}
\begin{solbox}{A}
\textbf{\color{brand}Correct answer.} \textcolor{okgreen}{\textbf{(A)}} $8$

\textbf{\color{brand}Concepts required.} Heisenberg-picture light cone of a local observable in a local circuit.

\textbf{\color{brand}Relevant formulae.} Support after $p$ layers on a ring: $2 + 2p$ qubits (one new node on each side per cost layer).

\textbf{\color{brand}Step-by-step explanation.} Work in the Heisenberg picture: $\langle Z_uZ_v\rangle = \bra{+}^{\otimes n}U^\dagger Z_uZ_vU\ket{+}^{\otimes n}$. Conjugating by a mixer layer
(single-qubit gates) keeps the support. Conjugating by a cost layer can only add the neighbours of the current support (an
$e^{-i\gamma Z_aZ_b/2}$ acts trivially unless it touches the support). Each layer enlarges the support by one node on each side
of the ring: $\{u, v\}$ (2 qubits) $\to$ 4 $\to$ 6 $\to$ 8 after $p = 3$ layers, i.e.\ $2p + 2 = 8$ qubits.
This \emph{light cone} is why fixed-$p$ QAOA is local, and why large instances of bounded degree can be analysed with small
subgraphs.

\textbf{\color{brand}Why the other options are incorrect.}
\begin{itemize}[leftmargin=2.4em, itemsep=1pt, topsep=2pt]
\item[\textbf{(B)}] Gates far from the edge cancel in $U^\dagger Z_uZ_vU$; only the light cone matters.
\item[\textbf{(C)}] $7$ counts the edge's two nodes plus only one side growing per layer.
\item[\textbf{(D)}] $6$ is the light cone for $p = 2$.
\end{itemize}
\textbf{\color{brand}Key takeaway.} At fixed $p$, QAOA is local: an edge sees only its $p$-neighbourhood.
\end{solbox}
\qsep

\needspace{13\baselineskip}
\begin{qitem}{62}
\qmeta{Single-correct MCQ}{Ring at $p = 2$: $5/6$}
For even rings $C_n$ with $n \ge 2p + 2$, numerical optimisation of the angles gives the best $p = 2$ approximation ratio
\par\smallskip\noindent\begin{tabularx}{\linewidth}{@{}*{4}{L}@{}}
\optlabel{A}\,$3/4$ & \correctlabel{B}\,$5/6$ & \optlabel{C}\,$7/8$ & \optlabel{D}\,$1$
\end{tabularx}
\end{qitem}
\begin{solbox}{B}
\textbf{\color{brand}Correct answer.} \textcolor{okgreen}{\textbf{(B)}} $5/6$

\textbf{\color{brand}Concepts required.} Light-cone independence of $n$; numerical optimisation of ring QAOA.

\textbf{\color{brand}Relevant formulae.} Ring values (numerical): $r_1 = \frac34$, $r_2 = \frac56$, $r_3 = \frac78$, consistent with $r_p = \frac{2p+1}{2p+2}$.

\textbf{\color{brand}Step-by-step explanation.} For $n \ge 2p + 2$ the light cone of an edge (previous question) is a path of $2p + 2$ nodes that does not close around the
ring, so every edge sees the same local problem and the ratio does not depend on $n$. Optimising $F_2$ numerically for
$C_6$, $C_8$ and $C_{10}$ (\texttt{verify\_answers.py}) gives $r = 0.83333 = 5/6$; for $p = 1$ it gives $3/4$ (Q.48) and for $p = 3$, $7/8$.
These values follow the pattern $r_p = \frac{2p+1}{2p+2}$ conjectured by Farhi et al.\ for the ring: the ratio improves with depth
but reaches $1$ only as $p \to \infty$ for large rings.

\textbf{\color{brand}Why the other options are incorrect.}
\begin{itemize}[leftmargin=2.4em, itemsep=1pt, topsep=2pt]
\item[\textbf{(A)}] $3/4$ is the $p = 1$ value (Q.48).
\item[\textbf{(C)}] $7/8$ is the $p = 3$ value.
\item[\textbf{(D)}] $1$ would need $\langle Z_uZ_v\rangle = -1$ for every edge, which a light cone of 6 qubits cannot achieve on a large ring.
\end{itemize}
\textbf{\color{brand}Key takeaway.} On rings each extra layer closes part of the remaining gap, but never all of it.
\end{solbox}
\qsep

\needspace{13\baselineskip}
\begin{qitem}{63}
\qmeta{Single-correct MCQ}{Parameter shift vs.\ finite differences}
For the default circuit ($n = 15$, $|E| = 30$, $p = 2$), compare two ways of computing the gradient of $F_2$ with respect to
all $2p$ angles: (i) the exact parameter-shift rule applied to every rotation gate separately (each $R_Z$ and $R_X$ gate is shifted
by $\pm\pi/2$ on its own), and (ii) a forward finite-difference gradient as used by L-BFGS-B. The numbers of circuit evaluations are
\par\smallskip\noindent\begin{tabularx}{\linewidth}{@{}*{4}{L}@{}}
\optlabel{A}\,$8$ and $5$ & \optlabel{B}\,$180$ and $4$ & \optlabel{C}\,$90$ and $5$ & \correctlabel{D}\,$180$ and $5$
\end{tabularx}
\end{qitem}
\begin{solbox}{D}
\textbf{\color{brand}Correct answer.} \textcolor{okgreen}{\textbf{(D)}} $180$ and $5$

\textbf{\color{brand}Concepts required.} Parameter-shift rule per gate and the chain rule; finite differences.

\textbf{\color{brand}Relevant formulae.} $\partial_\theta f = \tfrac12[f(\theta + \tfrac\pi2) - f(\theta - \tfrac\pi2)]$ per gate; $\ \partial_{\gamma_l}F = \sum_e w_e\partial_{\theta_e}F$; forward difference: $1 + 2p$ runs.

\textbf{\color{brand}Step-by-step explanation.} \textbf{(i)} $\gamma_l$ enters $|E|$ gates $R_Z(\gamma_lw_e)$; by the chain rule $\partial F/\partial\gamma_l = \sum_e w_e\,\partial F/\partial\theta_e$, and each
$\partial F/\partial\theta_e$ needs two runs (shifted by $\pm\pi/2$, Q.16). So $2|E|$ runs per $\gamma_l$. Likewise $\beta_l$ enters $n$ gates
$R_X(2\beta_l)$: $2n$ runs. Total $2p(|E| + n) = 2\cdot2\cdot(30 + 15) = 180$.
A single two-term shift of $\gamma_l$ as a whole is \emph{not} exact: its generator $H_C$ has many different eigenvalue gaps.
\textbf{(ii)} $f(\boldsymbol\theta)$ once plus one shifted run per parameter: $2p + 1 = 5$ (the notebook's table: ``$2p+1$ circuit runs per gradient'').
Finite differences are cheaper but approximate and sensitive to shot noise.

\textbf{\color{brand}Why the other options are incorrect.}
\begin{itemize}[leftmargin=2.4em, itemsep=1pt, topsep=2pt]
\item[\textbf{(A)}] $8$ applies a two-term shift to each of the $2p$ angles as a whole, which is not exact because $\gamma_l$ multiplies all $|E|$ gates.
\item[\textbf{(B)}] $4$ forgets the evaluation at the unshifted point.
\item[\textbf{(C)}] $90$ counts one run per gate instead of two shifted runs.
\end{itemize}
\textbf{\color{brand}Key takeaway.} Exact gradients of shared angles cost $2\times$(number of gates that use them).
\end{solbox}
\qsep

\needspace{13\baselineskip}
\begin{qitem}{64}
\qmeta{Single-correct MCQ}{COBYLA \texttt{rhobeg}}
The notebook optimises the angles with SciPy's \texttt{minimize(..., method="COBYLA")} and the options
\texttt{\{"maxiter": 80, "rhobeg": 0.2\}}. The option \texttt{rhobeg = 0.2} sets
\begin{itemize}[leftmargin=2.6em, labelsep=0.4em, itemsep=2pt, topsep=3pt]
\item[\optlabel{A}] a learning rate multiplying the gradient
\item[\optlabel{B}] the maximum number of iterations divided by 400
\item[\optlabel{C}] the tolerance at which COBYLA declares convergence
\item[\correctlabel{D}] the initial size of the trust region, i.e.\ the first steps change the angles by about $0.2$ rad
\end{itemize}
\end{qitem}
\begin{solbox}{D}
\textbf{\color{brand}Correct answer.} \textcolor{okgreen}{\textbf{(D)}} the initial size of the trust region, i.e.\ the first steps change the angles by about $0.2$ rad

\textbf{\color{brand}Concepts required.} COBYLA: derivative-free trust-region optimisation with linear models.

\textbf{\color{brand}Relevant formulae.} Trust-region radius $\rho$: starts at \texttt{rhobeg}, shrinks to \texttt{tol}; each step stays within $\rho$ of the current point.

\textbf{\color{brand}Step-by-step explanation.} COBYLA (Constrained Optimization BY Linear Approximation) is derivative-free: it builds linear models of the objective
from function values at a simplex of points and takes steps inside a trust region of radius $\rho$. \texttt{rhobeg} is the
starting radius: the first probes move each angle by about $0.2$ rad. $\rho$ then shrinks towards \texttt{tol}/\texttt{rhoend} as the
search converges. No gradients (and therefore no learning rate) are used. With only 80 evaluations a too small \texttt{rhobeg}
would explore little; a too large one would jump across the oscillating landscape.

\textbf{\color{brand}Why the other options are incorrect.}
\begin{itemize}[leftmargin=2.4em, itemsep=1pt, topsep=2pt]
\item[\textbf{(A)}] COBYLA uses no gradients, so there is no learning rate.
\item[\textbf{(B)}] \texttt{maxiter} is a separate option ($80$).
\item[\textbf{(C)}] The final accuracy is set by \texttt{tol} (rhoend), not by \texttt{rhobeg}.
\end{itemize}
\textbf{\color{brand}Key takeaway.} \texttt{rhobeg} is the first step size; choose it comparable to the scale of the landscape's features.
\end{solbox}
\qsep

\needspace{13\baselineskip}
\begin{qitem}{65}
\qmeta{Single-correct MCQ}{Shot noise of the objective}
With \texttt{OBJECTIVE\_SHOTS = 1000}, $\langle H_C\rangle$ is estimated from 1000 shots. If the cut values of the sampled
bitstrings have a standard deviation of $2.5$, the standard error of the estimated objective is about
\par\smallskip\noindent\begin{tabularx}{\linewidth}{@{}*{4}{L}@{}}
\optlabel{A}\,$0.0063$ & \optlabel{B}\,$0.0025$ & \correctlabel{C}\,$0.079$ & \optlabel{D}\,$2.5$
\end{tabularx}
\end{qitem}
\begin{solbox}{C}
\textbf{\color{brand}Correct answer.} \textcolor{okgreen}{\textbf{(C)}} $0.079$

\textbf{\color{brand}Concepts required.} Standard error of a mean estimated from shots.

\textbf{\color{brand}Relevant formulae.} $\mathrm{SE} = \sigma/\sqrt S$.

\textbf{\color{brand}Step-by-step explanation.} The shot estimate is the sample mean of $H_C(z) = \tfrac W2 - C(z)$; its spread equals that of $C(z)$, $\sigma = 2.5$.
Standard error of a mean of $S$ independent samples: $\sigma/\sqrt S = 2.5/\sqrt{1000} = 2.5/31.6 \approx 0.079$.
Differences between neighbouring optimizer steps near the optimum are often of this size or smaller, so the optimizer starts
to chase noise; that is why the notebook warns that shot-based objectives ``make optimisation harder''.

\textbf{\color{brand}Why the other options are incorrect.}
\begin{itemize}[leftmargin=2.4em, itemsep=1pt, topsep=2pt]
\item[\textbf{(A)}] $0.0063 = \sigma^2/S$ is the variance of the mean, not its standard error.
\item[\textbf{(B)}] $0.0025 = \sigma/S$ divides by $S$ instead of $\sqrt S$.
\item[\textbf{(D)}] $2.5$ is the spread of single shots, not of their mean.
\end{itemize}
\textbf{\color{brand}Key takeaway.} $1000$ shots give an objective accurate to about $\pm0.08$ --- noise the optimizer must live with.
\end{solbox}
\qsep

\needspace{13\baselineskip}
\begin{qitem}{66}
\qmeta{Single-correct MCQ}{Brute-force memory}
For $n = 22$ (\texttt{BRUTE\_FORCE\_MAX\_NODES}), the notebook's \texttt{brute\_force\_maxcut} creates the arrays
\texttt{index = np.arange(2**N, dtype=np.int64)} and \texttt{values = np.zeros(2**N)}. Together they occupy about
\par\smallskip\noindent\begin{tabularx}{\linewidth}{@{}*{4}{L}@{}}
\optlabel{A}\,$34$ MB & \correctlabel{B}\,$67$ MB & \optlabel{C}\,$0.5$ MB & \optlabel{D}\,$4.2$ MB
\end{tabularx}
\end{qitem}
\begin{solbox}{B}
\textbf{\color{brand}Correct answer.} \textcolor{okgreen}{\textbf{(B)}} $67$ MB

\textbf{\color{brand}Concepts required.} Memory of NumPy arrays: entries $\times$ bytes per entry.

\textbf{\color{brand}Relevant formulae.} $2^{22} = 4\,194\,304$; int64 and float64: $8$ bytes; total $2\times8\times2^{22}$ bytes.

\textbf{\color{brand}Step-by-step explanation.} $2^{22} = 4\,194\,304$ entries per array. \texttt{int64} and \texttt{float64} both use $8$ bytes:
$2 \times 8 \times 4\,194\,304 = 67\,108\,864$ bytes $\approx 67$ MB (plus temporaries of the same size created per edge while
evaluating \texttt{(index >> a) \^{} (index >> b)}). $34$ MB is one array; $4.2$ MB is the number of \emph{entries} read as bytes.
Each additional node doubles both memory and time, which is why the exact reference stops at 22 nodes.

\textbf{\color{brand}Why the other options are incorrect.}
\begin{itemize}[leftmargin=2.4em, itemsep=1pt, topsep=2pt]
\item[\textbf{(A)}] $34$ MB is one of the two arrays.
\item[\textbf{(C)}] $0.5$ MB is the 15-qubit state vector, unrelated.
\item[\textbf{(D)}] $4.2$ MB is the number of entries ($4.2$ million), not bytes.
\end{itemize}
\textbf{\color{brand}Key takeaway.} Brute force costs $O(2^n)$ memory and $O(|E|\,2^n)$ time: fine at 22 nodes, hopeless at 40.
\end{solbox}
\qsep

\needspace{13\baselineskip}
\begin{qitem}{67}
\qmeta{Single-correct MCQ}{Multi-GPU state vector}
A 36-qubit QAOA state vector is simulated in fp64. If it is split evenly over a power-of-two number of 80 GB GPUs (as the
multi-GPU state-vector target does), ignoring workspace, the minimum number of GPUs is
\par\smallskip\noindent\begin{tabularx}{\linewidth}{@{}*{4}{L}@{}}
\optlabel{A}\,$8$ & \optlabel{B}\,$14$ & \optlabel{C}\,$32$ & \correctlabel{D}\,$16$
\end{tabularx}
\end{qitem}
\begin{solbox}{D}
\textbf{\color{brand}Correct answer.} \textcolor{okgreen}{\textbf{(D)}} $16$

\textbf{\color{brand}Concepts required.} Distributing a state vector over a power-of-two number of devices.

\textbf{\color{brand}Relevant formulae.} memory $= 16\times2^{36}$ B $\approx 1.10$ TB; per GPU $= 1.10$ TB$/2^k \le 80$ GB $\Rightarrow k \ge 4$.

\textbf{\color{brand}Step-by-step explanation.} Memory: $16\times2^{36} = 1.0995\times10^{12}$ bytes $\approx 1.10$ TB. Dividing by $80$ GB: $13.7$, so at least $14$ GPUs by memory
alone. The state is distributed in power-of-two slices ($2^k$ GPUs each holding $2^{36-k}$ amplitudes): $2^3 = 8$ GPUs would need
$137$ GB each (too much), $2^4 = 16$ GPUs need $68.7$ GB each (fits). So 16 GPUs, launched with one MPI rank per GPU
(\texttt{BACKEND = "multi-gpu"}, \texttt{mpiexec}). These GPU paths are implemented in the notebook but were not executed on the
CPU-only test machine.

\textbf{\color{brand}Why the other options are incorrect.}
\begin{itemize}[leftmargin=2.4em, itemsep=1pt, topsep=2pt]
\item[\textbf{(A)}] $8$ GPUs would need $137$ GB each.
\item[\textbf{(B)}] $14$ GPUs suffice by memory, but the state is split into power-of-two slices.
\item[\textbf{(C)}] $32$ is more than necessary; $16$ already fit.
\end{itemize}
\textbf{\color{brand}Key takeaway.} Each extra qubit doubles the number of GPUs needed.
\end{solbox}
\qsep

\needspace{13\baselineskip}
\begin{qitem}{68}
\qmeta{Matching lists}{Backend / method $\to$ target}
Match each setting of \texttt{BACKEND} and \texttt{SIMULATION\_METHOD} in \texttt{QAOA.ipynb} (List-I) with what the notebook's
\texttt{resolve\_backend} does (List-II).
\par\smallskip\noindent{\renewcommand{\arraystretch}{1.35}%
\begin{tabularx}{\linewidth}{@{}c>{\raggedright\arraybackslash}p{0.33\linewidth}@{\hspace{1.2em}}cL@{}}
\toprule
\multicolumn{2}{@{}l}{\textbf{List-I}} & \multicolumn{2}{l@{}}{\textbf{List-II}} \\ \midrule
\textbf{(P)} & \texttt{"cpu"} + \texttt{"statevector"} & \textbf{(1)} & target \texttt{tensornet-mps} \\
\textbf{(Q)} & \texttt{"cpu"} + \texttt{"density\_matrix"} & \textbf{(2)} & stops with an error: the MPS simulator needs an NVIDIA GPU \\
\textbf{(R)} & \texttt{"gpu"} + \texttt{"mps"} & \textbf{(3)} & target \texttt{qpp-cpu} \\
\textbf{(S)} & \texttt{"cpu"} + \texttt{"mps"} & \textbf{(4)} & target \texttt{density-matrix-cpu} \\
\bottomrule
\end{tabularx}}
\par\smallskip\noindent\begin{tabularx}{\linewidth}{@{}*{2}{L}@{}}
\optlabel{A}\,P$\to$4, Q$\to$3, R$\to$1, S$\to$2 & \optlabel{B}\,P$\to$3, Q$\to$1, R$\to$4, S$\to$2 \\[3pt]
\optlabel{C}\,P$\to$3, Q$\to$4, R$\to$2, S$\to$1 & \correctlabel{D}\,P$\to$3, Q$\to$4, R$\to$1, S$\to$2
\end{tabularx}
\end{qitem}
\begin{solbox}{D}
\textbf{\color{brand}Correct answer.} \textcolor{okgreen}{\textbf{(D)}} P$\to$3, Q$\to$4, R$\to$1, S$\to$2

\textbf{\color{brand}Concepts required.} The notebook's backend table and its explicit error for unsupported combinations.

\textbf{\color{brand}Relevant formulae.} \texttt{\_TARGET\_TABLE[(backend, method)]} $\to$ CUDA-Q target; missing keys raise \texttt{ValueError} with \texttt{\_UNSUPPORTED\_HINTS}.

\textbf{\color{brand}Step-by-step explanation.} From \texttt{\_TARGET\_TABLE} and \texttt{\_UNSUPPORTED\_HINTS} in the shared utilities (QAOA allows the methods
\texttt{statevector}, \texttt{mps}, \texttt{tensornet}, \texttt{density\_matrix}):
\begin{itemize}\itemsep0pt
\item P: \texttt{("cpu", "statevector")} $\to$ \texttt{qpp-cpu} (OpenMP state vector, fp64) (3) --- the target of the recorded run.
\item Q: \texttt{("cpu", "density\_matrix")} $\to$ \texttt{density-matrix-cpu} (4); memory $16\times4^n$ bytes.
\item R: \texttt{("gpu", "mps")} $\to$ \texttt{tensornet-mps} (1).
\item S: \texttt{("cpu", "mps")} is not in the table; the hint says ``CUDA-Q's MPS simulator (target 'tensornet-mps') needs an
NVIDIA GPU'' and a \texttt{ValueError} is raised (2).
\end{itemize}

\textbf{\color{brand}Why the other options are incorrect.}
\begin{itemize}[leftmargin=2.4em, itemsep=1pt, topsep=2pt]
\item[\textbf{(A)}] P$\to$4, Q$\to$3 swaps the two CPU targets; \texttt{density\_matrix} selects \texttt{density-matrix-cpu}.
\item[\textbf{(B)}] Q$\to$1 maps the density matrix to an MPS target; these are different simulation methods.
\item[\textbf{(C)}] R$\to$2, S$\to$1 reverses which MPS combination works: MPS (\texttt{tensornet-mps}) needs a GPU.
\end{itemize}
\textbf{\color{brand}Key takeaway.} CPU: \texttt{qpp-cpu} / \texttt{density-matrix-cpu}; GPU: \texttt{nvidia}, \texttt{tensornet-mps}, \texttt{tensornet}.
\end{solbox}
\qsep

\needspace{13\baselineskip}
\begin{qitem}{69}
\qmeta{Single-correct MCQ}{When MPS is efficient}
For which instance is a matrix-product-state (MPS) simulation of a depth-$p$ QAOA circuit most likely to be efficient (small
bond dimension), with the qubits ordered along a line $0, 1, \dots, n-1$?
\begin{itemize}[leftmargin=2.6em, labelsep=0.4em, itemsep=2pt, topsep=3pt]
\item[\optlabel{A}] Any graph, because QAOA states are always product states before measurement
\item[\optlabel{B}] A complete graph $K_n$: its symmetry makes the MPS exact with bond dimension 1
\item[\correctlabel{C}] A ring $C_n$ with small $p$: every cut of the line is crossed by only 2 edges
\item[\optlabel{D}] A dense random graph with edge probability $0.5$: more edges average out the entanglement
\end{itemize}
\end{qitem}
\begin{solbox}{C}
\textbf{\color{brand}Correct answer.} \textcolor{okgreen}{\textbf{(C)}} A ring $C_n$ with small $p$: every cut of the line is crossed by only 2 edges

\textbf{\color{brand}Concepts required.} Entanglement across a bipartition bounds the MPS bond dimension.

\textbf{\color{brand}Relevant formulae.} $\chi_k \le 2^{p\,m_k}$, $m_k$ = number of edges crossing the cut between qubits $k$ and $k+1$.

\textbf{\color{brand}Step-by-step explanation.} An MPS across a cut between qubits $k$ and $k+1$ needs a bond dimension at least the Schmidt rank there. Each cost gate
$e^{-i\theta Z_aZ_b/2}$ that straddles the cut can at most double that rank, so after $p$ layers $\chi \le 2^{p\cdot m_k}$, where $m_k$ is the
number of edges crossing the cut. For a ring in natural order $m_k = 2$ for every cut, so $\chi \le 4^p$: small for small $p$.
Dense random graphs have $m_k = \Theta(n^2)$ crossings: the bond dimension can grow exponentially. $K_n$ is not product-like
either: its cost layer entangles all pairs (Q.28). The cost layer creates entanglement, so QAOA states are not product states.

\textbf{\color{brand}Why the other options are incorrect.}
\begin{itemize}[leftmargin=2.4em, itemsep=1pt, topsep=2pt]
\item[\textbf{(A)}] The cost layer entangles qubits (Q.28); QAOA states are generally entangled.
\item[\textbf{(B)}] $K_n$ has $\approx n^2/4$ edges across the middle cut; its QAOA state is not a product state.
\item[\textbf{(D)}] More edges mean more entangling gates across each cut, hence \emph{larger} bond dimension.
\end{itemize}
\textbf{\color{brand}Key takeaway.} MPS likes 1D-like graphs at low depth; dense graphs need state vectors or tensor networks.
\end{solbox}
\qsep

\needspace{13\baselineskip}
\begin{qitem}{70}
\qmeta{Assertion--Reason}{Shot-based objective on hardware}
\textbf{Assertion (A):} When \texttt{BACKEND = "hardware"}, the notebook estimates the objective from shots even if
\texttt{OBJECTIVE\_SHOTS = None}.\\[1mm]
\textbf{Reason (R):} Hardware (and hardware-emulation) targets only return measurement samples; they cannot give a state vector or an
exact expectation value.
\begin{itemize}[leftmargin=2.6em, labelsep=0.4em, itemsep=2pt, topsep=3pt]
\item[\correctlabel{A}] Both Assertion (A) and Reason (R) are true, and R is the correct explanation of A
\item[\optlabel{B}] Both A and R are true, but R is \emph{not} the correct explanation of A
\item[\optlabel{C}] A is true, but R is false
\item[\optlabel{D}] A is false, but R is true
\end{itemize}
\end{qitem}
\begin{solbox}{A}
\textbf{\color{brand}Correct answer.} \textcolor{okgreen}{\textbf{(A)}} Both Assertion (A) and Reason (R) are true, and R is the correct explanation of A

\textbf{\color{brand}Concepts required.} Hardware returns samples only; the code switches to shot-based objectives automatically.

\textbf{\color{brand}Relevant formulae.} $\widehat{\langle H_C\rangle} = \tfrac1S\sum_{\text{shots}} H_C(z)$, $\ S = $ \texttt{SHOTS} on hardware.

\textbf{\color{brand}Step-by-step explanation.} \textbf{A is true:} the code sets
\texttt{objective\_shots = OBJECTIVE\_SHOTS if OBJECTIVE\_SHOTS is not None else (SHOTS if backend\_cfg["device"] == "hardware" else None)},
so on hardware $\langle H_C\rangle$ is estimated with \texttt{SHOTS} $= 4000$ shots per evaluation.
\textbf{R is true and is the reason:} a QPU only produces bitstrings; exact expectation values exist only in simulators
(README: ``Hardware targets support sampling only''). Answer: (A).

\textbf{\color{brand}Why the other options are incorrect.}
\begin{itemize}[leftmargin=2.4em, itemsep=1pt, topsep=2pt]
\item[\textbf{(B)}] needs A and R both true but R not explaining A; R is exactly why the code makes this choice.
\item[\textbf{(C)}] needs A true and R false; R is true.
\item[\textbf{(D)}] needs A false and R true; A is true.
\end{itemize}
\textbf{\color{brand}Key takeaway.} On hardware every objective value is a noisy shot estimate.
\end{solbox}
\qsep

\needspace{13\baselineskip}
\begin{qitem}{71}
\qmeta{Single-correct MCQ}{Depolarising-noise model}
A simple global-depolarising noise model treats the noisy QAOA output as a mixture: with probability $F = (1 - \varepsilon)^{N_{\text{CNOT}}}$
the ideal state, otherwise a uniformly random bitstring. For the default circuit ($N_{\text{CNOT}} = 120$, ideal $\langle C\rangle = 20.53$,
$W/2 = 15$) and $\varepsilon = 0.005$ per CNOT, the predicted noisy $\langle C\rangle$ is about
\par\smallskip\noindent\begin{tabularx}{\linewidth}{@{}*{4}{L}@{}}
\optlabel{A}\,$17.2$ & \optlabel{B}\,$15.0$ & \correctlabel{C}\,$18.0$ & \optlabel{D}\,$11.3$
\end{tabularx}
\end{qitem}
\begin{solbox}{C}
\textbf{\color{brand}Correct answer.} \textcolor{okgreen}{\textbf{(C)}} $18.0$

\textbf{\color{brand}Concepts required.} Global depolarising model: mixture of the ideal output with the uniform distribution.

\textbf{\color{brand}Relevant formulae.} $F = (1-\varepsilon)^{N_{\text{CNOT}}}$; $\ \langle C\rangle_{\text{noisy}} = F\langle C\rangle_{\text{ideal}} + (1 - F)\tfrac W2$.

\textbf{\color{brand}Step-by-step explanation.} $F = 0.995^{120} = e^{120\ln0.995} = e^{-0.6015} \approx 0.548$.
Mixture: $\langle C\rangle_{\text{noisy}} = F\,\langle C\rangle_{\text{ideal}} + (1 - F)\,\tfrac W2 = 0.548\cdot20.53 + 0.452\cdot15 = 11.25 + 6.78 \approx 18.0$.
$11.3$ forgets the random part (a random bitstring still cuts $W/2$ on average); $17.2$ uses the linearisation
$F \approx 1 - 120\varepsilon = 0.4$, which is poor when $120\varepsilon$ is not small. The model is illustrative only (real noise is
not exactly global depolarising), but it shows how noise pulls QAOA towards the random baseline and limits useful depth.

\textbf{\color{brand}Why the other options are incorrect.}
\begin{itemize}[leftmargin=2.4em, itemsep=1pt, topsep=2pt]
\item[\textbf{(A)}] $17.2$ uses $F \approx 1 - 120\varepsilon = 0.40$; the exact $F = 0.548$.
\item[\textbf{(B)}] $15.0$ is the fully depolarised limit $F = 0$.
\item[\textbf{(D)}] $11.3$ keeps only $F\langle C\rangle_{\text{ideal}}$; the random part still contributes $W/2$ per shot.
\end{itemize}
\textbf{\color{brand}Key takeaway.} Noise pulls $\langle C\rangle$ towards $W/2$ exponentially in the number of two-qubit gates.
\end{solbox}
\qsep

\needspace{13\baselineskip}
\begin{qitem}{72}
\qmeta{Matching lists}{QAOA variants}
Match each QAOA variant (List-I) with its key idea (List-II).
\par\smallskip\noindent{\renewcommand{\arraystretch}{1.35}%
\begin{tabularx}{\linewidth}{@{}c>{\raggedright\arraybackslash}p{0.33\linewidth}@{\hspace{1.2em}}cL@{}}
\toprule
\multicolumn{2}{@{}l}{\textbf{List-I}} & \multicolumn{2}{l@{}}{\textbf{List-II}} \\ \midrule
\textbf{(P)} & warm-start QAOA & \textbf{(1)} & problem-specific mixers (e.g.\ XY mixers) that keep the state inside the feasible subspace of a constrained problem \\
\textbf{(Q)} & ADAPT-QAOA & \textbf{(2)} & use measured correlations $\langle Z_iZ_j\rangle$ to fix and eliminate variables, then solve the smaller instance \\
\textbf{(R)} & Quantum Alternating Operator Ansatz (Hadfield et al.) & \textbf{(3)} & start from a state biased towards a classical (e.g.\ relaxed/SDP) solution instead of $\ket+^{\otimes n}$ \\
\textbf{(S)} & recursive QAOA (RQAOA) & \textbf{(4)} & grow the mixer layer by layer from an operator pool, choosing the operator with the largest energy gradient \\
\bottomrule
\end{tabularx}}
\par\smallskip\noindent\begin{tabularx}{\linewidth}{@{}*{2}{L}@{}}
\optlabel{A}\,P$\to$3, Q$\to$1, R$\to$4, S$\to$2 & \optlabel{B}\,P$\to$3, Q$\to$4, R$\to$2, S$\to$1 \\[3pt]
\correctlabel{C}\,P$\to$3, Q$\to$4, R$\to$1, S$\to$2 & \optlabel{D}\,P$\to$4, Q$\to$3, R$\to$1, S$\to$2
\end{tabularx}
\end{qitem}
\begin{solbox}{C}
\textbf{\color{brand}Correct answer.} \textcolor{okgreen}{\textbf{(C)}} P$\to$3, Q$\to$4, R$\to$1, S$\to$2

\textbf{\color{brand}Concepts required.} Main families of QAOA extensions.

\textbf{\color{brand}Relevant formulae.} Warm start: $\ket{\psi_0} = \bigotimes_i R_Y(\theta_i)\ket0$ from a relaxed solution; RQAOA: fix $z_i = \operatorname{sign}\langle Z_iZ_j\rangle\,z_j$ for the largest $|\langle Z_iZ_j\rangle|$.

\textbf{\color{brand}Step-by-step explanation.} \begin{itemize}\itemsep0pt
\item P: warm-start QAOA (Egger, Mare\v{c}ek, Woerner 2021) changes the initial state (and matching mixer) using a classical relaxation (3).
\item Q: ADAPT-QAOA (Zhu et al.\ 2022) selects mixers adaptively from a pool by gradient (4).
\item R: the Quantum Alternating Operator Ansatz (Hadfield et al.\ 2019) generalises the mixer to respect constraints (1).
\item S: RQAOA (Bravyi et al.\ 2020) eliminates variables using the strongest correlations, recursively (2).
\end{itemize}
None of these is implemented in \texttt{QAOA.ipynb}, which uses the standard $X$ mixer and $\ket+^{\otimes n}$.

\textbf{\color{brand}Why the other options are incorrect.}
\begin{itemize}[leftmargin=2.4em, itemsep=1pt, topsep=2pt]
\item[\textbf{(A)}] Q$\to$1, R$\to$4: feasibility-preserving mixers are Hadfield's idea; ADAPT grows mixers from a pool.
\item[\textbf{(B)}] R$\to$2, S$\to$1 swaps the constrained-mixer ansatz and RQAOA.
\item[\textbf{(D)}] P$\to$4, Q$\to$3 swaps warm start and ADAPT: ADAPT changes the mixer, warm start changes the initial state.
\end{itemize}
\textbf{\color{brand}Key takeaway.} Variants change the start state, the mixer, or the outer loop --- the notebook implements standard QAOA.
\end{solbox}
\qsep

\needspace{13\baselineskip}
\begin{qitem}{73}
\qmeta{Single-correct MCQ}{INTERP initialisation}
The INTERP initialisation (Zhou et al.\ 2020) builds depth-$(p+1)$ starting angles from optimised depth-$p$ angles by
$\gamma_i^{(p+1)} = \tfrac{i-1}{p}\,\gamma_{i-1}^{(p)} + \tfrac{p-i+1}{p}\,\gamma_i^{(p)}$ for $i = 1, \dots, p+1$ (with $\gamma_0^{(p)} = \gamma_{p+1}^{(p)} = 0$).
From $\boldsymbol\gamma^{(2)} = (0.4, 0.8)$ it gives $\boldsymbol\gamma^{(3)} =$
\par\smallskip\noindent\begin{tabularx}{\linewidth}{@{}*{2}{L}@{}}
\correctlabel{A}\,$(0.4,\ 0.6,\ 0.8)$ & \optlabel{B}\,$(0.27,\ 0.53,\ 0.8)$ \\[3pt]
\optlabel{C}\,$(0.4,\ 0.8,\ 0.8)$ & \optlabel{D}\,$(0.4,\ 0.8,\ 0)$
\end{tabularx}
\end{qitem}
\begin{solbox}{A}
\textbf{\color{brand}Correct answer.} \textcolor{okgreen}{\textbf{(A)}} $(0.4,\ 0.6,\ 0.8)$

\textbf{\color{brand}Concepts required.} INTERP: linear interpolation of an optimised schedule onto $p+1$ points.

\textbf{\color{brand}Relevant formulae.} $\gamma_i^{(p+1)} = \tfrac{i-1}{p}\gamma_{i-1}^{(p)} + \tfrac{p-i+1}{p}\gamma_i^{(p)}$, $\ \gamma_0 = \gamma_{p+1} = 0$.

\textbf{\color{brand}Step-by-step explanation.} $p = 2$.
$i = 1$: $\tfrac02\gamma_0 + \tfrac22\gamma_1 = 0.4$.
$i = 2$: $\tfrac12\gamma_1 + \tfrac12\gamma_2 = \tfrac12(0.4 + 0.8) = 0.6$.
$i = 3$: $\tfrac22\gamma_2 + \tfrac02\gamma_3 = 0.8$.
So $(0.4, 0.6, 0.8)$: the optimised schedule is linearly interpolated onto one more point, keeping its smooth shape. This is a
general strategy from the literature; the notebook itself starts every depth from its fixed ramp.

\textbf{\color{brand}Why the other options are incorrect.}
\begin{itemize}[leftmargin=2.4em, itemsep=1pt, topsep=2pt]
\item[\textbf{(B)}] $(0.27, 0.53, 0.8)$ is a linear ramp from $0$, which discards the optimised first angle.
\item[\textbf{(C)}] $(0.4, 0.8, 0.8)$ copies the last angle; INTERP gives the middle point $0.6$.
\item[\textbf{(D)}] $(0.4, 0.8, 0)$ appends a zero layer --- equivalent to depth $p$, a valid but different start.
\end{itemize}
\textbf{\color{brand}Key takeaway.} Good angles at depth $p$ are smooth schedules; interpolate them to start depth $p+1$.
\end{solbox}
\qsep

\needspace{13\baselineskip}
\begin{qitem}{74}
\qmeta{Statement I / II}{Heuristic reference: upper bound on $r$}
For graphs with more nodes than \texttt{BRUTE\_FORCE\_MAX\_NODES}, the notebook replaces the exact optimum by the best cut found
by \texttt{local\_search\_maxcut} and still reports $r = \langle C\rangle/\texttt{max\_cut}$.\\[1mm]
\textbf{Statement I:} In that case the reported \texttt{max\_cut} is guaranteed to be at least the true $C_{\max}$.\\[1mm]
\textbf{Statement II:} In that case the reported ratio is at least the true ratio $\langle C\rangle/C_{\max}$.
\begin{itemize}[leftmargin=2.6em, labelsep=0.4em, itemsep=2pt, topsep=3pt]
\item[\optlabel{A}] Both Statement I and Statement II are true
\item[\optlabel{B}] Both Statement I and Statement II are false
\item[\optlabel{C}] Statement I is true, but Statement II is false
\item[\correctlabel{D}] Statement I is false, but Statement II is true
\end{itemize}
\end{qitem}
\begin{solbox}{D}
\textbf{\color{brand}Correct answer.} \textcolor{okgreen}{\textbf{(D)}} Statement I is false, but Statement II is true

\textbf{\color{brand}Concepts required.} A heuristic lower bound on $C_{\max}$ makes the reported ratio an upper bound.

\textbf{\color{brand}Relevant formulae.} $\texttt{max\_cut}_{\text{heur}} \le C_{\max}$ $\Rightarrow$ $\langle C\rangle/\texttt{max\_cut}_{\text{heur}} \ge \langle C\rangle/C_{\max}$.

\textbf{\color{brand}Step-by-step explanation.} The heuristic returns the cut of an actual split, so $\texttt{max\_cut} \le C_{\max}$ (it can only miss the optimum, never exceed it).
\textbf{I is false} (the inequality goes the other way).
\textbf{II is true:} dividing the same positive $\langle C\rangle$ by a smaller (or equal) number gives a larger (or equal) value:
$\langle C\rangle/\texttt{max\_cut} \ge \langle C\rangle/C_{\max}$. The reported ratio is an \emph{upper bound} on the true one, which is why
the notebook labels the reference ``NOT guaranteed optimal''. Answer: (D).

\textbf{\color{brand}Why the other options are incorrect.}
\begin{itemize}[leftmargin=2.4em, itemsep=1pt, topsep=2pt]
\item[\textbf{(A)}] needs both statements true; I is false.
\item[\textbf{(B)}] needs both statements false; II is true.
\item[\textbf{(C)}] needs I true and II false; I is false and II is true.
\end{itemize}
\textbf{\color{brand}Key takeaway.} Without an exact reference, the reported ratio is optimistic.
\end{solbox}
\qsep

\needspace{13\baselineskip}
\begin{qitem}{75}
\qmeta{Single-correct MCQ}{Recorded timing per gate}
In the recorded run, the quantum execution time was $40.01$ s for the $80$ objective evaluations (\texttt{cudaq.observe}
calls) of the $225$-gate circuit on 15 qubits. The average time per evaluation and per gate are about
\par\smallskip\noindent\begin{tabularx}{\linewidth}{@{}*{2}{L}@{}}
\optlabel{A}\,$0.010$ s and $44\ \mu$s & \optlabel{B}\,$0.18$ s and $0.8$ ms \\[3pt]
\correctlabel{C}\,$0.50$ s and $2.2$ ms & \optlabel{D}\,$40$ s and $0.18$ s
\end{tabularx}
\end{qitem}
\begin{solbox}{C}
\textbf{\color{brand}Correct answer.} \textcolor{okgreen}{\textbf{(C)}} $0.50$ s and $2.2$ ms

\textbf{\color{brand}Concepts required.} Turning recorded timings into per-evaluation and per-gate costs; comparing with the calibration.

\textbf{\color{brand}Relevant formulae.} $t_{\text{eval}} = t_{\text{exec}}/N_{\text{eval}}$; $\ t_{\text{gate}} = t_{\text{eval}}/N_{\text{gates}}$.

\textbf{\color{brand}Step-by-step explanation.} Per evaluation: $40.01/80 \approx 0.50$ s. Per gate: $0.50/225 \approx 2.2\times10^{-3}$ s $= 2.2$ ms.
This matches the notebook's own calibration printed before execution: ``$1.0921$ ms/gate at 14 qubits, extrapolated $\times2^1$
to 15 qubits'' $\to 2.18$ ms/gate, and the estimate $\sim39.8$ s versus $40.84$ s measured for the whole cell.
($0.010$ s divides by the $4000$ shots, which play no role in exact \texttt{observe} calls.)

\textbf{\color{brand}Why the other options are incorrect.}
\begin{itemize}[leftmargin=2.4em, itemsep=1pt, topsep=2pt]
\item[\textbf{(A)}] Dividing by $4000$ shots is wrong: the exact \texttt{observe} calls use no shots.
\item[\textbf{(B)}] $0.18$ s divides the total by the number of gates instead of the number of evaluations.
\item[\textbf{(D)}] $40$ s is the total, not per evaluation.
\end{itemize}
\textbf{\color{brand}Key takeaway.} About $2.2$ ms per gate at 15 qubits on the 4-core CPU, matching the notebook's own calibration.
\end{solbox}
\qsep

\needspace{13\baselineskip}
\begin{qitem}{76}
\qmeta{Single-correct MCQ}{Run-time scaling to 25 nodes}
The notebook estimates run time as (number of evaluations) $\times$ (gates per circuit) $\times$ (time per gate), with the time per
gate scaling as $2^n$ for a state vector. Going from the default ($n = 15$, $|E| = 30$, $p = 2$) to $n = 25$ with $|E| = 90$ edges at
the same $p$ and evaluation budget, the run time grows by a factor of about
\par\smallskip\noindent\begin{tabularx}{\linewidth}{@{}*{4}{L}@{}}
\optlabel{A}\,$2.9\times10^{6}$ & \correctlabel{B}\,$2.8\times10^{3}$ & \optlabel{C}\,$2.7$ & \optlabel{D}\,$1.0\times10^{3}$
\end{tabularx}
\end{qitem}
\begin{solbox}{B}
\textbf{\color{brand}Correct answer.} \textcolor{okgreen}{\textbf{(B)}} $2.8\times10^{3}$

\textbf{\color{brand}Concepts required.} The notebook's runtime model: gates $\times$ per-gate cost $\propto 2^n$.

\textbf{\color{brand}Relevant formulae.} $\frac{T_{25}}{T_{15}} = \frac{25 + 2(3\cdot90 + 25)}{15 + 2(3\cdot30 + 15)}\times2^{10} = \frac{615}{225}\times1024$.

\textbf{\color{brand}Step-by-step explanation.} Gates: $n + p(3|E| + n)$: $15 + 2(90 + 15) = 225$ before, $25 + 2(270 + 25) = 615$ after; ratio $615/225 = 2.73$.
Time per gate: $2^{25-15} = 1024$ times larger. Total factor $2.73\times1024 \approx 2.8\times10^3$, i.e.\ the recorded
$\approx 40$ s would become $\approx 31$ hours on the same 4-core CPU --- the reason the notebook suggests a GPU for 24--30 nodes.
$1.0\times10^3$ ignores the extra gates; $2.7$ ignores the exponential state size; $2.9\times10^6$ uses the $4^{\Delta n}$ scaling of a
density matrix. (The model is the notebook's rough estimate, quoted as accurate to about $\pm3\times$.)

\textbf{\color{brand}Why the other options are incorrect.}
\begin{itemize}[leftmargin=2.4em, itemsep=1pt, topsep=2pt]
\item[\textbf{(A)}] $2.9\times10^6$ uses $4^{\Delta n}$, the density-matrix scaling; this is a state-vector run.
\item[\textbf{(C)}] $2.7$ is only the gate ratio; it ignores that each gate acts on a $1024\times$ larger state.
\item[\textbf{(D)}] $1.0\times10^3$ is only the $2^{10}$ factor; the circuit also has $2.73\times$ more gates.
\end{itemize}
\textbf{\color{brand}Key takeaway.} Ten more qubits cost three orders of magnitude on a state-vector simulator.
\end{solbox}
\qsep

\needspace{13\baselineskip}
\begin{qitem}{77}
\qmeta{Single-correct MCQ}{Quantum advantage: state of knowledge}
Which statement about QAOA for Max-Cut reflects the current state of knowledge?
\begin{itemize}[leftmargin=2.6em, labelsep=0.4em, itemsep=2pt, topsep=3pt]
\item[\correctlabel{A}] No quantum advantage has been demonstrated: the classical Goemans--Williamson algorithm guarantees $0.878$ in polynomial time, and fixed-depth QAOA is limited by locality
\item[\optlabel{B}] The recorded result $r = 0.82$ on 15 nodes proves a quantum speed-up
\item[\optlabel{C}] QAOA solves Max-Cut exactly in polynomial time for any fixed $p$
\item[\optlabel{D}] QAOA with $p = 1$ beats Goemans--Williamson on every 3-regular graph
\end{itemize}
\end{qitem}
\begin{solbox}{A}
\textbf{\color{brand}Correct answer.} \textcolor{okgreen}{\textbf{(A)}} No quantum advantage has been demonstrated: the classical Goemans--Williamson algorithm guarantees $0.878$ in polynomial time, and fixed-depth QAOA is limited by locality

\textbf{\color{brand}Concepts required.} Known guarantees and limitations of QAOA versus classical algorithms.

\textbf{\color{brand}Relevant formulae.} GW: $0.878$ (polynomial, classical); QAOA $p = 1$, 3-regular: $0.6924$.

\textbf{\color{brand}Step-by-step explanation.} Known facts: $p = 1$ QAOA guarantees $0.6924$ on 3-regular graphs (Q.49), below the Goemans--Williamson $0.878$. Locality
(light cone, Q.61) limits any fixed-$p$ QAOA: e.g.\ Bravyi et al.\ (2020) showed that on certain bipartite regular graphs low-depth
QAOA stays far from the optimum. Max-Cut is NP-hard, so no polynomial exact algorithm is expected. A 15-node simulation that a
laptop brute-forces in $0.01$ s cannot show a speed-up; QAOA's value is as a well-understood, tunable heuristic and a benchmark.

\textbf{\color{brand}Why the other options are incorrect.}
\begin{itemize}[leftmargin=2.4em, itemsep=1pt, topsep=2pt]
\item[\textbf{(B)}] A 15-node instance is brute-forced in $0.01$ s; small simulations cannot show a speed-up.
\item[\textbf{(C)}] No fixed-$p$ QAOA is known to solve an NP-hard problem exactly; that would imply surprising complexity results.
\item[\textbf{(D)}] $0.6924 < 0.878$: $p = 1$ QAOA has the weaker worst-case guarantee.
\end{itemize}
\textbf{\color{brand}Key takeaway.} Treat QAOA as a heuristic with honest benchmarks, not as a proven speed-up.
\end{solbox}
\qsep

\needspace{13\baselineskip}
\begin{qitem}{78}
\qmeta{Single-correct MCQ}{Custom weighted edges}
With \texttt{GRAPH = "custom"} and \texttt{CUSTOM\_EDGES = [(0, 1, 1.0), (1, 2, 2.5), (0, 2, -1.0)]} (3 nodes), the maximum cut and
the random-split baseline printed by the notebook are
\par\smallskip\noindent\begin{tabularx}{\linewidth}{@{}*{4}{L}@{}}
\optlabel{A}\,$3.5$ and $2.25$ & \correctlabel{B}\,$3.5$ and $1.25$ & \optlabel{C}\,$2.5$ and $1.25$ & \optlabel{D}\,$4.5$ and $2.25$
\end{tabularx}
\end{qitem}
\begin{solbox}{B}
\textbf{\color{brand}Correct answer.} \textcolor{okgreen}{\textbf{(B)}} $3.5$ and $1.25$

\textbf{\color{brand}Concepts required.} Weighted (including negative) Max-Cut; the baseline uses the signed total weight.

\textbf{\color{brand}Relevant formulae.} $C(z) = \sum w_{ij}[z_i\neq z_j]$; baseline $W/2$ with $W = \sum w_{ij}$ (signed).

\textbf{\color{brand}Step-by-step explanation.} Negative weights are allowed (``weights may be any real numbers''); a cut edge with weight $-1$ \emph{lowers} the cut.
The non-trivial partitions isolate one node: $\{0\}$: edges $(0,1), (0,2)$: $1 - 1 = 0$. $\{1\}$: $(0,1), (1,2)$: $1 + 2.5 = 3.5$.
$\{2\}$: $(1,2), (0,2)$: $2.5 - 1 = 1.5$. So $C_{\max} = 3.5$ (strings $010$ and $101$).
Baseline $W/2$ with $W = 1 + 2.5 - 1 = 2.5$: $1.25$. The rest of the pipeline is unchanged: the edge gets \texttt{rz(gamma * w)} with
$w = -1$ and $H_C$ the term $-0.5\,Z_0Z_2$.

\textbf{\color{brand}Why the other options are incorrect.}
\begin{itemize}[leftmargin=2.4em, itemsep=1pt, topsep=2pt]
\item[\textbf{(A)}] The baseline uses the signed $W = 2.5$, not $\sum|w| = 4.5$.
\item[\textbf{(C)}] $2.5$ is $W$, not the maximum cut (isolating node 1 gives $3.5$).
\item[\textbf{(D)}] $4.5$ and $2.25$ use $|w|$, as if the negative edge also counted positively.
\end{itemize}
\textbf{\color{brand}Key takeaway.} Negative weights reward keeping nodes together; check every split, not just the obvious ones.
\end{solbox}
\qsep

\needspace{20\baselineskip}
\begin{passagebox}{Paragraph for Questions 79 and 80}
\emph{Previously recorded run of \texttt{QAOA.ipynb}} (default settings, CUDA-Q 0.16.0, target \texttt{qpp-cpu}):
15 nodes, 30 unit edges ($W = 30$), exact optimum $C_{\max} = 25$ with the four optimal bitstrings listed in A14,
$p = 2$, COBYLA with 80 evaluations, best expected cut $\langle C\rangle = 20.5336$, $4000$ shots. The four most frequent
bitstrings were all optimal:
$\texttt{110010001001101}$ (68 shots), $\texttt{001101110110010}$ (67), $\texttt{110010001011101}$ (56), $\texttt{001101110100010}$ (49);
in total 240 shots were optimal.
\end{passagebox}
\needspace{13\baselineskip}
\begin{qitem}{79}
\qmeta{Paragraph based}{Recorded run: partition probability}
Estimated from the recorded counts, the probability of sampling the partition described by $\texttt{110010001001101}$
(either this bitstring or its complement) is
\par\smallskip\noindent\begin{tabularx}{\linewidth}{@{}*{4}{L}@{}}
\correctlabel{A}\,$0.0338$ & \optlabel{B}\,$0.0600$ & \optlabel{C}\,$0.0170$ & \optlabel{D}\,$0.0263$
\end{tabularx}
\end{qitem}
\begin{solbox}{A}
\textbf{\color{brand}Correct answer.} \textcolor{okgreen}{\textbf{(A)}} $0.0338$

\textbf{\color{brand}Concepts required.} A partition corresponds to a bitstring and its complement; estimating probabilities from counts.

\textbf{\color{brand}Relevant formulae.} $\hat P(\{z, \bar z\}) = (n_z + n_{\bar z})/S$.

\textbf{\color{brand}Step-by-step explanation.} The complement of $\texttt{110010001001101}$ is $\texttt{001101110110010}$ (flip every bit), which is the second most frequent
string. Both describe the same partition, so the estimate is $(68 + 67)/4000 = 135/4000 = 0.03375 \approx 0.0338$.
The other optimal partition ($\texttt{110010001011101}$ / $\texttt{001101110100010}$) has $(56 + 49)/4000 = 0.02625$; together
$0.03375 + 0.02625 = 0.06 = P_{\text{opt}}$. The nearly equal counts within each pair ($68$ vs $67$, $56$ vs $49$) reflect the exact
symmetry $|\braket{z|\psi}| = |\braket{\bar z|\psi}|$ of the QAOA state, up to shot noise.

\textbf{\color{brand}Why the other options are incorrect.}
\begin{itemize}[leftmargin=2.4em, itemsep=1pt, topsep=2pt]
\item[\textbf{(B)}] $0.0600$ is the probability of \emph{both} optimal partitions together.
\item[\textbf{(C)}] $0.0170 = 68/4000$ counts only one of the two bitstrings of the partition.
\item[\textbf{(D)}] $0.0263$ is the other optimal partition.
\end{itemize}
\textbf{\color{brand}Key takeaway.} Add $z$ and $\bar z$ to get a partition's probability; the two counts should agree up to shot noise.
\end{solbox}
\qsep

\needspace{13\baselineskip}
\begin{qitem}{80}
\qmeta{Paragraph based}{Recorded run: improvement over random}
A stricter figure of merit measures the improvement over random guessing:
$\eta = \dfrac{\langle C\rangle - W/2}{C_{\max} - W/2}$. For the recorded run, $\eta$ is
\par\smallskip\noindent\begin{tabularx}{\linewidth}{@{}*{4}{L}@{}}
\optlabel{A}\,$0.600$ & \correctlabel{B}\,$0.553$ & \optlabel{C}\,$0.221$ & \optlabel{D}\,$0.821$
\end{tabularx}
\end{qitem}
\begin{solbox}{B}
\textbf{\color{brand}Correct answer.} \textcolor{okgreen}{\textbf{(B)}} $0.553$

\textbf{\color{brand}Concepts required.} Normalising performance relative to random guessing.

\textbf{\color{brand}Relevant formulae.} $\eta = \dfrac{\langle C\rangle - W/2}{C_{\max} - W/2}$; $\ r_{\text{random}} = \dfrac{W/2}{C_{\max}}$.

\textbf{\color{brand}Step-by-step explanation.} $\eta = \dfrac{20.5336 - 15}{25 - 15} = \dfrac{5.5336}{10} \approx 0.553$.
The usual ratio $r = 20.5336/25 = 0.821$ gives random guessing already $15/25 = 0.6$, so $r$ looks better than the actual
progress. $\eta = 0$ for random guessing and $1$ for the optimum: QAOA at $p = 2$ closed about $55\,\%$ of the gap between random
guessing and the optimum.

\textbf{\color{brand}Why the other options are incorrect.}
\begin{itemize}[leftmargin=2.4em, itemsep=1pt, topsep=2pt]
\item[\textbf{(A)}] $0.600$ is the ratio of random guessing.
\item[\textbf{(C)}] $0.221 = 0.821 - 0.6$ subtracts ratios instead of normalising by the remaining gap.
\item[\textbf{(D)}] $0.821$ is the ordinary approximation ratio $r$.
\end{itemize}
\textbf{\color{brand}Key takeaway.} Quote $r$ together with the random baseline; $\eta$ shows the real progress.
\end{solbox}
\qsep

\end{document}